% ============================================================
%  Analyse Réelle I — Notes de Cours
%  Licence L1 — Premier semestre
%  « Des suites aux dérivées »
% ============================================================
\documentclass[12pt,a4paper,openany]{book}

% --- Encodage et langue ---
\usepackage[utf8]{inputenc}
\usepackage[T1]{fontenc}
\usepackage[french]{babel}

% --- Mise en page ---
\usepackage{booktabs}
\usepackage[margin=2.5cm,top=3cm,bottom=3cm]{geometry}
\usepackage{fancyhdr}
\pagestyle{fancy}
\fancyhf{}
\fancyhead[LE]{\leftmark}
\fancyhead[RO]{\rightmark}
\fancyfoot[C]{\thepage}
\renewcommand{\headrulewidth}{0.4pt}

% --- Mathématiques ---
\usepackage{amsmath,amssymb,amsthm,mathtools}

% --- Graphiques ---
\usepackage{tikz}
\usetikzlibrary{arrows.meta,positioning,calc,decorations.pathreplacing,patterns,shapes.geometric}

% --- Boîtes colorées ---
\usepackage[most,theorems,skins,breakable]{tcolorbox}

% --- Hyperliens et index ---
\usepackage[colorlinks=true,linkcolor=blue!60!black,citecolor=green!50!black,urlcolor=blue!70!black]{hyperref}
\usepackage{makeidx}
\makeindex

% --- Listes ---
\usepackage{enumitem}

% --- Couleurs ---
\usepackage{xcolor}
\definecolor{defblue}{RGB}{0,51,102}
\definecolor{defbluebg}{RGB}{220,235,252}
\definecolor{thmred}{RGB}{120,20,20}
\definecolor{thmbg}{RGB}{255,240,240}
\definecolor{propgreen}{RGB}{20,100,20}
\definecolor{propbg}{RGB}{235,250,235}
\definecolor{rembg}{RGB}{255,255,230}
\definecolor{exbg}{RGB}{245,245,245}
\definecolor{errbg}{RGB}{255,230,230}
\definecolor{exobg}{RGB}{240,248,255}

% ============================================================
%  Environnements de théorèmes (amsthm + tcolorbox wrapping)
% ============================================================
\theoremstyle{definition}
\newtheorem{definition}{Définition}[chapter]
\newtheorem{theoreme}{Théorème}[chapter]
\newtheorem{proposition}{Proposition}[chapter]
\newtheorem{lemme}{Lemme}[chapter]
\newtheorem{corollaire}{Corollaire}[chapter]
\newtheorem{remarque}{Remarque}[chapter]
\newtheorem{exemple}{Exemple}[chapter]
\newtheorem{exercice}{Exercice}[chapter]

\tcolorboxenvironment{definition}{enhanced,breakable,colback=defbluebg,colframe=defblue,left=3mm,right=3mm,top=2mm,bottom=2mm,before skip=10pt,after skip=10pt}
\tcolorboxenvironment{theoreme}{enhanced,breakable,colback=thmbg,colframe=thmred,left=3mm,right=3mm,top=2mm,bottom=2mm,before skip=10pt,after skip=10pt}
\tcolorboxenvironment{proposition}{enhanced,breakable,colback=propbg,colframe=propgreen,left=3mm,right=3mm,top=2mm,bottom=2mm,before skip=10pt,after skip=10pt}
\tcolorboxenvironment{lemme}{enhanced,breakable,colback=propbg,colframe=propgreen!70!black,left=3mm,right=3mm,top=2mm,bottom=2mm,before skip=10pt,after skip=10pt}
\tcolorboxenvironment{corollaire}{enhanced,breakable,colback=propbg,colframe=propgreen!50!black,left=3mm,right=3mm,top=2mm,bottom=2mm,before skip=10pt,after skip=10pt}
\tcolorboxenvironment{remarque}{enhanced,breakable,colback=rembg,colframe=orange!70!black,left=3mm,right=3mm,top=2mm,bottom=2mm,before skip=10pt,after skip=10pt}
\tcolorboxenvironment{exemple}{enhanced,breakable,colback=exbg,colframe=gray!70!black,left=3mm,right=3mm,top=2mm,bottom=2mm,before skip=10pt,after skip=10pt}
\tcolorboxenvironment{exercice}{enhanced,breakable,colback=exobg,colframe=blue!50!black,left=3mm,right=3mm,top=2mm,bottom=2mm,before skip=10pt,after skip=10pt}

% Boîte pour les erreurs fréquentes
\newtcolorbox{errmark}{%
  enhanced,breakable,
  colback=errbg,colframe=red!70!black,
  title={\textbf{Erreurs fréquentes}},
  fonttitle=\bfseries\color{red!80!black},
  left=3mm,right=3mm,top=2mm,bottom=2mm,
  before skip=10pt,after skip=10pt
}

% Boîte pour les résumés de chapitre
\newtcolorbox{resumchapitre}{%
  enhanced,breakable,
  colback=blue!5!white,colframe=blue!50!black,
  title={\textbf{Résumé du chapitre}},
  fonttitle=\bfseries,
  left=3mm,right=3mm,top=2mm,bottom=2mm,
  before skip=10pt,after skip=10pt
}

% ============================================================
%  Commandes personnalisées
% ============================================================
\newcommand{\N}{\mathbb{N}}
\newcommand{\Z}{\mathbb{Z}}
\newcommand{\Q}{\mathbb{Q}}
\newcommand{\R}{\mathbb{R}}
\newcommand{\C}{\mathbb{C}}
\newcommand{\eps}{\varepsilon}
\newcommand{\abs}[1]{\left\lvert #1 \right\rvert}
\newcommand{\norm}[1]{\left\lVert #1 \right\rVert}
\newcommand{\sumark}[1]{\sup\left(#1\right)}
\newcommand{\infmark}[1]{\inf\left(#1\right)}

% ============================================================
\begin{document}

% ============================================================
%  PAGE DE TITRE
% ============================================================
\begin{titlepage}
\begin{tikzpicture}[remember picture, overlay]
  \fill[left color=blue!15!white, right color=blue!5!white]
    (current page.south west) rectangle (current page.north east);
  \fill[color=orange!70!yellow]
    ([yshift=-2cm]current page.north west) rectangle ([yshift=-2.3cm]current page.north east);
  \fill[color=blue!80!black, rounded corners=8pt]
    ([xshift=3cm, yshift=-4cm]current page.north west) rectangle
    ([xshift=-3cm, yshift=-10cm]current page.north east);
  \node[anchor=center, text=white, font=\Huge\bfseries]
    at ([yshift=-6cm]current page.north) {Analyse R\'eelle I};
  \node[anchor=center, text=orange!80!yellow, font=\Large]
    at ([yshift=-7.5cm]current page.north) {Notes de Cours};
  \node[anchor=center, text=white!80, font=\large]
    at ([yshift=-8.8cm]current page.north) {Licence L2 --- 2025--2026};
  \node[anchor=center, text=blue!80!black, font=\Large\itshape]
    at ([yshift=-12cm]current page.north) {Ya\'e Ulrich Gaba};
  \draw[orange!70!yellow, line width=2pt]
    ([xshift=5cm, yshift=-13.5cm]current page.north west) --
    ([xshift=-5cm, yshift=-13.5cm]current page.north east);
  \node[anchor=center, text width=12cm, align=center, text=blue!60!black, font=\itshape]
    at ([yshift=-16cm]current page.north) {<<\,La rigueur est au math\'ematicien ce que la morale est \`a l'homme.\,>>\\[4pt]--- Henri Poincar\'e};
  \node[anchor=south, text=gray, font=\small]
    at ([yshift=1.5cm]current page.south) {\today};
\end{tikzpicture}
\end{titlepage}

% ============================================================
%  TABLE DES MATIÈRES
% ============================================================
\tableofcontents

% ============================================================
%  PRÉFACE
% ============================================================
\chapter*{Préface}
\addcontentsline{toc}{chapter}{Préface}
\markboth{Préface}{Préface}

Vous tenez entre vos mains un cours d'\emph{analyse réelle}. Si vous arrivez du lycée, vous avez déjà manipulé des limites, des dérivées, des intégrales~; vous avez peut-être même l'impression de << savoir >> ce que ces notions signifient. Ce cours va vous montrer que la réalité est plus subtile --- et plus belle --- que ce que vous imaginez.

\medskip

\textbf{Qu'est-ce que l'analyse~?} L'analyse mathématique est la branche des mathématiques qui étudie les processus de \emph{passage à la limite}. Quand on écrit $\lim_{n\to\infty} a_n = \ell$, qu'est-ce que cela signifie \emph{exactement}~? Quand on dit qu'une fonction est << continue >>, quelle est la définition \emph{précise}~? Pourquoi certaines fonctions sont-elles dérivables et d'autres non~? L'analyse répond à toutes ces questions avec une rigueur absolue.

\medskip

\textbf{Pourquoi des démonstrations~?} Au lycée, on vous a souvent présenté des résultats sans justification, ou avec des arguments intuitifs. À l'université, la règle change~: \emph{tout résultat doit être démontré}. Cela peut sembler difficile, voire décourageant au début. Mais la démonstration n'est pas un exercice de style~: c'est le seul moyen de s'assurer qu'un énoncé est \emph{vrai}. L'intuition est un guide précieux, mais elle peut tromper. L'histoire des mathématiques est remplie de conjectures << évidentes >> qui se sont révélées fausses.

\medskip

\textbf{Ce que vous allez apprendre.} Au fil de ce cours, vous découvrirez~:
\begin{itemize}[leftmargin=2cm]
  \item les techniques fondamentales de démonstration~: preuve directe, contraposée, raisonnement par l'absurde, récurrence~;
  \item la construction rigoureuse du corps des nombres réels $\R$ et l'axiome de la borne supérieure~;
  \item la théorie des suites réelles~: convergence, suites monotones, suites de Cauchy~;
  \item la continuité et ses conséquences profondes~: théorème des valeurs intermédiaires, théorème de Heine~;
  \item la dérivabilité et ses applications.
\end{itemize}

\medskip

\textbf{Conseils pour réussir.} L'analyse se travaille \emph{chaque jour}. Ne laissez pas s'accumuler les incompréhensions~: une notion non comprise au chapitre~2 rendra le chapitre~3 incompréhensible. Refaites les démonstrations du cours \emph{sans regarder vos notes}. Faites les exercices --- tous les exercices. Discutez avec vos camarades. Posez des questions. N'ayez pas peur de ne pas comprendre tout de suite~: l'analyse est une discipline qui demande du temps et de la maturation.

\medskip

\textbf{Un mot d'encouragement.} Si vous ressentez de l'anxiété face à l'abstraction et à la rigueur, sachez que c'est parfaitement normal. Tous les mathématiciens sont passés par cette étape. La difficulté que vous ressentez n'est pas un signe d'incompétence~: c'est le signe que vous êtes en train d'apprendre quelque chose de profondément nouveau. Faites-vous confiance, travaillez régulièrement, et vous verrez que la rigueur mathématique, loin d'être un obstacle, est une source de clarté et de satisfaction intellectuelle.

\vfill
\hfill\textit{Bon courage et bonne lecture.}


% ============================================================
%  NOTATIONS ET CONVENTIONS
% ============================================================
\chapter*{Notations et conventions}
\addcontentsline{toc}{chapter}{Notations et conventions}
\markboth{Notations et conventions}{Notations et conventions}

Nous rassemblons ici les notations utilisées tout au long de ce cours.

\bigskip

\begin{tabular}{@{}p{3.5cm}p{10cm}@{}}
\toprule
\textbf{Notation} & \textbf{Signification} \\
\midrule
$\N$ & Ensemble des entiers naturels $\{0,1,2,3,\ldots\}$. Dans ce cours, $0 \in \N$. \\[4pt]
$\N^*$ & $\N \setminus \{0\} = \{1,2,3,\ldots\}$. \\[4pt]
$\Z$ & Ensemble des entiers relatifs $\{\ldots,-2,-1,0,1,2,\ldots\}$. \\[4pt]
$\Q$ & Ensemble des nombres rationnels $\left\{\dfrac{p}{q} : p \in \Z,\, q \in \N^*\right\}$. \\[4pt]
$\R$ & Ensemble des nombres réels. \\[4pt]
$\C$ & Ensemble des nombres complexes. \\[4pt]
$\forall$ & << Pour tout >> (quantificateur universel). \\[4pt]
$\exists$ & << Il existe >> (quantificateur existentiel). \\[4pt]
$\exists!$ & << Il existe un unique >>. \\[4pt]
$\in$, $\notin$ & << appartient à >>, << n'appartient pas à >>. \\[4pt]
$\subset$, $\supset$ & Inclusion (large)~: $A \subset B$ signifie $\forall x,\, x \in A \Rightarrow x \in B$. \\[4pt]
$\cup$, $\cap$ & Réunion, intersection. \\[4pt]
$A \setminus B$ & Différence~: $\{x \in A : x \notin B\}$. \\[4pt]
$A^c$ & Complémentaire de $A$ (dans un ensemble ambiant fixé). \\[4pt]
$\emptyset$ & Ensemble vide. \\[4pt]
$[a,b]$ & Intervalle fermé $\{x \in \R : a \leq x \leq b\}$. \\[4pt]
$]a,b[$ & Intervalle ouvert $\{x \in \R : a < x < b\}$. \\[4pt]
$[a,b[$, $]a,b]$ & Intervalles semi-ouverts. \\[4pt]
$\abs{x}$ & Valeur absolue de $x \in \R$. \\[4pt]
$\sup A$, $\inf A$ & Borne supérieure, borne inférieure de $A \subset \R$. \\[4pt]
$\max A$, $\min A$ & Maximum, minimum de $A$ (quand ils existent). \\[4pt]
$\Rightarrow$ & Implication. \\[4pt]
$\Leftrightarrow$ & Équivalence. \\[4pt]
$\lnot P$ & Négation de la proposition $P$. \\[4pt]
\bottomrule
\end{tabular}

\bigskip

\noindent\textbf{Convention sur les quantificateurs.} Quand on écrit
\[
  \forall \eps > 0,\; \exists N \in \N,\; \forall n \geq N,\; \abs{a_n - \ell} < \eps,
\]
cela signifie~: << pour tout nombre réel $\eps$ strictement positif, il existe un entier naturel $N$ tel que, pour tout entier $n$ supérieur ou égal à $N$, on a $\abs{a_n - \ell} < \eps$ >>. L'ordre des quantificateurs est \emph{crucial}~: le changer modifie le sens de l'énoncé.


% ############################################################
%  CHAPITRE 1
% ############################################################
\chapter{Logique, ensembles et raisonnement mathématique}
\label{chap:logique}

\vspace{-0.5cm}
\begin{center}\itshape
  << Les mathématiques sont un exercice de la pensée, et la logique en est la grammaire. >>
\end{center}
\bigskip

Ce premier chapitre pose les fondations de tout le cours. Avant de parler de limites, de continuité ou de dérivées, nous devons apprendre à \emph{raisonner rigoureusement}. Nous allons donc étudier la logique mathématique, les ensembles et les principales méthodes de démonstration. Chacune de ces notions sera illustrée par de nombreux exemples.

% ============================================================
\section{Propositions logiques et connecteurs}
\label{sec:propositions}
% ============================================================

\begin{definition}[Proposition]\label{def:prop}
  \index{proposition}
  Une \textbf{proposition} (ou \textbf{assertion}) est un énoncé mathématique qui est soit \emph{vrai}, soit \emph{faux}. On dit que la proposition a une \textbf{valeur de vérité}.
\end{definition}

\begin{exemple}\label{ex:prop-exemples}
  \begin{enumerate}[label=(\alph*)]
    \item << $2 + 3 = 5$ >> est une proposition vraie.
    \item << $7$ est un nombre pair >> est une proposition fausse.
    \item << $\sqrt{2}$ est irrationnel >> est une proposition vraie (nous le démontrerons dans ce chapitre).
    \item << $x^2 \geq 0$ >> n'est \emph{pas} une proposition tant que $x$ n'est pas précisé~: c'est un \textbf{prédicat} (une proposition dépendant d'une variable).
  \end{enumerate}
\end{exemple}

\subsection{Les connecteurs logiques}

À partir de propositions simples, on construit des propositions composées à l'aide de \textbf{connecteurs logiques}\index{connecteur logique}.

\begin{definition}[Connecteurs logiques]\label{def:connecteurs}
  Soient $P$ et $Q$ deux propositions. On définit~:
  \begin{enumerate}[label=(\roman*)]
    \item \textbf{Négation}\index{négation}~: $\lnot P$ (lue << non $P$ >>) est vraie si et seulement si $P$ est fausse.
    \item \textbf{Conjonction}\index{conjonction}~: $P \land Q$ (lue << $P$ et $Q$ >>) est vraie si et seulement si $P$ et $Q$ sont toutes deux vraies.
    \item \textbf{Disjonction}\index{disjonction}~: $P \lor Q$ (lue << $P$ ou $Q$ >>) est vraie si et seulement si au moins l'une des deux propositions est vraie.
    \item \textbf{Implication}\index{implication}~: $P \Rightarrow Q$ (lue << $P$ implique $Q$ >> ou << si $P$ alors $Q$ >>) est fausse uniquement lorsque $P$ est vraie et $Q$ est fausse.
    \item \textbf{Équivalence}\index{équivalence}~: $P \Leftrightarrow Q$ (lue << $P$ si et seulement si $Q$ >>) est vraie lorsque $P$ et $Q$ ont la même valeur de vérité.
  \end{enumerate}
\end{definition}

\subsection{Tables de vérité}
\index{table de vérité}

Les \textbf{tables de vérité} donnent la valeur de vérité d'une proposition composée en fonction des valeurs de vérité de ses composantes. On note V pour << vrai >> et F pour << faux >>.

\bigskip
\begin{center}
\begin{tabular}{cc|c|c|c|c|c}
$P$ & $Q$ & $\lnot P$ & $P \land Q$ & $P \lor Q$ & $P \Rightarrow Q$ & $P \Leftrightarrow Q$ \\
\hline
V & V & F & V & V & V & V \\
V & F & F & F & V & F & F \\
F & V & V & F & V & V & F \\
F & F & V & F & F & V & V \\
\end{tabular}
\end{center}

\begin{remarque}[L'implication et le sens commun]\label{rem:impl-sens}
  Le fait que $P \Rightarrow Q$ soit \emph{vraie} quand $P$ est fausse surprend souvent. En mathématiques, l'implication << si $P$ alors $Q$ >> ne dit rien sur le cas où l'hypothèse $P$ n'est pas satisfaite. Par exemple, << si $2+2=5$, alors la Lune est un fromage >> est une implication \emph{vraie} au sens logique, car l'hypothèse est fausse.
\end{remarque}

\begin{remarque}[Implication et équivalence]\label{rem:impl-equiv}
  On a~: $P \Leftrightarrow Q$ est logiquement équivalent à $(P \Rightarrow Q) \land (Q \Rightarrow P)$. Autrement dit, pour montrer une équivalence, il faut démontrer les deux implications.
\end{remarque}


% ============================================================
\section{Quantificateurs}
\label{sec:quantificateurs}
% ============================================================

Les quantificateurs permettent de transformer un prédicat (qui dépend d'une variable) en une proposition.

\begin{definition}[Quantificateur universel]\label{def:quant-univ}
  \index{quantificateur!universel}
  Soit $P(x)$ un prédicat portant sur un élément $x$ d'un ensemble $E$. La proposition
  \[
    \forall x \in E,\; P(x)
  \]
  (lue << pour tout $x$ dans $E$, $P(x)$ >>) est \emph{vraie} si et seulement si $P(x)$ est vraie pour \emph{chaque} élément $x$ de $E$.
\end{definition}

\begin{definition}[Quantificateur existentiel]\label{def:quant-exist}
  \index{quantificateur!existentiel}
  La proposition
  \[
    \exists x \in E,\; P(x)
  \]
  (lue << il existe $x$ dans $E$ tel que $P(x)$ >>) est \emph{vraie} si et seulement si $P(x)$ est vraie pour \emph{au moins un} élément $x$ de $E$.
\end{definition}

\begin{exemple}\label{ex:quant-ex}
  \begin{enumerate}[label=(\alph*)]
    \item $\forall n \in \N,\; n^2 \geq 0$ est \textbf{vraie}.
    \item $\forall n \in \N,\; n^2 > 0$ est \textbf{fausse} (car $0^2 = 0$, qui n'est pas $> 0$).
    \item $\exists n \in \N,\; n^2 = 9$ est \textbf{vraie} (prendre $n = 3$).
    \item $\exists n \in \N,\; n^2 = 7$ est \textbf{fausse}.
  \end{enumerate}
\end{exemple}

\subsection{Négation de propositions quantifiées}
\label{subsec:negation-quant}

\index{négation!de quantificateurs}
La négation des quantificateurs est un point \emph{fondamental} que beaucoup d'étudiants maîtrisent mal. Retenez les deux règles suivantes~:

\begin{theoreme}[Négation des quantificateurs]\label{thm:neg-quant}
  \begin{enumerate}[label=(\roman*)]
    \item $\lnot\bigl(\forall x \in E,\; P(x)\bigr) \;\Leftrightarrow\; \exists x \in E,\; \lnot P(x)$.
    \item $\lnot\bigl(\exists x \in E,\; P(x)\bigr) \;\Leftrightarrow\; \forall x \in E,\; \lnot P(x)$.
  \end{enumerate}
  Autrement dit~: pour nier un << pour tout >>, on le remplace par un << il existe >> et on nie la propriété~; pour nier un << il existe >>, on le remplace par un << pour tout >> et on nie la propriété.
\end{theoreme}

\begin{remarque}[Règle mnémotechnique]\label{rem:mnemo-neg}
  Quand on nie, chaque quantificateur << bascule >>~: $\forall$ devient $\exists$ et $\exists$ devient $\forall$. Puis on nie la propriété finale. On procède de gauche à droite, un quantificateur à la fois.
\end{remarque}

\begin{exemple}[Négation de la convergence d'une suite]\label{ex:neg-convergence}
  La définition de << la suite $(a_n)$ converge vers $\ell$ >> s'écrit~:
  \[
    \forall \eps > 0,\; \exists N \in \N,\; \forall n \geq N,\; \abs{a_n - \ell} < \eps.
  \]
  Nions cette proposition \emph{étape par étape}~:
  \begin{enumerate}[label=\textbf{Étape \arabic*} :]
    \item On nie le premier quantificateur. $\forall \eps > 0$ devient $\exists \eps > 0$, et on nie le reste~:
    \[
      \exists \eps > 0,\; \lnot\Bigl(\exists N \in \N,\; \forall n \geq N,\; \abs{a_n - \ell} < \eps\Bigr).
    \]
    \item On nie le deuxième quantificateur. $\exists N \in \N$ devient $\forall N \in \N$, et on nie le reste~:
    \[
      \exists \eps > 0,\; \forall N \in \N,\; \lnot\Bigl(\forall n \geq N,\; \abs{a_n - \ell} < \eps\Bigr).
    \]
    \item On nie le troisième quantificateur. $\forall n \geq N$ devient $\exists n \geq N$, et on nie la propriété~:
    \[
      \exists \eps > 0,\; \forall N \in \N,\; \exists n \geq N,\; \lnot\bigl(\abs{a_n - \ell} < \eps\bigr).
    \]
    \item On nie l'inégalité~: $\lnot(\abs{a_n - \ell} < \eps)$ est $\abs{a_n - \ell} \geq \eps$.
  \end{enumerate}

  \medskip
  \textbf{Résultat final.} La suite $(a_n)$ \emph{ne converge pas} vers $\ell$ si et seulement si~:
  \[
    \boxed{\exists \eps > 0,\; \forall N \in \N,\; \exists n \geq N,\; \abs{a_n - \ell} \geq \eps.}
  \]
  En français~: << il existe un $\eps > 0$ tel que, quel que soit le rang $N$, on peut toujours trouver un terme $a_n$ (avec $n \geq N$) qui reste à distance au moins $\eps$ de $\ell$ >>.
\end{exemple}

\begin{errmark}
  \begin{itemize}
    \item \textbf{Erreur~:} écrire $\lnot(\forall x,\, P(x))$ comme $\forall x,\, \lnot P(x)$. La négation du << pour tout >> est << il existe >>, pas << pour tout >>.
    \item \textbf{Erreur~:} oublier de nier la propriété finale après avoir basculé tous les quantificateurs.
    \item \textbf{Erreur~:} nier $<$ en $>$ au lieu de $\geq$. La négation de $a < b$ est $a \geq b$, \emph{pas} $a > b$.
  \end{itemize}
\end{errmark}


% ============================================================
\section{Méthodes de démonstration}
\label{sec:methodes-demo}
% ============================================================

Nous présentons maintenant les principales méthodes de démonstration\index{démonstration} que vous utiliserez tout au long de vos études.

% -----------------------------------------------
\subsection{Preuve directe}
\label{subsec:preuve-directe}
\index{preuve directe}
% -----------------------------------------------

Pour montrer $P \Rightarrow Q$, on \emph{suppose} que $P$ est vraie et on en \emph{déduit} que $Q$ est vraie par une chaîne de déductions logiques.

\begin{exemple}[La somme de deux entiers pairs est paire]\label{ex:paire-somme}
  \textbf{Énoncé.} Soient $a, b \in \Z$. Si $a$ et $b$ sont pairs, alors $a + b$ est pair.

  \begin{proof}
    Supposons que $a$ et $b$ sont pairs. Par définition, il existe $k, l \in \Z$ tels que $a = 2k$ et $b = 2l$. Alors~:
    \[
      a + b = 2k + 2l = 2(k + l).
    \]
    Comme $k + l \in \Z$, le nombre $a + b$ est bien pair.
  \end{proof}
\end{exemple}

\begin{exemple}[Le carré d'un entier impair est impair]\label{ex:carre-impair}
  \textbf{Énoncé.} Soit $n \in \Z$. Si $n$ est impair, alors $n^2$ est impair.

  \begin{proof}
    Supposons que $n$ est impair. Il existe $k \in \Z$ tel que $n = 2k + 1$. Alors~:
    \[
      n^2 = (2k+1)^2 = 4k^2 + 4k + 1 = 2(2k^2 + 2k) + 1.
    \]
    En posant $m = 2k^2 + 2k \in \Z$, on obtient $n^2 = 2m + 1$, ce qui montre que $n^2$ est impair.
  \end{proof}
\end{exemple}

% -----------------------------------------------
\subsection{Preuve par contraposée}
\label{subsec:contraposee}
\index{contraposée}
% -----------------------------------------------

Pour montrer $P \Rightarrow Q$, on peut montrer sa \textbf{contraposée}~: $\lnot Q \Rightarrow \lnot P$. Ces deux implications sont logiquement équivalentes.

\begin{proposition}[Équivalence avec la contraposée]\label{prop:equiv-contraposee}
  Pour toutes propositions $P$ et $Q$~:
  \[
    (P \Rightarrow Q) \;\Leftrightarrow\; (\lnot Q \Rightarrow \lnot P).
  \]
\end{proposition}

\begin{proof}
  On vérifie par table de vérité. $P \Rightarrow Q$ est fausse uniquement quand $P$ est V et $Q$ est F. Dans ce cas, $\lnot Q$ est V et $\lnot P$ est F, donc $\lnot Q \Rightarrow \lnot P$ est aussi fausse. Dans tous les autres cas, les deux implications sont vraies.
\end{proof}

\begin{exemple}[Si $n^2$ est pair, alors $n$ est pair]\label{ex:carre-pair}
  \textbf{Énoncé.} Soit $n \in \Z$. Si $n^2$ est pair, alors $n$ est pair.

  \begin{proof}
    Montrons la contraposée~: << si $n$ est impair, alors $n^2$ est impair >>. C'est exactement l'Exemple~\ref{ex:carre-impair}. Donc la contraposée est vraie, et par conséquent l'implication originale l'est aussi.
  \end{proof}
\end{exemple}

\begin{exemple}\label{ex:contraposee-div}
  \textbf{Énoncé.} Soient $a, b \in \N$. Si $ab$ est impair, alors $a$ est impair et $b$ est impair.

  \begin{proof}
    Montrons la contraposée~: << si $a$ est pair ou $b$ est pair, alors $ab$ est pair >>.

    Supposons que $a$ est pair (le cas << $b$ est pair >> est analogue). Il existe $k \in \Z$ tel que $a = 2k$. Alors $ab = 2kb = 2(kb)$, qui est pair. Donc la contraposée est établie.
  \end{proof}
\end{exemple}

% -----------------------------------------------
\subsection{Preuve par l'absurde}
\label{subsec:absurde}
\index{raisonnement par l'absurde}
% -----------------------------------------------

Pour montrer qu'une proposition $P$ est vraie, on \emph{suppose que $P$ est fausse} (c'est-à-dire que $\lnot P$ est vraie) et on en déduit une \textbf{contradiction}~: on arrive à un énoncé de la forme $Q \land \lnot Q$.

\begin{exemple}[L'ensemble des nombres premiers est infini]\label{ex:premiers-infini}
  \textbf{Énoncé.} Il existe une infinité de nombres premiers.

  \begin{proof}
    Supposons, par l'absurde, qu'il n'existe qu'un nombre fini de nombres premiers. Notons-les $p_1, p_2, \ldots, p_k$. Considérons le nombre~:
    \[
      N = p_1 \cdot p_2 \cdots p_k + 1.
    \]
    Le nombre $N$ est supérieur à $1$, donc il admet au moins un diviseur premier $p$. Ce nombre $p$ doit être l'un des $p_i$ (car, par hypothèse, la liste $p_1, \ldots, p_k$ contient \emph{tous} les nombres premiers). Donc $p = p_i$ pour un certain $i$, et $p_i$ divise $N$.

    Or, $p_i$ divise aussi le produit $p_1 \cdot p_2 \cdots p_k$. Donc $p_i$ divise la différence~:
    \[
      N - p_1 \cdot p_2 \cdots p_k = 1.
    \]
    Ainsi $p_i$ divise $1$, ce qui est impossible car $p_i \geq 2$.

    Nous avons obtenu une contradiction. L'hypothèse de départ (nombre fini de premiers) est donc fausse. Il existe bien une infinité de nombres premiers.
  \end{proof}
\end{exemple}

\begin{exemple}[Irrationalité de $\sqrt{2}$]\label{ex:sqrt2}
  \index{irrationalité de $\sqrt{2}$}
  \textbf{Énoncé.} $\sqrt{2}$ est irrationnel, c'est-à-dire $\sqrt{2} \notin \Q$.

  \begin{proof}
    Supposons, par l'absurde, que $\sqrt{2} \in \Q$. Alors il existe $p \in \Z$ et $q \in \N^*$ tels que
    \[
      \sqrt{2} = \frac{p}{q},
    \]
    où la fraction $\frac{p}{q}$ est \emph{irréductible}, c'est-à-dire que $p$ et $q$ n'ont aucun facteur commun autre que $1$ (en particulier, $p$ et $q$ ne sont pas tous les deux pairs).

    En élevant au carré, on obtient~:
    \[
      2 = \frac{p^2}{q^2}, \quad \text{donc} \quad p^2 = 2q^2.
    \]

    \textbf{Étape 1~: $p$ est pair.} L'égalité $p^2 = 2q^2$ montre que $p^2$ est pair. D'après l'Exemple~\ref{ex:carre-pair}, si $p^2$ est pair, alors $p$ est pair. Donc il existe $k \in \Z$ tel que $p = 2k$.

    \textbf{Étape 2~: $q$ est pair.} En substituant $p = 2k$ dans $p^2 = 2q^2$~:
    \[
      (2k)^2 = 2q^2 \implies 4k^2 = 2q^2 \implies q^2 = 2k^2.
    \]
    Donc $q^2$ est pair, et par le même argument, $q$ est pair.

    \textbf{Contradiction.} Nous avons montré que $p$ et $q$ sont tous les deux pairs, ce qui contredit l'hypothèse que la fraction $\frac{p}{q}$ est irréductible.

    Conclusion~: $\sqrt{2}$ est irrationnel.
  \end{proof}
\end{exemple}

% -----------------------------------------------
\subsection{Raisonnement par récurrence}
\label{subsec:recurrence}
\index{récurrence}
% -----------------------------------------------

Le raisonnement par récurrence permet de démontrer qu'une propriété $P(n)$ est vraie pour tout $n \geq n_0$ (où $n_0 \in \N$ est un point de départ).

\begin{theoreme}[Principe de récurrence]\label{thm:recurrence}
  Soit $P(n)$ une propriété dépendant d'un entier $n \geq n_0$. Si~:
  \begin{enumerate}[label=(\roman*)]
    \item \textbf{Initialisation~:} $P(n_0)$ est vraie~;
    \item \textbf{Hérédité~:} pour tout $n \geq n_0$, $P(n) \Rightarrow P(n+1)$~;
  \end{enumerate}
  alors $P(n)$ est vraie pour tout $n \geq n_0$.
\end{theoreme}

\begin{remarque}[Analogie des dominos]\label{rem:dominos}
  Imaginez une file infinie de dominos numérotés $n_0, n_0+1, n_0+2, \ldots$ L'initialisation consiste à pousser le premier domino. L'hérédité assure que si le domino $n$ tombe, le domino $n+1$ tombe aussi. Conclusion~: tous les dominos tombent.
\end{remarque}

\begin{exemple}[Somme des $n$ premiers entiers]\label{ex:somme-entiers}
  \textbf{Énoncé.} Pour tout $n \geq 1$, $\displaystyle\sum_{k=1}^{n} k = \frac{n(n+1)}{2}$.

  \begin{proof}
    Notons $P(n)$ la propriété << $\sum_{k=1}^{n} k = \frac{n(n+1)}{2}$ >>.

    \textbf{Initialisation~:} Pour $n = 1$, le membre de gauche vaut $1$, et le membre de droite vaut $\frac{1 \cdot 2}{2} = 1$. Donc $P(1)$ est vraie.

    \textbf{Hérédité~:} Soit $n \geq 1$. Supposons que $P(n)$ est vraie (c'est l'\textbf{hypothèse de récurrence}). Montrons $P(n+1)$, c'est-à-dire~: $\sum_{k=1}^{n+1} k = \frac{(n+1)(n+2)}{2}$.

    On a~:
    \[
      \sum_{k=1}^{n+1} k = \left(\sum_{k=1}^{n} k\right) + (n+1) = \frac{n(n+1)}{2} + (n+1)
    \]
    par hypothèse de récurrence. On factorise~:
    \[
      = (n+1)\left(\frac{n}{2} + 1\right) = (n+1) \cdot \frac{n+2}{2} = \frac{(n+1)(n+2)}{2}.
    \]
    Donc $P(n+1)$ est vraie.

    Par le principe de récurrence, $P(n)$ est vraie pour tout $n \geq 1$.
  \end{proof}
\end{exemple}

\begin{exemple}[Inégalité de Bernoulli]\label{ex:bernoulli}
  \index{inégalité de Bernoulli}
  \textbf{Énoncé.} Pour tout $x \geq -1$ et tout $n \in \N$, $(1+x)^n \geq 1 + nx$.

  \begin{proof}
    Fixons $x \geq -1$. Soit $P(n)$ la propriété << $(1+x)^n \geq 1+nx$ >>.

    \textbf{Initialisation~:} Pour $n = 0$~: $(1+x)^0 = 1$ et $1 + 0 \cdot x = 1$. On a bien $1 \geq 1$. $P(0)$ est vraie.

    \textbf{Hérédité~:} Soit $n \geq 0$. Supposons $P(n)$ vraie~: $(1+x)^n \geq 1+nx$. Alors~:
    \[
      (1+x)^{n+1} = (1+x)^n \cdot (1+x) \geq (1+nx)(1+x)
    \]
    car $(1+x) \geq 0$ (puisque $x \geq -1$). En développant~:
    \[
      (1+nx)(1+x) = 1 + x + nx + nx^2 = 1 + (n+1)x + nx^2 \geq 1 + (n+1)x
    \]
    car $nx^2 \geq 0$. Donc $(1+x)^{n+1} \geq 1 + (n+1)x$, c'est-à-dire $P(n+1)$.

    Par le principe de récurrence, $P(n)$ est vraie pour tout $n \in \N$.
  \end{proof}
\end{exemple}

\begin{theoreme}[Récurrence forte]\label{thm:recurrence-forte}
  \index{récurrence!forte}
  Soit $P(n)$ une propriété dépendant d'un entier $n \geq n_0$. Si~:
  \begin{enumerate}[label=(\roman*)]
    \item $P(n_0)$ est vraie~;
    \item pour tout $n \geq n_0$, $\bigl(P(n_0) \land P(n_0+1) \land \cdots \land P(n)\bigr) \Rightarrow P(n+1)$~;
  \end{enumerate}
  alors $P(n)$ est vraie pour tout $n \geq n_0$.
\end{theoreme}

\begin{exemple}[Tout entier $n \geq 2$ admet un diviseur premier]\label{ex:diviseur-premier}
  \begin{proof}
    Soit $P(n)$ la propriété << $n$ admet un diviseur premier >> pour $n \geq 2$.

    \textbf{Initialisation~:} $P(2)$ est vraie car $2$ est lui-même premier et se divise lui-même.

    \textbf{Hérédité (récurrence forte)~:} Soit $n \geq 2$. Supposons que $P(k)$ est vraie pour tout $2 \leq k \leq n$. Montrons $P(n+1)$.

    \emph{Cas 1~:} Si $n+1$ est premier, alors $n+1$ est son propre diviseur premier. $P(n+1)$ est vraie.

    \emph{Cas 2~:} Si $n+1$ n'est pas premier, alors $n+1 = ab$ avec $2 \leq a, b \leq n$. Par hypothèse de récurrence (forte), $a$ admet un diviseur premier $p$. Comme $p$ divise $a$ et $a$ divise $n+1$, le nombre $p$ divise $n+1$. Donc $P(n+1)$ est vraie.

    Par récurrence forte, $P(n)$ est vraie pour tout $n \geq 2$.
  \end{proof}
\end{exemple}


% ============================================================
\section{Ensembles}
\label{sec:ensembles}
\index{ensemble}
% ============================================================

\begin{definition}[Ensemble]\label{def:ensemble}
  Un \textbf{ensemble} est une collection d'objets appelés \textbf{éléments}. On note $x \in E$ pour dire que $x$ est un élément de $E$, et $x \notin E$ dans le cas contraire.
\end{definition}

\subsection{Opérations sur les ensembles}

\begin{definition}[Opérations ensemblistes]\label{def:ops-ens}
  \index{réunion}\index{intersection}\index{complémentaire}
  Soient $A$ et $B$ deux sous-ensembles d'un ensemble $E$.
  \begin{enumerate}[label=(\roman*)]
    \item \textbf{Réunion~:} $A \cup B = \{x \in E : x \in A \text{ ou } x \in B\}$.
    \item \textbf{Intersection~:} $A \cap B = \{x \in E : x \in A \text{ et } x \in B\}$.
    \item \textbf{Différence~:} $A \setminus B = \{x \in E : x \in A \text{ et } x \notin B\}$.
    \item \textbf{Complémentaire~:} $A^c = E \setminus A = \{x \in E : x \notin A\}$.
    \item \textbf{Produit cartésien~:}\index{produit cartésien} $A \times B = \{(a,b) : a \in A,\, b \in B\}$.
  \end{enumerate}
\end{definition}

\bigskip
\begin{center}
\begin{tikzpicture}[scale=0.8]
  % Réunion
  \begin{scope}[shift={(-5,0)}]
    \draw[thick,fill=blue!15] (-0.6,0) circle (1.2);
    \draw[thick,fill=blue!15] (0.6,0) circle (1.2);
    \node at (-1.0,0.0) {$A$};
    \node at (1.0,0.0) {$B$};
    \node at (0,-2) {$A \cup B$};
  \end{scope}
  % Intersection
  \begin{scope}[shift={(0,0)}]
    \draw[thick] (-0.6,0) circle (1.2);
    \draw[thick] (0.6,0) circle (1.2);
    \begin{scope}
      \clip (-0.6,0) circle (1.2);
      \fill[blue!25] (0.6,0) circle (1.2);
    \end{scope}
    \node at (-1.0,0.0) {$A$};
    \node at (1.0,0.0) {$B$};
    \node at (0,-2) {$A \cap B$};
  \end{scope}
  % Différence
  \begin{scope}[shift={(5,0)}]
    \draw[thick,fill=blue!15] (-0.6,0) circle (1.2);
    \draw[thick,fill=white] (0.6,0) circle (1.2);
    \node at (-1.0,0.0) {$A$};
    \node at (1.0,0.0) {$B$};
    \node at (0,-2) {$A \setminus B$};
  \end{scope}
\end{tikzpicture}
\end{center}

\begin{theoreme}[Lois de De Morgan]\label{thm:de-morgan}
  \index{lois de De Morgan}
  Soient $A$ et $B$ deux sous-ensembles d'un ensemble $E$. Alors~:
  \begin{enumerate}[label=(\roman*)]
    \item $(A \cup B)^c = A^c \cap B^c$,
    \item $(A \cap B)^c = A^c \cup B^c$.
  \end{enumerate}
\end{theoreme}

\begin{proof}
  Démontrons (i). On va montrer la double inclusion.

  \textbf{Sens $\subset$~:} Soit $x \in (A \cup B)^c$. Par définition du complémentaire, $x \notin A \cup B$. Cela signifie que $x \notin A$ \emph{et} $x \notin B$ (car si $x$ appartenait à l'un des deux, il serait dans la réunion). Donc $x \in A^c$ et $x \in B^c$, c'est-à-dire $x \in A^c \cap B^c$.

  \textbf{Sens $\supset$~:} Soit $x \in A^c \cap B^c$. Alors $x \in A^c$ (donc $x \notin A$) et $x \in B^c$ (donc $x \notin B$). Comme $x$ n'est ni dans $A$ ni dans $B$, on a $x \notin A \cup B$, donc $x \in (A \cup B)^c$.

  La preuve de (ii) est analogue (ou bien on applique (i) à $A^c$ et $B^c$ puis on prend le complémentaire).
\end{proof}


% ============================================================
\section{Applications}
\label{sec:applications}
\index{application}
% ============================================================

\begin{definition}[Application]\label{def:application}
  Soient $E$ et $F$ deux ensembles. Une \textbf{application} (ou \textbf{fonction}) de $E$ dans $F$ est une relation qui, à chaque élément $x \in E$, associe un \emph{unique} élément $f(x) \in F$. On écrit $f : E \to F$, $x \mapsto f(x)$.
\end{definition}

\begin{definition}[Injection, surjection, bijection]\label{def:inj-surj-bij}
  \index{injection}\index{surjection}\index{bijection}
  Soit $f : E \to F$ une application.
  \begin{enumerate}[label=(\roman*)]
    \item $f$ est \textbf{injective} si~: $\forall x_1, x_2 \in E,\; f(x_1) = f(x_2) \Rightarrow x_1 = x_2$.
    \item $f$ est \textbf{surjective} si~: $\forall y \in F,\; \exists x \in E,\; f(x) = y$.
    \item $f$ est \textbf{bijective} si elle est à la fois injective et surjective.
  \end{enumerate}
\end{definition}

\bigskip
\begin{center}
\begin{tikzpicture}[
  >=Stealth,
  every node/.style={font=\small},
  ens/.style={ellipse, draw, minimum width=2cm, minimum height=3cm, thick}
]
  % Injection
  \begin{scope}[shift={(-5,0)}]
    \node[ens] (E1) at (0,0) {};
    \node[ens] (F1) at (2.5,0) {};
    \node[above] at (0,1.7) {$E$};
    \node[above] at (2.5,1.7) {$F$};
    \node[below] at (1.25,-2) {\textbf{Injection}};
    \foreach \y/\label in {0.7/a, 0/b, -0.7/c} {
      \fill (0,\y) circle (2pt) node[left] {$\label$};
    }
    \foreach \y/\label in {0.9/1, 0.3/2, -0.3/3, -0.9/4} {
      \fill (2.5,\y) circle (2pt) node[right] {$\label$};
    }
    \draw[->] (0.15,0.7) -- (2.35,0.9);
    \draw[->] (0.15,0) -- (2.35,-0.3);
    \draw[->] (0.15,-0.7) -- (2.35,-0.9);
  \end{scope}

  % Surjection
  \begin{scope}[shift={(0,0)}]
    \node[ens] (E2) at (0,0) {};
    \node[ens] (F2) at (2.5,0) {};
    \node[above] at (0,1.7) {$E$};
    \node[above] at (2.5,1.7) {$F$};
    \node[below] at (1.25,-2) {\textbf{Surjection}};
    \foreach \y/\label in {0.9/a, 0.3/b, -0.3/c, -0.9/d} {
      \fill (0,\y) circle (2pt) node[left] {$\label$};
    }
    \foreach \y/\label in {0.7/1, 0/2, -0.7/3} {
      \fill (2.5,\y) circle (2pt) node[right] {$\label$};
    }
    \draw[->] (0.15,0.9) -- (2.35,0.7);
    \draw[->] (0.15,0.3) -- (2.35,0.7);
    \draw[->] (0.15,-0.3) -- (2.35,0);
    \draw[->] (0.15,-0.9) -- (2.35,-0.7);
  \end{scope}

  % Bijection
  \begin{scope}[shift={(5,0)}]
    \node[ens] (E3) at (0,0) {};
    \node[ens] (F3) at (2.5,0) {};
    \node[above] at (0,1.7) {$E$};
    \node[above] at (2.5,1.7) {$F$};
    \node[below] at (1.25,-2) {\textbf{Bijection}};
    \foreach \y/\label in {0.7/a, 0/b, -0.7/c} {
      \fill (0,\y) circle (2pt) node[left] {$\label$};
    }
    \foreach \y/\label in {0.7/1, 0/2, -0.7/3} {
      \fill (2.5,\y) circle (2pt) node[right] {$\label$};
    }
    \draw[->] (0.15,0.7) -- (2.35,0.7);
    \draw[->] (0.15,0) -- (2.35,0);
    \draw[->] (0.15,-0.7) -- (2.35,-0.7);
  \end{scope}
\end{tikzpicture}
\end{center}

\begin{exemple}\label{ex:exemples-inj-surj}
  \begin{enumerate}[label=(\alph*)]
    \item $f : \R \to \R$, $x \mapsto x^2$ n'est ni injective (car $f(1) = f(-1) = 1$) ni surjective (car $-1$ n'a pas d'antécédent).
    \item $f : \R \to \R^+$, $x \mapsto x^2$ est surjective mais pas injective.
    \item $f : \R^+ \to \R^+$, $x \mapsto x^2$ est bijective.
    \item $f : \R \to \R$, $x \mapsto 2x + 1$ est bijective (sa réciproque est $g(y) = \frac{y-1}{2}$).
  \end{enumerate}
\end{exemple}


% ============================================================
\section{Relations d'équivalence et relations d'ordre}
\label{sec:relations}
% ============================================================

\begin{definition}[Relation d'équivalence]\label{def:rel-equiv}
  \index{relation d'équivalence}
  Une relation $\mathcal{R}$ sur un ensemble $E$ est une \textbf{relation d'équivalence} si elle est~:
  \begin{enumerate}[label=(\roman*)]
    \item \textbf{réflexive~:} $\forall x \in E,\; x \mathrel{\mathcal{R}} x$~;
    \item \textbf{symétrique~:} $\forall x, y \in E,\; x \mathrel{\mathcal{R}} y \Rightarrow y \mathrel{\mathcal{R}} x$~;
    \item \textbf{transitive~:} $\forall x, y, z \in E,\; (x \mathrel{\mathcal{R}} y \text{ et } y \mathrel{\mathcal{R}} z) \Rightarrow x \mathrel{\mathcal{R}} z$.
  \end{enumerate}
\end{definition}

\begin{exemple}\label{ex:equiv-mod}
  Sur $\Z$, la relation << $a \equiv b \pmod{n}$ >> (c'est-à-dire $n \mid (a-b)$) est une relation d'équivalence pour tout $n \in \N^*$.
\end{exemple}

\begin{definition}[Relation d'ordre]\label{def:rel-ordre}
  \index{relation d'ordre}
  Une relation $\leq$ sur un ensemble $E$ est une \textbf{relation d'ordre} si elle est~:
  \begin{enumerate}[label=(\roman*)]
    \item \textbf{réflexive~:} $\forall x \in E,\; x \leq x$~;
    \item \textbf{antisymétrique~:} $\forall x, y \in E,\; (x \leq y \text{ et } y \leq x) \Rightarrow x = y$~;
    \item \textbf{transitive~:} $\forall x, y, z \in E,\; (x \leq y \text{ et } y \leq z) \Rightarrow x \leq z$.
  \end{enumerate}
  L'ordre est \textbf{total} si~: $\forall x, y \in E,\; x \leq y \text{ ou } y \leq x$.
\end{definition}

\begin{exemple}\label{ex:ordres}
  \begin{enumerate}[label=(\alph*)]
    \item L'ordre usuel $\leq$ sur $\R$ est un ordre total.
    \item L'inclusion $\subset$ sur $\mathcal{P}(E)$ est un ordre partiel (non total en général).
    \item La divisibilité $\mid$ sur $\N^*$ est un ordre partiel.
  \end{enumerate}
\end{exemple}


% ============================================================
\section*{Erreurs fréquentes du chapitre}
\addcontentsline{toc}{section}{Erreurs fréquentes du chapitre}
% ============================================================

\begin{errmark}
  \begin{enumerate}
    \item \textbf{Confondre $\Rightarrow$ et $\Leftrightarrow$.} Écrire $\Leftrightarrow$ quand on ne prouve qu'un seul sens est une faute grave. Vérifiez toujours si vos implications sont réversibles.
    \item \textbf{Mauvaise négation des quantificateurs.} La négation de << $\forall x,\, P(x)$ >> n'est \emph{pas} << $\forall x,\, \lnot P(x)$ >> mais << $\exists x,\, \lnot P(x)$ >>. Revoyez la section~\ref{subsec:negation-quant}.
    \item \textbf{Oublier l'initialisation dans une récurrence.} Sans initialisation, la récurrence ne prouve rien.
    \item \textbf{Confondre << supposons $P(n)$ >> et << montrons $P(n)$ >> dans l'hérédité.} On \emph{suppose} $P(n)$ (hypothèse de récurrence) et on \emph{montre} $P(n+1)$.
    \item \textbf{Confondre $\in$ et $\subset$.} Un élément \emph{appartient} à un ensemble~; un ensemble est \emph{inclus} dans un autre.
  \end{enumerate}
\end{errmark}


% ============================================================
\section*{Exercices}
\addcontentsline{toc}{section}{Exercices}
% ============================================================

\begin{exercice}[Connecteurs logiques]\label{exo:exo-connecteurs}
  Construire la table de vérité de la proposition $(P \Rightarrow Q) \Leftrightarrow (\lnot P \lor Q)$. Qu'observe-t-on~?
\end{exercice}

\begin{exercice}[Négation]\label{exo:exo-negation1}
  Écrire la négation de chacune des propositions suivantes~:
  \begin{enumerate}[label=(\alph*)]
    \item $\forall x \in \R,\; x^2 \geq 0$.
    \item $\exists n \in \N,\; n^2 = n$.
    \item $\forall x \in \R,\; \exists y \in \R,\; x + y = 0$.
    \item $\forall \eps > 0,\; \exists \delta > 0,\; \forall x \in \R,\; \abs{x - a} < \delta \Rightarrow \abs{f(x) - f(a)} < \eps$.
  \end{enumerate}
\end{exercice}

\begin{exercice}[Preuve directe]\label{exo:exo-directe}
  Montrer que pour tout $n \in \N$, le nombre $n^2 + n$ est pair.
\end{exercice}

\begin{exercice}[Contraposée]\label{exo:exo-contra1}
  Montrer par contraposée~: si $n^2$ est divisible par $3$, alors $n$ est divisible par $3$.
\end{exercice}

\begin{exercice}[Absurde]\label{exo:exo-absurde1}
  Montrer par l'absurde que $\sqrt{3}$ est irrationnel.
\end{exercice}

\begin{exercice}[Récurrence (facile)]\label{exo:exo-recurrence1}
  Montrer par récurrence que pour tout $n \geq 1$~:
  \[
    \sum_{k=1}^{n} k^2 = \frac{n(n+1)(2n+1)}{6}.
  \]
\end{exercice}

\begin{exercice}[Récurrence (moyen)]\label{exo:exo-recurrence2}
  Montrer par récurrence que pour tout $n \geq 1$, $7^n - 1$ est divisible par $6$.
\end{exercice}

\begin{exercice}[Récurrence forte]\label{exo:exo-recurrence-forte}
  Montrer par récurrence forte que tout entier $n \geq 2$ peut s'écrire comme produit de nombres premiers. (C'est une partie du théorème fondamental de l'arithmétique.)
\end{exercice}

\begin{exercice}[Ensembles]\label{exo:exo-ensembles}
  Soient $A$, $B$, $C$ trois sous-ensembles d'un ensemble $E$.
  \begin{enumerate}[label=(\alph*)]
    \item Montrer que $A \cap (B \cup C) = (A \cap B) \cup (A \cap C)$.
    \item Montrer que $A \setminus (B \cap C) = (A \setminus B) \cup (A \setminus C)$.
  \end{enumerate}
\end{exercice}

\begin{exercice}[Injection et surjection]\label{exo:exo-inj-surj}
  Soit $f : E \to F$ et $g : F \to G$ deux applications.
  \begin{enumerate}[label=(\alph*)]
    \item Montrer que si $g \circ f$ est injective, alors $f$ est injective.
    \item Montrer que si $g \circ f$ est surjective, alors $g$ est surjective.
    \item Donner un exemple montrant que $g \circ f$ injective n'implique pas $g$ injective.
  \end{enumerate}
\end{exercice}

\begin{exercice}[Relation d'équivalence]\label{exo:exo-rel-equiv}
  Sur $\Z \times \Z^*$, on définit la relation $(a,b) \mathrel{\mathcal{R}} (c,d) \Leftrightarrow ad = bc$. Montrer que $\mathcal{R}$ est une relation d'équivalence. Interpréter les classes d'équivalence.
\end{exercice}

\begin{exercice}[Défi]\label{exo:exo-defi-ch1}
  Montrer que pour tout entier $n \geq 1$, on a~:
  \[
    \sum_{k=1}^{n} \frac{1}{k(k+1)} = \frac{n}{n+1}.
  \]
  \emph{Indications~: on peut utiliser la récurrence, ou bien la décomposition en éléments simples.}
\end{exercice}


% ============================================================
\section*{Résumé du chapitre}
\addcontentsline{toc}{section}{Résumé du chapitre}
% ============================================================

\begin{resumchapitre}
  \begin{itemize}
    \item Une \textbf{proposition} est un énoncé vrai ou faux. On combine les propositions avec les connecteurs $\lnot$, $\land$, $\lor$, $\Rightarrow$, $\Leftrightarrow$.
    \item Les \textbf{quantificateurs} $\forall$ et $\exists$ transforment des prédicats en propositions. La négation échange $\forall$ et $\exists$.
    \item \textbf{Méthodes de démonstration~:}
      \begin{itemize}
        \item \emph{Preuve directe~:} on suppose $P$ et on déduit $Q$.
        \item \emph{Contraposée~:} on montre $\lnot Q \Rightarrow \lnot P$.
        \item \emph{Absurde~:} on suppose $\lnot P$ et on aboutit à une contradiction.
        \item \emph{Récurrence~:} initialisation + hérédité.
      \end{itemize}
    \item \textbf{Ensembles~:} opérations $\cup$, $\cap$, $\setminus$, complémentaire, lois de De Morgan.
    \item \textbf{Applications~:} injection (éléments distincts ont des images distinctes), surjection (tout élément de l'arrivée est atteint), bijection (les deux à la fois).
    \item \textbf{Relations~:} relation d'équivalence (réflexive, symétrique, transitive) et relation d'ordre (réflexive, antisymétrique, transitive).
  \end{itemize}
\end{resumchapitre}


% ############################################################
%  CHAPITRE 2
% ############################################################
\chapter{Le corps des réels --- Axiomes et complétude}
\label{chap:reels}

\vspace{-0.5cm}
\begin{center}\itshape
  << Dieu a créé les entiers~; tout le reste est l'{\oe}uvre de l'homme. >> --- Leopold Kronecker
\end{center}
\bigskip

Ce chapitre est le c{\oe}ur fondateur de l'analyse réelle. Nous allons construire le cadre dans lequel tout le reste du cours se déploie~: le corps des nombres réels $\R$, muni de son axiome de complétude. C'est cet axiome qui distingue $\R$ de $\Q$ et qui rend possible la théorie des limites.

% ============================================================
\section{Rappels sur $\N$, $\Z$, $\Q$~: pourquoi $\Q$ ne suffit pas}
\label{sec:rappels-Q}
% ============================================================

Les nombres rationnels $\Q$ forment un corps ordonné~: on peut y additionner, multiplier, diviser (sauf par zéro), et comparer. Pourtant, $\Q$ présente un \emph{défaut fondamental}~: il a des << trous >>.

\begin{proposition}[$\sqrt{2}$ n'est pas rationnel]\label{prop:sqrt2-rappel}
  Il n'existe aucun nombre rationnel $r$ tel que $r^2 = 2$.
\end{proposition}

\begin{proof}
  C'est exactement la preuve de l'Exemple~\ref{ex:sqrt2} du Chapitre~\ref{chap:logique}.
\end{proof}

Ce résultat montre que l'équation $x^2 = 2$ n'a pas de solution dans $\Q$. Géométriquement, si l'on place les rationnels sur la droite numérique, il y a un << trou >> à l'endroit où devrait se trouver $\sqrt{2}$. Les nombres réels comblent ces trous.

\begin{remarque}[Insuffisance de $\Q$ pour l'analyse]\label{rem:Q-insuffisant}
  Le problème va plus loin que l'absence de $\sqrt{2}$. De nombreux résultats fondamentaux de l'analyse --- comme le théorème des valeurs intermédiaires --- sont \emph{faux} dans $\Q$. Par exemple, la fonction $f : \Q \to \Q$ définie par $f(x) = x^2 - 2$ change de signe entre $1$ et $2$ (car $f(1) = -1 < 0$ et $f(2) = 2 > 0$), mais elle ne s'annule en aucun point de $\Q$.
\end{remarque}


% ============================================================
\section{Axiomes de corps ordonné}
\label{sec:axiomes}
\index{corps ordonné}
% ============================================================

Les nombres réels $\R$ sont définis axiomatiquement comme un \textbf{corps commutatif totalement ordonné} vérifiant l'\textbf{axiome de la borne supérieure}.

\begin{definition}[Corps commutatif]\label{def:corps}
  \index{corps commutatif}
  Un \textbf{corps commutatif} est un ensemble $K$ muni de deux opérations $+$ (addition) et $\cdot$ (multiplication) vérifiant~:

  \medskip
  \textbf{Axiomes de l'addition~:}
  \begin{enumerate}[label=(A\arabic*)]
    \item \textbf{Associativité~:} $\forall a,b,c \in K,\; (a+b)+c = a+(b+c)$.
    \item \textbf{Commutativité~:} $\forall a,b \in K,\; a+b = b+a$.
    \item \textbf{Élément neutre~:} $\exists\, 0 \in K,\; \forall a \in K,\; a + 0 = a$.
    \item \textbf{Opposé~:} $\forall a \in K,\; \exists\, (-a) \in K,\; a + (-a) = 0$.
  \end{enumerate}

  \medskip
  \textbf{Axiomes de la multiplication~:}
  \begin{enumerate}[label=(M\arabic*)]
    \item \textbf{Associativité~:} $\forall a,b,c \in K,\; (a \cdot b) \cdot c = a \cdot (b \cdot c)$.
    \item \textbf{Commutativité~:} $\forall a,b \in K,\; a \cdot b = b \cdot a$.
    \item \textbf{Élément neutre~:} $\exists\, 1 \in K,\; 1 \neq 0,\; \forall a \in K,\; a \cdot 1 = a$.
    \item \textbf{Inverse~:} $\forall a \in K \setminus \{0\},\; \exists\, a^{-1} \in K,\; a \cdot a^{-1} = 1$.
  \end{enumerate}

  \medskip
  \textbf{Distributivité~:}
  \begin{enumerate}[label=(D)]
    \item $\forall a,b,c \in K,\; a \cdot (b + c) = a \cdot b + a \cdot c$.
  \end{enumerate}
\end{definition}

\begin{definition}[Corps ordonné]\label{def:corps-ordonne}
  \index{corps ordonné}
  Un \textbf{corps ordonné} est un corps commutatif $(K, +, \cdot)$ muni d'une relation d'ordre total $\leq$ compatible avec les opérations~:
  \begin{enumerate}[label=(O\arabic*)]
    \item \textbf{Compatibilité avec l'addition~:} $\forall a,b,c \in K,\; a \leq b \Rightarrow a + c \leq b + c$.
    \item \textbf{Compatibilité avec la multiplication~:} $\forall a,b \in K,\; (0 \leq a \text{ et } 0 \leq b) \Rightarrow 0 \leq a \cdot b$.
  \end{enumerate}
\end{definition}

\begin{remarque}\label{rem:Q-et-R-corps}
  $\Q$ et $\R$ sont tous deux des corps ordonnés. Ce qui distingue $\R$ de $\Q$, c'est l'axiome de la borne supérieure, que nous énoncerons dans la section~\ref{sec:axiome-sup}.
\end{remarque}


% ============================================================
\section{Valeur absolue}
\label{sec:valeur-absolue}
\index{valeur absolue}
% ============================================================

\begin{definition}[Valeur absolue]\label{def:val-abs}
  Pour tout $x \in \R$, la \textbf{valeur absolue} de $x$ est~:
  \[
    \abs{x} = \begin{cases} x & \text{si } x \geq 0, \\ -x & \text{si } x < 0. \end{cases}
  \]
  Géométriquement, $\abs{x}$ est la distance de $x$ à l'origine sur la droite réelle, et $\abs{x - y}$ est la distance entre $x$ et $y$.
\end{definition}

\begin{proposition}[Propriétés de la valeur absolue]\label{prop:prop-val-abs}
  Pour tous $x, y \in \R$ et tout $r \geq 0$~:
  \begin{enumerate}[label=(\roman*)]
    \item $\abs{x} \geq 0$, et $\abs{x} = 0 \Leftrightarrow x = 0$.
    \item $\abs{-x} = \abs{x}$.
    \item $\abs{xy} = \abs{x}\abs{y}$.
    \item $\abs{x} \leq r \Leftrightarrow -r \leq x \leq r$.
    \item $-\abs{x} \leq x \leq \abs{x}$.
  \end{enumerate}
\end{proposition}

\begin{theoreme}[Inégalité triangulaire]\label{thm:ineg-triang}
  \index{inégalité triangulaire}
  Pour tous $x, y \in \R$~:
  \[
    \abs{x + y} \leq \abs{x} + \abs{y}.
  \]
\end{theoreme}

\begin{proof}
  On sait que $-\abs{x} \leq x \leq \abs{x}$ et $-\abs{y} \leq y \leq \abs{y}$ (propriété (v) ci-dessus). En additionnant ces deux encadrements~:
  \[
    -(\abs{x} + \abs{y}) \leq x + y \leq \abs{x} + \abs{y}.
  \]
  Posons $r = \abs{x} + \abs{y} \geq 0$. On a $-r \leq x+y \leq r$, ce qui, par la propriété (iv), équivaut à $\abs{x+y} \leq r = \abs{x} + \abs{y}$.
\end{proof}

\begin{corollaire}[Inégalité triangulaire inversée]\label{cor:ineg-triang-inv}
  Pour tous $x, y \in \R$~:
  \[
    \bigl|\abs{x} - \abs{y}\bigr| \leq \abs{x - y}.
  \]
\end{corollaire}

\begin{proof}
  Par l'inégalité triangulaire appliquée à $x = (x-y) + y$~:
  \[
    \abs{x} = \abs{(x-y) + y} \leq \abs{x-y} + \abs{y},
  \]
  donc $\abs{x} - \abs{y} \leq \abs{x-y}$. En échangeant les rôles de $x$ et $y$~:
  \[
    \abs{y} - \abs{x} \leq \abs{y-x} = \abs{x-y}.
  \]
  Les deux inégalités $\abs{x} - \abs{y} \leq \abs{x-y}$ et $-(\abs{x} - \abs{y}) \leq \abs{x-y}$ donnent, par la propriété (iv), $\bigl|\abs{x} - \abs{y}\bigr| \leq \abs{x-y}$.
\end{proof}


% ============================================================
\section{Parties bornées, borne supérieure, borne inférieure}
\label{sec:bornes}
% ============================================================

\begin{definition}[Majorant, minorant, partie bornée]\label{def:majorant-minorant}
  \index{majorant}\index{minorant}\index{partie bornée}
  Soit $A \subset \R$ une partie non vide.
  \begin{enumerate}[label=(\roman*)]
    \item Un réel $M$ est un \textbf{majorant} de $A$ si~: $\forall a \in A,\; a \leq M$.
    \item Un réel $m$ est un \textbf{minorant} de $A$ si~: $\forall a \in A,\; a \geq m$.
    \item $A$ est \textbf{majorée} si elle admet au moins un majorant. $A$ est \textbf{minorée} si elle admet au moins un minorant. $A$ est \textbf{bornée} si elle est à la fois majorée et minorée.
  \end{enumerate}
\end{definition}

\begin{definition}[Maximum et minimum]\label{def:max-min}
  \index{maximum}\index{minimum}
  Soit $A \subset \R$ non vide.
  \begin{enumerate}[label=(\roman*)]
    \item Si $A$ admet un majorant $M$ qui appartient à $A$, on dit que $M$ est le \textbf{maximum} de $A$~: $M = \max A$.
    \item Si $A$ admet un minorant $m$ qui appartient à $A$, on dit que $m$ est le \textbf{minimum} de $A$~: $m = \min A$.
  \end{enumerate}
\end{definition}

\begin{definition}[Borne supérieure]\label{def:borne-sup}
  \index{borne supérieure}
  Soit $A \subset \R$ une partie non vide et majorée. La \textbf{borne supérieure} de $A$, notée $\sup A$, est le \emph{plus petit des majorants} de $A$. Autrement dit, $s = \sup A$ si et seulement si~:
  \begin{enumerate}[label=(\roman*)]
    \item $s$ est un majorant de $A$~: $\forall a \in A,\; a \leq s$~;
    \item $s$ est le plus petit majorant~: $\forall M \in \R,\; \bigl(\forall a \in A,\; a \leq M\bigr) \Rightarrow s \leq M$.
  \end{enumerate}
  De manière équivalente, la condition (ii) peut être remplacée par~:
  \begin{enumerate}[label=(\roman*')]
    \setcounter{enumi}{1}
    \item tout réel strictement inférieur à $s$ n'est pas un majorant~: $\forall \eps > 0,\; \exists a \in A,\; a > s - \eps$.
  \end{enumerate}
\end{definition}

\begin{remarque}[Distinction cruciale entre sup et max]\label{rem:sup-vs-max}
  Le maximum d'un ensemble, s'il existe, est un élément de l'ensemble. La borne supérieure, elle, peut ne pas appartenir à l'ensemble. Par exemple~:
  \begin{itemize}
    \item $A = \{1, 2, 3\}$~: $\max A = \sup A = 3 \in A$.
    \item $A = {]0, 1[}$~: $\sup A = 1$, mais $1 \notin A$, donc $A$ n'a pas de maximum.
  \end{itemize}
  En revanche, quand le maximum existe, il coïncide avec le sup.
\end{remarque}

\begin{definition}[Borne inférieure]\label{def:borne-inf}
  \index{borne inférieure}
  Soit $A \subset \R$ une partie non vide et minorée. La \textbf{borne inférieure} de $A$, notée $\inf A$, est le \emph{plus grand des minorants} de $A$.
\end{definition}

\begin{theoreme}[Caractérisation séquentielle de la borne supérieure]\label{thm:caract-sup}
  \index{borne supérieure!caractérisation}
  Soit $A \subset \R$ non vide et majorée, et soit $s \in \R$. Alors~:
  \[
    s = \sup A \quad \Longleftrightarrow \quad
    \begin{cases}
      (1)\; \forall a \in A,\; a \leq s, \\[4pt]
      (2)\; \forall \eps > 0,\; \exists a \in A,\; a > s - \eps.
    \end{cases}
  \]
\end{theoreme}

\begin{proof}
  \textbf{Sens $\Rightarrow$~:} Supposons $s = \sup A$.

  \emph{Condition (1)~:} Par définition, $s$ est un majorant de $A$, donc $\forall a \in A,\, a \leq s$.

  \emph{Condition (2)~:} Soit $\eps > 0$. Le réel $s - \eps$ est strictement inférieur à $s$. Or, $s$ est le \emph{plus petit} majorant de $A$, donc $s - \eps$ n'est pas un majorant de $A$. Cela signifie qu'il existe $a \in A$ tel que $a > s - \eps$.

  \medskip
  \textbf{Sens $\Leftarrow$~:} Supposons que $s$ vérifie (1) et (2). Montrons que $s = \sup A$.

  La condition (1) dit que $s$ est un majorant de $A$. Il reste à montrer que $s$ est le plus petit majorant.

  Soit $M$ un majorant quelconque de $A$. Montrons que $s \leq M$. Raisonnons par l'absurde~: supposons $s > M$. Posons $\eps = s - M > 0$. Par la condition (2), il existe $a \in A$ tel que $a > s - \eps = s - (s - M) = M$. Mais cela contredit le fait que $M$ est un majorant de $A$ (car on a trouvé $a \in A$ avec $a > M$).

  Donc $s \leq M$. Comme c'est vrai pour tout majorant $M$, le réel $s$ est le plus petit majorant, c'est-à-dire $s = \sup A$.
\end{proof}

\begin{center}
\begin{tikzpicture}[scale=1.5,>=Stealth]
  % Droite réelle
  \draw[thick,->] (-0.5,0) -- (5,0);
  % Ensemble A
  \draw[very thick,blue] (1,0) -- (3.5,0);
  \fill[blue] (1,0) circle (1.5pt);
  \draw[blue,fill=white] (3.5,0) circle (1.5pt);
  \node[below,blue] at (2.25,-0.15) {$A$};
  % sup
  \draw[thick,red,dashed] (3.5,-0.3) -- (3.5,0.3);
  \node[above,red] at (3.5,0.3) {$\sup A$};
  % s - eps
  \draw[thick,orange,dashed] (2.8,-0.25) -- (2.8,0.25);
  \node[below,orange] at (2.8,-0.25) {$s-\eps$};
  % Accolade
  \draw[decorate,decoration={brace,amplitude=5pt,raise=3pt}] (2.8,0) -- (3.5,0) node[midway,above=8pt] {\small $\exists\, a \in A$ ici};
\end{tikzpicture}
\end{center}


% ============================================================
\section{Axiome de la borne supérieure}
\label{sec:axiome-sup}
\index{axiome de la borne supérieure}
% ============================================================

Nous énonçons maintenant l'axiome qui distingue $\R$ de $\Q$ et qui est la pierre angulaire de toute l'analyse réelle.

\begin{theoreme}[Axiome de la borne supérieure (complétude de $\R$)]\label{thm:completude}
  \index{complétude}
  Toute partie non vide et majorée de $\R$ admet une borne supérieure dans $\R$.
\end{theoreme}

Cet axiome affirme que la droite réelle n'a pas de << trous >>. Contrairement à ce qui se passe dans $\Q$, toute partie bornée de $\R$ possède un sup \emph{dans $\R$}.

\begin{remarque}[Pourquoi $\Q$ ne vérifie pas cet axiome]\label{rem:Q-pas-complet}
  Considérons l'ensemble
  \[
    A = \{q \in \Q : q \geq 0 \text{ et } q^2 < 2\}.
  \]
  Cet ensemble est non vide (par exemple $1 \in A$) et majoré dans $\Q$ (par exemple par $2$). Si la borne supérieure existait dans $\Q$, appelons-la $s$. On pourrait montrer que $s^2 = 2$ (car $s^2 < 2$ et $s^2 > 2$ mènent à des contradictions). Mais nous avons prouvé qu'aucun rationnel ne vérifie $s^2 = 2$. Donc $A$ n'a pas de borne supérieure dans $\Q$.

  En revanche, dans $\R$, la borne supérieure de $A$ est $\sqrt{2}$.
\end{remarque}

\begin{remarque}[Borne inférieure]\label{rem:inf-existe}
  On déduit de l'axiome de la borne supérieure que toute partie non vide et \emph{minorée} de $\R$ admet une borne inférieure dans $\R$. En effet, si $A$ est non vide et minorée, alors $-A = \{-a : a \in A\}$ est non vide et majorée, donc $\sup(-A)$ existe, et on a $\inf A = -\sup(-A)$.
\end{remarque}


% ============================================================
\section{Conséquences fondamentales de la complétude}
\label{sec:consequences}
% ============================================================

L'axiome de la borne supérieure a des conséquences profondes. Nous en démontrons les plus importantes.

% -----------------------------------------------
\subsection{Propriété d'Archimède}
\index{propriété d'Archimède}
% -----------------------------------------------

\begin{theoreme}[Propriété d'Archimède]\label{thm:archimede}
  L'ensemble $\N$ n'est pas majoré dans $\R$. Autrement dit~:
  \[
    \forall x \in \R,\; \exists n \in \N,\; n > x.
  \]
  De manière équivalente~: pour tous $a > 0$ et $b > 0$, il existe $n \in \N$ tel que $na > b$.
\end{theoreme}

\begin{proof}
  Raisonnons par l'absurde. Supposons que $\N$ est majoré dans $\R$. Comme $\N$ est non vide, l'axiome de la borne supérieure garantit l'existence de $s = \sup \N \in \R$.

  Considérons $s - 1$. Puisque $s - 1 < s$ et que $s$ est le \emph{plus petit} majorant de $\N$, le réel $s - 1$ n'est pas un majorant de $\N$. Donc il existe $n_0 \in \N$ tel que $n_0 > s - 1$, c'est-à-dire $n_0 + 1 > s$.

  Or, $n_0 + 1 \in \N$ (car $\N$ est stable par successeur). Donc $n_0 + 1$ est un élément de $\N$ qui dépasse $s$, ce qui contredit le fait que $s$ est un majorant de $\N$.

  Cette contradiction montre que l'hypothèse est fausse~: $\N$ n'est pas majoré.
\end{proof}

\begin{remarque}\label{rem:archimede-utile}
  La propriété d'Archimède sera utilisée constamment~: elle permet d'affirmer que $\frac{1}{n}$ peut être rendu aussi petit que l'on veut (en prenant $n$ assez grand). C'est la base de tous les arguments en $\eps$.
\end{remarque}

% -----------------------------------------------
\subsection{Densité de $\Q$ dans $\R$}
\index{densité de $\Q$ dans $\R$}
% -----------------------------------------------

\begin{theoreme}[Densité de $\Q$ dans $\R$]\label{thm:densite-Q}
  Entre deux réels distincts, il existe toujours un nombre rationnel. Autrement dit~:
  \[
    \forall x, y \in \R,\; x < y \;\Rightarrow\; \exists r \in \Q,\; x < r < y.
  \]
\end{theoreme}

\begin{proof}
  Soient $x, y \in \R$ avec $x < y$. On cherche $r = \frac{p}{q} \in \Q$ tel que $x < \frac{p}{q} < y$, c'est-à-dire $qx < p < qy$.

  \textbf{Étape 1~: choix de $q$.} Comme $y - x > 0$, par la propriété d'Archimède (Théorème~\ref{thm:archimede}), il existe $q \in \N^*$ tel que $q(y - x) > 1$, c'est-à-dire $qy - qx > 1$.

  \textbf{Étape 2~: choix de $p$.} Nous avons $qy > qx + 1$. L'intervalle $]qx, qy[$ a une longueur strictement supérieure à $1$, donc il contient un entier $p \in \Z$. Plus précisément, on peut prendre $p = \lfloor qx \rfloor + 1$ (le plus petit entier strictement supérieur à $qx$). Alors~:
  \begin{itemize}
    \item $p > qx$ par construction (car $p = \lfloor qx \rfloor + 1 > qx$)~;
    \item $p = \lfloor qx \rfloor + 1 \leq qx + 1 < qy$ (car $qy > qx + 1$).
  \end{itemize}

  \textbf{Étape 3~: conclusion.} On a $qx < p < qy$. En divisant par $q > 0$~: $x < \frac{p}{q} < y$. Le nombre $r = \frac{p}{q}$ est rationnel et vérifie $x < r < y$.
\end{proof}

% -----------------------------------------------
\subsection{Densité de $\R \setminus \Q$ dans $\R$}
\index{densité des irrationnels}
% -----------------------------------------------

\begin{theoreme}[Densité des irrationnels]\label{thm:densite-irr}
  Entre deux réels distincts, il existe toujours un nombre irrationnel~:
  \[
    \forall x, y \in \R,\; x < y \;\Rightarrow\; \exists t \in \R \setminus \Q,\; x < t < y.
  \]
\end{theoreme}

\begin{proof}
  Soient $x, y \in \R$ avec $x < y$. Considérons les réels $\frac{x}{\sqrt{2}}$ et $\frac{y}{\sqrt{2}}$. On a $\frac{x}{\sqrt{2}} < \frac{y}{\sqrt{2}}$. Par densité de $\Q$ dans $\R$ (Théorème~\ref{thm:densite-Q}), il existe $r \in \Q$ tel que~:
  \[
    \frac{x}{\sqrt{2}} < r < \frac{y}{\sqrt{2}}.
  \]
  Posons $t = r\sqrt{2}$. Alors $x < t < y$. Si $r = 0$, on a $t = 0 \in \Q$, ce qui ne convient pas~; dans ce cas, on recommence avec $r \neq 0$ (ce qui est possible car l'intervalle $]\frac{x}{\sqrt{2}}, \frac{y}{\sqrt{2}}[$ contient une infinité de rationnels).

  Pour $r \neq 0$~: si $t = r\sqrt{2}$ était rationnel, alors $\sqrt{2} = \frac{t}{r}$ serait rationnel (quotient de deux rationnels), ce qui est absurde. Donc $t$ est irrationnel.
\end{proof}

% -----------------------------------------------
\subsection{Partie entière}
\index{partie entière}
% -----------------------------------------------

\begin{theoreme}[Existence et unicité de la partie entière]\label{thm:partie-entiere}
  Pour tout $x \in \R$, il existe un unique $n \in \Z$ tel que $n \leq x < n + 1$. Cet entier est appelé la \textbf{partie entière} de $x$ et noté $\lfloor x \rfloor$ ou $E(x)$.
\end{theoreme}

\begin{proof}
  \textbf{Existence.}

  \emph{Cas $x \geq 0$~:} Par la propriété d'Archimède, l'ensemble $\{k \in \N : k > x\}$ est non vide. Par le bon ordre de $\N$, il admet un plus petit élément $k_0$. Posons $n = k_0 - 1$. Alors $n \leq x$ (sinon $n > x$ et $n < k_0$, ce qui contredit la minimalité de $k_0$) et $k_0 = n + 1 > x$.

  \emph{Cas $x < 0$~:} On applique le résultat précédent à $-x > 0$~: il existe $m \in \Z$ avec $m \leq -x < m+1$, d'où $-m-1 < x \leq -m$. Si $x = -m$, alors $n = -m$ convient. Sinon $x < -m$ et on pose $n = -m-1$, d'où $n < x < n+1$.

  \textbf{Unicité.} Si $n$ et $n'$ vérifient $n \leq x < n+1$ et $n' \leq x < n'+1$, alors $n < n'+1$ et $n' < n+1$, ce qui donne $\abs{n - n'} < 1$. Comme $n$ et $n'$ sont entiers, on a $n = n'$.
\end{proof}


% ============================================================
\section{Théorème des intervalles emboîtés}
\label{sec:intervalles-emboites}
\index{intervalles emboîtés}
% ============================================================

\begin{theoreme}[Théorème des intervalles emboîtés]\label{thm:intervalles-emboites}
  Soit $(I_n)_{n \in \N}$ une suite d'intervalles fermés bornés $I_n = [a_n, b_n]$ telle que~:
  \begin{enumerate}[label=(\roman*)]
    \item $I_{n+1} \subset I_n$ pour tout $n \in \N$ (les intervalles sont emboîtés)~;
    \item $\lim_{n \to \infty} (b_n - a_n) = 0$ (les longueurs tendent vers $0$).
  \end{enumerate}
  Alors l'intersection $\displaystyle\bigcap_{n=0}^{\infty} I_n$ contient exactement un point.
\end{theoreme}

\begin{proof}
  L'emboîtement $I_{n+1} \subset I_n$ signifie que $a_n \leq a_{n+1}$ et $b_{n+1} \leq b_n$ pour tout $n$. Autrement dit, la suite $(a_n)$ est croissante et la suite $(b_n)$ est décroissante. De plus, $a_n \leq b_n$ pour tout $n$, et même $a_n \leq b_0$ et $b_n \geq a_0$ pour tout $n$.

  \textbf{Existence.} La suite $(a_n)$ est croissante et majorée (par $b_0$). Par l'axiome de la borne supérieure, elle admet une borne supérieure~: posons $\ell = \sup\{a_n : n \in \N\}$.

  Montrons que $\ell \in I_n$ pour tout $n$. Fixons $n \in \N$. On a~:
  \begin{itemize}
    \item $\ell \geq a_n$ car $\ell$ est un majorant de $\{a_k : k \in \N\}$ et $a_n$ est dans cet ensemble~;
    \item $\ell \leq b_n$ car $b_n$ est un majorant de $\{a_k : k \in \N\}$ (en effet, pour tout $k$, si $k \leq n$ alors $a_k \leq a_n \leq b_n$, et si $k > n$ alors $a_k \leq b_k \leq b_n$)~; comme $\ell$ est le \emph{plus petit} majorant, $\ell \leq b_n$.
  \end{itemize}
  Donc $a_n \leq \ell \leq b_n$, c'est-à-dire $\ell \in [a_n, b_n] = I_n$. Ceci étant vrai pour tout $n$, on a $\ell \in \bigcap_{n=0}^{\infty} I_n$.

  \textbf{Unicité.} Supposons qu'il existe $\ell' \in \bigcap_{n=0}^{\infty} I_n$ avec $\ell' \neq \ell$. Pour tout $n$, on a $a_n \leq \ell \leq b_n$ et $a_n \leq \ell' \leq b_n$, donc~:
  \[
    \abs{\ell - \ell'} \leq b_n - a_n.
  \]
  Comme $\lim_{n \to \infty}(b_n - a_n) = 0$, on obtient $\abs{\ell - \ell'} \leq 0$, d'où $\ell = \ell'$. Contradiction.

  L'intersection contient donc exactement un point.
\end{proof}

\begin{center}
\begin{tikzpicture}[scale=1.2,>=Stealth]
  % Axe
  \draw[thick,->] (-0.5,0) -- (7,0);

  % Intervalles emboîtés
  \foreach \n/\a/\b/\y in {0/0.5/6/1.8, 1/1.2/5.2/1.2, 2/1.8/4.5/0.6, 3/2.3/4.0/0.0} {
    \draw[thick,blue] (\a,\y) -- (\b,\y);
    \fill[blue] (\a,\y) circle (1.5pt);
    \fill[blue] (\b,\y) circle (1.5pt);
    \node[left,blue,font=\small] at (\a-0.1,\y) {$a_{\n}$};
    \node[right,blue,font=\small] at (\b+0.1,\y) {$b_{\n}$};
    \node[right,font=\small] at (7,\y) {$I_{\n}$};
  }

  % Point limite
  \node at (3.15,-0.6) {$\vdots$};
  \fill[red] (3.15,-1.0) circle (2pt);
  \node[below,red] at (3.15,-1.1) {$\ell$};
\end{tikzpicture}
\end{center}


% ============================================================
\section*{Remarque historique}
\addcontentsline{toc}{section}{Remarque historique}
% ============================================================

\index{Dedekind}\index{Cantor}
La construction rigoureuse des nombres réels a été accomplie indépendamment par Richard \textbf{Dedekind} (1831--1916) et Georg \textbf{Cantor} (1845--1918) dans les années 1870.

Dedekind a défini les réels par des \textbf{coupures}~: un nombre réel est une partition de $\Q$ en deux ensembles $A$ et $B$ tels que tout élément de $A$ est inférieur à tout élément de $B$. Par exemple, le réel $\sqrt{2}$ correspond à la coupure $A = \{q \in \Q : q < 0 \text{ ou } q^2 < 2\}$, $B = \Q \setminus A$.

Cantor, quant à lui, a construit $\R$ comme l'ensemble des \textbf{classes d'équivalence de suites de Cauchy} de rationnels. Deux suites de Cauchy $(a_n)$ et $(b_n)$ sont équivalentes si $\lim_{n \to \infty}(a_n - b_n) = 0$.

Ces deux constructions aboutissent au même objet~: un corps ordonné complet, unique à isomorphisme près. Dans ce cours, nous adoptons l'approche axiomatique~: nous \emph{postulons} l'existence de $\R$ avec l'axiome de la borne supérieure, et nous en déduisons toutes les conséquences.


% ============================================================
\section*{Erreurs fréquentes du chapitre}
\addcontentsline{toc}{section}{Erreurs fréquentes du chapitre}
% ============================================================

\begin{errmark}
  \begin{enumerate}
    \item \textbf{Confondre max et sup.} Le maximum d'un ensemble est un élément \emph{de} l'ensemble qui majore tous les autres. La borne supérieure est le plus petit majorant, mais elle peut ne pas appartenir à l'ensemble. Par exemple, $\sup\, ]0,1[ \,= 1$, mais $\max\, ]0,1[$ n'existe pas.
    \item \textbf{Croire que la borne sup est toujours dans l'ensemble.} C'est faux~: $\sup\left\{1 - \frac{1}{n} : n \geq 1\right\} = 1$, mais $1$ n'est pas de la forme $1 - \frac{1}{n}$.
    \item \textbf{Confondre << majoré >> et << admet un maximum >>.} L'ensemble $]0,1[$ est majoré (par $1$, ou par $2$, ou par $1000$), mais n'a pas de maximum.
    \item \textbf{Oublier les hypothèses de l'axiome.} L'axiome de la borne supérieure concerne les parties \emph{non vides et majorées}. Il ne dit rien sur les parties non majorées (comme $\N$) ni sur l'ensemble vide.
    \item \textbf{Appliquer l'axiome dans $\Q$.} L'axiome de la borne supérieure est \emph{faux} dans $\Q$. C'est précisément ce qui distingue $\R$ de $\Q$.
  \end{enumerate}
\end{errmark}


% ============================================================
\section*{Exercices}
\addcontentsline{toc}{section}{Exercices}
% ============================================================

\begin{exercice}[Valeur absolue]\label{exo:exo-val-abs}
  Montrer que pour tous $x, y \in \R$~:
  \begin{enumerate}[label=(\alph*)]
    \item $\abs{x - y} \leq \abs{x} + \abs{y}$.
    \item $\abs{\abs{x} - \abs{y}} \leq \abs{x + y}$.
    \item $\abs{x + y}^2 + \abs{x - y}^2 = 2(\abs{x}^2 + \abs{y}^2)$.
  \end{enumerate}
\end{exercice}

\begin{exercice}[Inégalité triangulaire généralisée]\label{exo:exo-ineg-triang}
  Montrer par récurrence que pour tous $x_1, x_2, \ldots, x_n \in \R$~:
  \[
    \abs{x_1 + x_2 + \cdots + x_n} \leq \abs{x_1} + \abs{x_2} + \cdots + \abs{x_n}.
  \]
\end{exercice}

\begin{exercice}[Borne supérieure (calculs)]\label{exo:exo-sup-calculs}
  Déterminer, si elles existent, les bornes supérieure et inférieure des ensembles suivants. Préciser si le max ou le min existe.
  \begin{enumerate}[label=(\alph*)]
    \item $A = \left\{\dfrac{n}{n+1} : n \in \N\right\}$.
    \item $B = \left\{\dfrac{(-1)^n}{n} : n \in \N^*\right\}$.
    \item $C = \left\{\dfrac{1}{n} + \dfrac{1}{m} : n, m \in \N^*\right\}$.
    \item $D = \{x \in \R : x^2 < 3\}$.
  \end{enumerate}
\end{exercice}

\begin{exercice}[Caractérisation du sup]\label{exo:exo-caract-sup}
  Soit $A \subset \R$ non vide et majorée. Montrer que~:
  \[
    s = \sup A \quad \Longleftrightarrow \quad
    s \text{ est un majorant de } A \text{ et } \forall n \in \N^*,\; \exists a_n \in A,\; a_n > s - \tfrac{1}{n}.
  \]
\end{exercice}

\begin{exercice}[Propriétés du sup]\label{exo:exo-prop-sup}
  Soient $A, B \subset \R$ non vides et majorées. On pose $A + B = \{a + b : a \in A,\, b \in B\}$.
  Montrer que $\sup(A + B) = \sup A + \sup B$.
\end{exercice}

\begin{exercice}\label{exo:exo-sup-produit}
  Soient $A, B \subset \R$ non vides avec $A, B \subset \R^+$ (éléments positifs). On pose $A \cdot B = \{ab : a \in A,\, b \in B\}$. Montrer que $\sup(A \cdot B) = \sup A \cdot \sup B$.
\end{exercice}

\begin{exercice}[Densité]\label{exo:exo-densite}
  Montrer qu'entre deux réels distincts, il existe une infinité de rationnels et une infinité d'irrationnels.
\end{exercice}

\begin{exercice}[Partie entière]\label{exo:exo-partie-entiere}
  Soit $x \in \R$ et $n \in \Z$. Montrer que~:
  \begin{enumerate}[label=(\alph*)]
    \item $\lfloor x + n \rfloor = \lfloor x \rfloor + n$.
    \item $\lfloor x \rfloor + \lfloor -x \rfloor = \begin{cases} 0 & \text{si } x \in \Z, \\ -1 & \text{sinon.} \end{cases}$
  \end{enumerate}
\end{exercice}

\begin{exercice}[Application de la propriété d'Archimède]\label{exo:exo-archimede}
  Montrer que $\inf\left\{\dfrac{1}{n} : n \in \N^*\right\} = 0$.
\end{exercice}

\begin{exercice}[Intervalles emboîtés (application)]\label{exo:exo-emboites}
  On pose $I_n = \left[\dfrac{n}{n+1},\, \dfrac{n+1}{n}\right]$ pour $n \geq 1$.
  \begin{enumerate}[label=(\alph*)]
    \item Vérifier que $I_{n+1} \subset I_n$ pour tout $n \geq 1$.
    \item Calculer $\lim_{n \to \infty}(b_n - a_n)$ où $a_n = \frac{n}{n+1}$ et $b_n = \frac{n+1}{n}$.
    \item Déterminer $\bigcap_{n=1}^{\infty} I_n$.
  \end{enumerate}
\end{exercice}

\begin{exercice}[Existence de $\sqrt{2}$ via la complétude]\label{exo:exo-sqrt2-sup}
  On pose $A = \{x \in \R : x \geq 0 \text{ et } x^2 \leq 2\}$.
  \begin{enumerate}[label=(\alph*)]
    \item Montrer que $A$ est non vide et majorée.
    \item Poser $s = \sup A$. Montrer que $s^2 = 2$ en raisonnant par l'absurde sur les cas $s^2 < 2$ et $s^2 > 2$.
  \end{enumerate}
  \emph{Indication pour (b)~: si $s^2 < 2$, montrer qu'on peut trouver $\eps > 0$ tel que $(s+\eps)^2 < 2$, ce qui contredit le fait que $s$ est un majorant de $A$.}
\end{exercice}

\begin{exercice}[Défi]\label{exo:exo-defi-ch2}
  Soit $(a_n)_{n \geq 0}$ la suite définie par $a_0 = 1$ et $a_{n+1} = \frac{1}{2}\left(a_n + \frac{2}{a_n}\right)$.
  \begin{enumerate}[label=(\alph*)]
    \item Montrer que $a_n > 0$ pour tout $n$.
    \item Montrer que $a_n^2 \geq 2$ pour tout $n \geq 1$.
    \item Montrer que la suite $(a_n)_{n \geq 1}$ est décroissante.
    \item En déduire que $(a_n)$ converge, et déterminer sa limite.
  \end{enumerate}
  \emph{(C'est la méthode de Héron pour calculer $\sqrt{2}$.)}
\end{exercice}


% ============================================================
\section*{Résumé du chapitre}
\addcontentsline{toc}{section}{Résumé du chapitre}
% ============================================================

\begin{resumchapitre}
  \begin{itemize}
    \item $\Q$ ne suffit pas pour l'analyse~: il a des << trous >> (par exemple, $\sqrt{2} \notin \Q$).
    \item $\R$ est un \textbf{corps commutatif totalement ordonné complet}. Les axiomes de corps et d'ordre sont partagés avec $\Q$~; ce qui distingue $\R$ est l'\textbf{axiome de la borne supérieure}.
    \item La \textbf{valeur absolue} vérifie l'inégalité triangulaire~: $\abs{x+y} \leq \abs{x} + \abs{y}$.
    \item La \textbf{borne supérieure} d'une partie $A$ est le plus petit majorant. Caractérisation~: $s = \sup A$ si et seulement si $s$ est majorant et $\forall \eps > 0$, $\exists a \in A$, $a > s - \eps$.
    \item \textbf{Axiome de la borne supérieure~:} toute partie non vide et majorée de $\R$ admet un sup dans $\R$.
    \item \textbf{Conséquences~:}
      \begin{itemize}
        \item Propriété d'Archimède~: $\N$ n'est pas majoré.
        \item Densité de $\Q$ dans $\R$~: entre deux réels, il y a toujours un rationnel.
        \item Densité de $\R \setminus \Q$~: entre deux réels, il y a toujours un irrationnel.
        \item Existence et unicité de la partie entière.
      \end{itemize}
    \item \textbf{Intervalles emboîtés~:} une suite décroissante d'intervalles fermés bornés dont les longueurs tendent vers $0$ a une intersection réduite à un point.
  \end{itemize}
\end{resumchapitre}

% ============================================================
%  Analyse Réelle I — Chapitres 3 et 4
%  Licence L1, Mathématiques
%  Partie 2 : Suites numériques, Séries numériques
% ============================================================

% ============================================================
%  CHAPITRE 3 : SUITES NUMÉRIQUES
% ============================================================

\chapter{Suites numériques}\label{chap:suites}
\index{suite numérique}

\section*{Introduction}

Considérons les nombres décimaux successifs qui approchent $\sqrt{2}$ :
\[
  1 ;\quad 1{,}4 ;\quad 1{,}41 ;\quad 1{,}414 ;\quad 1{,}4142 ;\quad \ldots
\]
Chaque terme est plus proche de $\sqrt{2}$ que le précédent. On \emph{sent} que ces
nombres <<~s'approchent~>> de $\sqrt{2}$, mais que signifie cette intuition de manière
rigoureuse ? Comment quantifier le fait qu'une suite de nombres réels <<~tend~>> vers
une valeur ? C'est précisément l'objet de ce chapitre : donner un sens mathématique
précis à la notion de convergence d'une suite, et en déduire un arsenal de résultats
qui constitue le socle de toute l'analyse.

\medskip

L'outil fondamental est la définition dite <<~$\eps$-$N$~>>. Elle transforme l'intuition
vague de <<~s'approcher~>> en une assertion logique vérifiable. Maîtriser cette définition
--- et savoir l'utiliser dans des démonstrations --- est l'objectif principal de ce chapitre.

% ============================================================
\section{Définitions fondamentales}\label{sec:suites:def}

\begin{definition}[Suite réelle]\label{def:suite-reelle}
  \index{suite!définition}
  On appelle \textbf{suite réelle} (ou \textbf{suite numérique}) toute application
  \[
    u : \N \longrightarrow \R, \qquad n \longmapsto u(n) = u_n.
  \]
  On note la suite $(u_n)_{n \in \N}$, ou simplement $(u_n)$. Le réel $u_n$ s'appelle
  le \textbf{terme général} (ou \textbf{terme d'indice $n$}) de la suite.
\end{definition}

\begin{remarque}
  On peut aussi définir une suite à partir d'un rang $n_0 \in \N$ : $(u_n)_{n \geq n_0}$.
  Cela ne change rien aux questions de convergence, car celle-ci ne dépend que du
  comportement <<~à partir d'un certain rang~>>.
\end{remarque}

\begin{exemple}\label{ex:suites-exemples}
  \begin{enumerate}[label=(\roman*)]
    \item $u_n = \dfrac{1}{n+1}$ pour $n \geq 0$. Les premiers termes sont $1, \frac{1}{2}, \frac{1}{3}, \frac{1}{4}, \ldots$
    \item $u_n = (-1)^n$ pour $n \geq 0$. Les premiers termes sont $1, -1, 1, -1, \ldots$
    \item $u_n = \dfrac{n}{n+1}$ pour $n \geq 0$. Les premiers termes sont $0, \frac{1}{2}, \frac{2}{3}, \frac{3}{4}, \ldots$
    \item Suite définie par récurrence : $u_0 = 1$, $u_{n+1} = \dfrac{u_n + 2}{2}$ pour $n \geq 0$.
  \end{enumerate}
\end{exemple}

% ============================================================
\section{Convergence d'une suite}\label{sec:suites:convergence}

Nous arrivons à la définition centrale de ce chapitre. Prenons le temps de bien la comprendre.

\begin{definition}[Convergence d'une suite]\label{def:convergence-suite}
  \index{suite!convergence}\index{convergence!d'une suite}\index{limite!d'une suite}
  Soit $(u_n)_{n \in \N}$ une suite réelle et soit $\ell \in \R$. On dit que la suite
  $(u_n)$ \textbf{converge} vers $\ell$, ou que $\ell$ est la \textbf{limite} de $(u_n)$,
  et l'on écrit
  \[
    \lim_{n \to +\infty} u_n = \ell \qquad \text{ou} \qquad u_n \xrightarrow[n \to +\infty]{} \ell,
  \]
  si la propriété suivante est vérifiée :
  \[
    \boxed{\forall\, \eps > 0, \quad \exists\, N \in \N, \quad \forall\, n \geq N, \quad \abs{u_n - \ell} < \eps.}
  \]
  Une suite qui converge est dite \textbf{convergente}. Une suite qui ne converge pas
  est dite \textbf{divergente}.
\end{definition}

\medskip

\noindent\textbf{Décortiquons cette définition quantificateur par quantificateur :}

\begin{itemize}[leftmargin=2cm]
  \item[\boldmath$\forall\, \eps > 0$] \textbf{<<~Pour tout $\eps > 0$~>>} --- c'est-à-dire
    \emph{quel que soit} le niveau de précision qu'on nous impose, aussi petit soit-il.
    On peut penser à $\eps$ comme une <<~tolérance~>> : on nous demande que les termes
    de la suite soient à distance au plus $\eps$ de $\ell$.

  \item[\boldmath$\exists\, N \in \N$] \textbf{<<~il existe un rang $N$~>>} --- à partir
    duquel la suite reste dans la bande de tolérance. Ce rang $N$ \emph{dépend} de $\eps$ :
    plus $\eps$ est petit, plus $N$ sera grand en général. On écrit parfois $N(\eps)$ pour
    souligner cette dépendance.

  \item[\boldmath$\forall\, n \geq N$] \textbf{<<~pour tout entier $n \geq N$~>>} --- c'est-à-dire
    à partir du rang $N$ et pour \emph{tous} les termes suivants, sans exception.

  \item[\boldmath$\abs{u_n - \ell} < \eps$] \textbf{<<~$u_n$ est à distance strictement
    inférieure à $\eps$ de $\ell$~>>} --- autrement dit $u_n \in \,]\ell - \eps,\, \ell + \eps[\,$.
\end{itemize}

\medskip

\noindent\textbf{En résumé :} la suite converge vers $\ell$ si, quelle que soit la
<<~fenêtre~>> $]\ell - \eps,\, \ell + \eps[$ qu'on place autour de $\ell$, tous les
termes de la suite finissent par tomber dans cette fenêtre (à partir d'un certain rang).

\bigskip

% --- TikZ : bande ε autour de la limite ---
\begin{center}
\begin{tikzpicture}[scale=1.0, >=Stealth]
  % Axes
  \draw[->] (-0.5,0) -- (11,0) node[right] {$n$};
  \draw[->] (0,-0.5) -- (0,4.5) node[above] {$u_n$};

  % Limite
  \def\ell{2.5}
  \def\epsval{0.6}
  \draw[dashed, blue!60!black, thick] (-0.3,\ell) -- (10.8,\ell) node[right] {$\ell$};

  % Bande epsilon
  \fill[blue!8] (-0.3,\ell-\epsval) rectangle (10.8,\ell+\epsval);
  \draw[blue!50!black, densely dashed] (-0.3,\ell+\epsval) -- (10.8,\ell+\epsval);
  \draw[blue!50!black, densely dashed] (-0.3,\ell-\epsval) -- (10.8,\ell-\epsval);
  \draw[<->, blue!60!black] (10.5,\ell) -- node[right] {$\eps$} (10.5,\ell+\epsval);
  \draw[<->, blue!60!black] (10.5,\ell) -- node[right] {$\eps$} (10.5,\ell-\epsval);

  % Rang N
  \def\Nval{4}
  \draw[red!70!black, thick, dashed] (\Nval,-0.3) -- (\Nval,4.3);
  \node[red!70!black, below] at (\Nval,-0.3) {$N$};

  % Termes de la suite
  \foreach \n/\y in {0/0.5, 1/3.8, 2/1.2, 3/3.5, 4/2.9, 5/2.2, 6/2.7, 7/2.35,
    8/2.55, 9/2.45, 10/2.52} {
    \fill[black] (\n, \y) circle (2.5pt);
    \ifnum\n<\Nval
    \else
      \draw[green!50!black, thin] (\n, \y) -- (\n, \ell);
    \fi
  }

  % Annotations
  \node[red!70!black, above] at (2,4.2) {\small termes hors de la bande possibles};
  \draw[->, red!70!black, thin] (2,4.1) -- (1,3.8);
  \node[green!50!black] at (8,4.2) {\small tous dans $]\ell-\eps, \ell+\eps[$};
  \draw[->, green!50!black, thin] (8,4.0) -- (8,2.6);
\end{tikzpicture}
\end{center}

\begin{remarque}\label{rem:eps-N-remarques}
  Quelques observations importantes :
  \begin{enumerate}[label=(\alph*)]
    \item On peut remplacer $\abs{u_n - \ell} < \eps$ par $\abs{u_n - \ell} \leq \eps$
      sans changer la définition (exercice~\ref{exo:eps-large}).
    \item On peut remplacer <<~$\forall\, \eps > 0$~>> par <<~$\forall\, \eps \in \{1/k : k \in \N^*\}$~>>
      sans changer la définition.
    \item Le rang $N$ n'est pas unique : si $N$ convient, tout $N' \geq N$ convient aussi.
    \item Un nombre \emph{fini} de termes de la suite peuvent être arbitrairement loin de $\ell$ :
      seul le comportement <<~à partir d'un certain rang~>> importe.
  \end{enumerate}
\end{remarque}

% ============================================================
\section{Premiers exemples de preuves $\eps$-$N$}\label{sec:suites:exemples-eps-N}

\subsection{Exemple 1 : $\lim_{n \to +\infty} \frac{1}{n} = 0$}
\index{suite!1/n@$1/n$}

C'est l'exemple fondateur. Détaillons la démarche en deux étapes : le \textbf{brouillon}
(recherche de $N$) puis la \textbf{rédaction formelle}.

\medskip

\noindent\textbf{Brouillon (analyse).}
On veut $\abs{\frac{1}{n} - 0} < \eps$, c'est-à-dire $\frac{1}{n} < \eps$ (pour $n \geq 1$).
Or $\frac{1}{n} < \eps \iff n > \frac{1}{\eps}$.
Donc il suffit de choisir $N$ tel que $N > \frac{1}{\eps}$, par exemple
$N = \lfloor \frac{1}{\eps} \rfloor + 1$ (partie entière plus un).

\medskip

\noindent\textbf{Rédaction formelle (synthèse).}

\begin{proof}
  Soit $\eps > 0$. Par la propriété archimédienne de $\R$, il existe $N \in \N^*$
  tel que $N > \frac{1}{\eps}$. Alors, pour tout $n \geq N$, on a $n \geq N > \frac{1}{\eps}$,
  d'où
  \[
    \abs{\frac{1}{n} - 0} = \frac{1}{n} \leq \frac{1}{N} < \eps.
  \]
  Ceci montre que $\lim_{n \to +\infty} \frac{1}{n} = 0$.
\end{proof}

\subsection{Exemple 2 : $\lim_{n \to +\infty} \frac{n}{n+1} = 1$}
\index{suite!n/(n+1)@$n/(n+1)$}

\noindent\textbf{Brouillon.}
On calcule :
\[
  \abs{\frac{n}{n+1} - 1} = \abs{\frac{n - (n+1)}{n+1}} = \abs{\frac{-1}{n+1}} = \frac{1}{n+1}.
\]
On veut $\frac{1}{n+1} < \eps$, soit $n+1 > \frac{1}{\eps}$, soit $n > \frac{1}{\eps} - 1$.
On choisit $N \in \N$ tel que $N > \frac{1}{\eps} - 1$.

\medskip

\noindent\textbf{Rédaction formelle.}

\begin{proof}
  Soit $\eps > 0$. Choisissons $N \in \N$ tel que $N > \frac{1}{\eps} - 1$
  (un tel $N$ existe par la propriété archimédienne). Pour tout $n \geq N$, on a
  $n + 1 > N + 1 > \frac{1}{\eps}$, donc
  \[
    \abs{\frac{n}{n+1} - 1} = \frac{1}{n+1} < \eps.
  \]
  Ainsi $\lim_{n \to +\infty} \frac{n}{n+1} = 1$.
\end{proof}

\subsection{Exemple 3 : $\bigl((-1)^n\bigr)$ diverge}
\index{suite!(-1)\^{}n@$(-1)^n$}

\begin{proposition}\label{prop:moins-un-puiss-n-diverge}
  La suite $((-1)^n)_{n \in \N}$ ne converge pas.
\end{proposition}

\noindent\textbf{Brouillon.}
Supposons par l'absurde qu'elle converge vers $\ell \in \R$.
Avec $\eps = \frac{1}{2}$, à partir d'un rang $N$, tous les termes sont dans
$]\ell - \frac{1}{2},\, \ell + \frac{1}{2}[$. Mais $u_{2k} = 1$ et $u_{2k+1} = -1$
pour tout $k$. Si les deux valeurs sont dans un intervalle de longueur $1$, alors
$\abs{1 - (-1)} = 2 \leq 1$, contradiction.

\begin{proof}
  Supposons par l'absurde que $((-1)^n)$ converge vers un réel $\ell$.
  Prenons $\eps = \frac{1}{2}$. Il existe $N \in \N$ tel que pour tout $n \geq N$,
  $\abs{(-1)^n - \ell} < \frac{1}{2}$.

  Choisissons $n_1 \geq N$ pair et $n_2 \geq N$ impair (par exemple $n_1 = 2\lceil N/2 \rceil$
  et $n_2 = n_1 + 1$). Alors :
  \[
    \abs{1 - \ell} = \abs{(-1)^{n_1} - \ell} < \tfrac{1}{2}
    \qquad \text{et} \qquad
    \abs{-1 - \ell} = \abs{(-1)^{n_2} - \ell} < \tfrac{1}{2}.
  \]
  Par l'inégalité triangulaire :
  \[
    2 = \abs{1 - (-1)} = \abs{(1 - \ell) + (\ell - (-1))}
      \leq \abs{1 - \ell} + \abs{-1 - \ell}
      < \tfrac{1}{2} + \tfrac{1}{2} = 1.
  \]
  On obtient $2 < 1$, ce qui est absurde. Donc $((-1)^n)$ diverge.
\end{proof}

% ============================================================
\section{Propriétés fondamentales}\label{sec:suites:proprietes}

\subsection{Unicité de la limite}

\begin{proposition}[Unicité de la limite]\label{prop:unicite-limite}
  \index{limite!unicité}
  Si une suite $(u_n)$ converge, alors sa limite est unique.
\end{proposition}

\begin{proof}
  Supposons que $(u_n)$ converge vers $\ell$ et vers $\ell'$ avec $\ell \neq \ell'$.
  Posons $\eps = \frac{\abs{\ell - \ell'}}{2} > 0$.

  Par convergence vers $\ell$, il existe $N_1 \in \N$ tel que pour tout $n \geq N_1$,
  $\abs{u_n - \ell} < \eps$.

  Par convergence vers $\ell'$, il existe $N_2 \in \N$ tel que pour tout $n \geq N_2$,
  $\abs{u_n - \ell'} < \eps$.

  Posons $N = \max(N_1, N_2)$. Pour tout $n \geq N$, par l'inégalité triangulaire :
  \[
    \abs{\ell - \ell'} = \abs{(\ell - u_n) + (u_n - \ell')}
      \leq \abs{u_n - \ell} + \abs{u_n - \ell'}
      < \eps + \eps = 2\eps = \abs{\ell - \ell'}.
  \]
  On obtient $\abs{\ell - \ell'} < \abs{\ell - \ell'}$, ce qui est absurde.
  Donc $\ell = \ell'$.
\end{proof}

\subsection{Suites bornées}

\begin{definition}[Suite bornée]\label{def:suite-bornee}
  \index{suite!bornée}
  Une suite $(u_n)$ est dite \textbf{bornée} s'il existe $M \geq 0$ tel que
  $\abs{u_n} \leq M$ pour tout $n \in \N$.
\end{definition}

\begin{proposition}\label{prop:cv-implique-bornee}
  \index{suite!convergente implique bornée}
  Toute suite convergente est bornée.
\end{proposition}

\begin{proof}
  Soit $(u_n)$ une suite convergeant vers $\ell \in \R$. Appliquons la définition de
  convergence avec $\eps = 1$ : il existe $N \in \N$ tel que pour tout $n \geq N$,
  $\abs{u_n - \ell} < 1$, d'où $\abs{u_n} \leq \abs{u_n - \ell} + \abs{\ell} < 1 + \abs{\ell}$.

  Pour les termes de rang $n < N$, ils sont en nombre fini. Posons
  \[
    M = \max\bigl(\abs{u_0},\, \abs{u_1},\, \ldots,\, \abs{u_{N-1}},\, 1 + \abs{\ell}\bigr).
  \]
  Alors pour tout $n \in \N$, $\abs{u_n} \leq M$.
\end{proof}

\begin{remarque}\label{rem:bornee-pas-cv}
  \textbf{Attention !} La réciproque est fausse : la suite $((-1)^n)$ est bornée
  (par $M = 1$) mais ne converge pas (proposition~\ref{prop:moins-un-puiss-n-diverge}).
  Ne confondez \emph{jamais} <<~bornée~>> et <<~convergente~>>.
\end{remarque}

% ============================================================
\section{Théorème des gendarmes}\label{sec:suites:gendarmes}

\begin{theoreme}[Théorème des gendarmes]\label{thm:gendarmes}
  \index{théorème!des gendarmes}\index{encadrement|see{théorème des gendarmes}}
  Soient $(u_n)$, $(v_n)$ et $(w_n)$ trois suites réelles telles que :
  \begin{enumerate}[label=(\roman*)]
    \item à partir d'un certain rang $n_0$, on a $u_n \leq v_n \leq w_n$ ;
    \item $\lim_{n \to +\infty} u_n = \lim_{n \to +\infty} w_n = \ell$.
  \end{enumerate}
  Alors $(v_n)$ converge et $\lim_{n \to +\infty} v_n = \ell$.
\end{theoreme}

\begin{proof}
  Soit $\eps > 0$.

  Puisque $u_n \to \ell$, il existe $N_1 \in \N$ tel que pour tout $n \geq N_1$ :
  $\abs{u_n - \ell} < \eps$, c'est-à-dire $\ell - \eps < u_n < \ell + \eps$.

  Puisque $w_n \to \ell$, il existe $N_2 \in \N$ tel que pour tout $n \geq N_2$ :
  $\abs{w_n - \ell} < \eps$, c'est-à-dire $\ell - \eps < w_n < \ell + \eps$.

  Posons $N = \max(n_0, N_1, N_2)$. Pour tout $n \geq N$, on a :
  \[
    \ell - \eps < u_n \leq v_n \leq w_n < \ell + \eps.
  \]
  Donc $\abs{v_n - \ell} < \eps$. Ceci étant vrai pour tout $\eps > 0$,
  la suite $(v_n)$ converge vers $\ell$.
\end{proof}

\begin{exemple}\label{ex:gendarmes-sin}
  Montrons que $\lim_{n \to +\infty} \frac{\sin(n)}{n} = 0$.

  Pour tout $n \geq 1$, on a $-1 \leq \sin(n) \leq 1$, donc
  $\frac{-1}{n} \leq \frac{\sin(n)}{n} \leq \frac{1}{n}$.
  Or $\lim \frac{-1}{n} = \lim \frac{1}{n} = 0$. Par le théorème des gendarmes,
  $\lim_{n \to +\infty} \frac{\sin(n)}{n} = 0$.
\end{exemple}

% ============================================================
\section{Opérations sur les limites}\label{sec:suites:operations}

\begin{theoreme}[Opérations sur les limites]\label{thm:operations-limites}
  \index{limite!opérations}
  Soient $(u_n)$ et $(v_n)$ deux suites convergentes, de limites respectives $\ell$ et $\ell'$.
  Soit $\lambda \in \R$. Alors :
  \begin{enumerate}[label=(\roman*)]
    \item $(\lambda u_n)$ converge vers $\lambda \ell$ ;
    \item $(u_n + v_n)$ converge vers $\ell + \ell'$ ;
    \item $(u_n \cdot v_n)$ converge vers $\ell \cdot \ell'$ ;
    \item si $\ell' \neq 0$, alors $(u_n / v_n)$ converge vers $\ell / \ell'$
      (à partir d'un certain rang, $v_n \neq 0$).
  \end{enumerate}
\end{theoreme}

\begin{proof}[Démonstration de (ii) : somme]
  Soit $\eps > 0$. Puisque $u_n \to \ell$, il existe $N_1$ tel que pour tout $n \geq N_1$,
  $\abs{u_n - \ell} < \frac{\eps}{2}$.
  Puisque $v_n \to \ell'$, il existe $N_2$ tel que pour tout $n \geq N_2$,
  $\abs{v_n - \ell'} < \frac{\eps}{2}$.

  Posons $N = \max(N_1, N_2)$. Pour tout $n \geq N$ :
  \[
    \abs{(u_n + v_n) - (\ell + \ell')}
      = \abs{(u_n - \ell) + (v_n - \ell')}
      \leq \abs{u_n - \ell} + \abs{v_n - \ell'}
      < \frac{\eps}{2} + \frac{\eps}{2} = \eps.
  \]
\end{proof}

\begin{proof}[Démonstration de (iii) : produit]
  On utilise l'identité :
  \[
    u_n v_n - \ell \ell' = u_n(v_n - \ell') + \ell'(u_n - \ell).
  \]
  Puisque $(u_n)$ converge, elle est bornée (proposition~\ref{prop:cv-implique-bornee}) :
  il existe $M > 0$ tel que $\abs{u_n} \leq M$ pour tout $n$.

  Soit $\eps > 0$. Il existe $N_1$ tel que pour tout $n \geq N_1$,
  $\abs{v_n - \ell'} < \frac{\eps}{2M}$.
  Il existe $N_2$ tel que pour tout $n \geq N_2$,
  $\abs{u_n - \ell} < \frac{\eps}{2(\abs{\ell'} + 1)}$.

  Posons $N = \max(N_1, N_2)$. Pour tout $n \geq N$ :
  \begin{align*}
    \abs{u_n v_n - \ell \ell'}
      &\leq \abs{u_n}\abs{v_n - \ell'} + \abs{\ell'}\abs{u_n - \ell} \\
      &\leq M \cdot \frac{\eps}{2M} + \abs{\ell'} \cdot \frac{\eps}{2(\abs{\ell'} + 1)} \\
      &= \frac{\eps}{2} + \frac{\abs{\ell'}\,\eps}{2(\abs{\ell'} + 1)} \\
      &< \frac{\eps}{2} + \frac{\eps}{2} = \eps.
  \end{align*}
\end{proof}

\begin{remarque}
  La démonstration de (iv) est plus technique ; on commence par montrer que
  $\frac{1}{v_n} \to \frac{1}{\ell'}$ lorsque $\ell' \neq 0$, puis on applique (iii).
  Le point délicat est de montrer que $\abs{v_n} \geq \frac{\abs{\ell'}}{2}$ à partir
  d'un certain rang, ce qui assure que $v_n \neq 0$ et permet de borner $\frac{1}{\abs{v_n}}$.
\end{remarque}

% ============================================================
\section{Suites monotones}\label{sec:suites:monotones}
\index{suite!monotone}

\begin{definition}[Suite monotone]\label{def:suite-monotone}
  Soit $(u_n)$ une suite réelle.
  \begin{itemize}
    \item $(u_n)$ est \textbf{croissante} si $u_{n+1} \geq u_n$ pour tout $n \in \N$.
    \item $(u_n)$ est \textbf{strictement croissante} si $u_{n+1} > u_n$ pour tout $n \in \N$.
    \item $(u_n)$ est \textbf{décroissante} si $u_{n+1} \leq u_n$ pour tout $n \in \N$.
    \item $(u_n)$ est \textbf{monotone} si elle est croissante ou décroissante.
  \end{itemize}
\end{definition}

\begin{theoreme}[Théorème de la limite monotone]\label{thm:limite-monotone}
  \index{théorème!de la limite monotone}\index{suite!monotone bornée}
  Toute suite croissante et majorée converge vers la borne supérieure de son ensemble de
  valeurs. Plus précisément :

  Si $(u_n)$ est croissante et majorée, alors $(u_n)$ converge et
  $\lim_{n \to +\infty} u_n = \sup\{u_n : n \in \N\}$.

  De même, toute suite décroissante et minorée converge vers l'infimum de son ensemble de valeurs.
\end{theoreme}

\begin{proof}
  Supposons $(u_n)$ croissante et majorée. L'ensemble $E = \{u_n : n \in \N\}$ est une
  partie non vide et majorée de $\R$. Par la propriété de la borne supérieure (axiome
  fondamental de $\R$), $\ell = \sup E$ existe dans $\R$.

  Montrons que $u_n \to \ell$. Soit $\eps > 0$. Puisque $\ell - \eps < \ell = \sup E$,
  le réel $\ell - \eps$ n'est pas un majorant de $E$. Il existe donc $N \in \N$ tel que
  $u_N > \ell - \eps$.

  Puisque $(u_n)$ est croissante, pour tout $n \geq N$ :
  \[
    \ell - \eps < u_N \leq u_n \leq \ell < \ell + \eps.
  \]
  Donc $\abs{u_n - \ell} < \eps$. Ceci étant vrai pour tout $\eps > 0$, on conclut que
  $\lim_{n \to +\infty} u_n = \ell$.

  Le cas décroissant minoré se traite de manière analogue avec $\inf$ à la place de $\sup$.
\end{proof}

\begin{remarque}
  Ce théorème repose de manière essentielle sur la \emph{complétude} de $\R$ (propriété de
  la borne supérieure). Il est faux dans $\Q$ : la suite des approximations décimales de
  $\sqrt{2}$ est croissante et majorée dans $\Q$, mais n'a pas de limite dans $\Q$.
\end{remarque}

% --- TikZ : suite monotone bornée ---
\begin{center}
\begin{tikzpicture}[scale=0.9, >=Stealth]
  \draw[->] (-0.5,0) -- (10.5,0) node[right] {$n$};
  \draw[->] (0,-0.5) -- (0,4.5) node[above] {$u_n$};

  % Borne sup
  \def\ell{3.5}
  \draw[dashed, red!70!black, thick] (-0.3,\ell) -- (10.3,\ell)
    node[right] {$\ell = \sup\{u_n\}$};

  % Points croissants
  \foreach \n/\y in {0/0.5, 1/1.3, 2/1.9, 3/2.4, 4/2.8, 5/3.05, 6/3.2,
    7/3.3, 8/3.37, 9/3.42, 10/3.45} {
    \fill[blue!70!black] (\n, \y) circle (2.5pt);
    \ifnum\n>0
      \pgfmathsetmacro{\prevn}{\n-1}
    \fi
  }
  % Relier les points
  \draw[blue!50!black, thin] (0,0.5) -- (1,1.3) -- (2,1.9) -- (3,2.4) --
    (4,2.8) -- (5,3.05) -- (6,3.2) -- (7,3.3) -- (8,3.37) -- (9,3.42) -- (10,3.45);

  \node[blue!70!black, below left] at (0,0.5) {$u_0$};
\end{tikzpicture}
\end{center}

% ============================================================
\section{Suites extraites et Bolzano-Weierstrass}\label{sec:suites:extraites}

\begin{definition}[Suite extraite]\label{def:suite-extraite}
  \index{suite!extraite}\index{sous-suite|see{suite extraite}}
  Soit $(u_n)_{n \in \N}$ une suite et soit $\varphi : \N \to \N$ une application
  strictement croissante. La suite $(u_{\varphi(n)})_{n \in \N}$ est appelée
  \textbf{suite extraite} (ou \textbf{sous-suite}) de $(u_n)$.
\end{definition}

\begin{remarque}
  Puisque $\varphi$ est strictement croissante, on a $\varphi(n) \geq n$ pour tout $n$
  (se démontre par récurrence). Intuitivement, une suite extraite <<~sélectionne~>>
  une infinité de termes de la suite initiale, en respectant l'ordre.
\end{remarque}

\begin{proposition}\label{prop:sous-suite-cv}
  Si $(u_n)$ converge vers $\ell$, alors toute suite extraite $(u_{\varphi(n)})$
  converge également vers $\ell$.
\end{proposition}

\begin{proof}
  Soit $\eps > 0$. Il existe $N$ tel que pour tout $n \geq N$, $\abs{u_n - \ell} < \eps$.
  Puisque $\varphi(n) \geq n$ pour tout $n$, pour tout $n \geq N$ on a $\varphi(n) \geq N$,
  donc $\abs{u_{\varphi(n)} - \ell} < \eps$.
\end{proof}

\begin{definition}[Valeur d'adhérence]\label{def:valeur-adherence}
  \index{valeur d'adhérence}
  Un réel $\ell$ est une \textbf{valeur d'adhérence} de la suite $(u_n)$ s'il est limite
  d'une suite extraite de $(u_n)$.
\end{definition}

\begin{exemple}
  La suite $u_n = (-1)^n$ possède deux valeurs d'adhérence : $1$ (limite de $(u_{2n})$)
  et $-1$ (limite de $(u_{2n+1})$).
\end{exemple}

\begin{theoreme}[Bolzano-Weierstrass]\label{thm:bolzano-weierstrass}
  \index{théorème!de Bolzano-Weierstrass}\index{Bolzano-Weierstrass}
  De toute suite bornée de réels, on peut extraire une sous-suite convergente.

  Autrement dit : toute suite bornée possède au moins une valeur d'adhérence.
\end{theoreme}

\begin{proof}
  Soit $(u_n)$ une suite bornée : il existe $a, b \in \R$ tels que $u_n \in [a, b]$
  pour tout $n \in \N$. On construit par récurrence une suite d'intervalles emboîtés
  et une suite d'indices.

  \textbf{Initialisation.} Posons $I_0 = [a_0, b_0] = [a, b]$.

  \textbf{Hérédité.} Supposons construit un intervalle $I_k = [a_k, b_k]$ de longueur
  $\frac{b - a}{2^k}$ contenant $u_n$ pour une infinité d'indices $n$. Coupons $I_k$
  en deux moitiés :
  \[
    I_k^- = \bigl[a_k,\, \tfrac{a_k + b_k}{2}\bigr]
    \qquad \text{et} \qquad
    I_k^+ = \bigl[\tfrac{a_k + b_k}{2},\, b_k\bigr].
  \]
  Au moins l'une des deux contient $u_n$ pour une infinité d'indices (car la réunion
  de deux ensembles finis est finie). Choisissons $I_{k+1}$ comme étant cette moitié
  (en cas d'ambiguïté, prenons $I_k^-$). On a $\abs{I_{k+1}} = \frac{b - a}{2^{k+1}}$.

  \textbf{Extraction.} On choisit $\varphi(0)$ tel que $u_{\varphi(0)} \in I_0$. Puis,
  par récurrence, on choisit $\varphi(k+1) > \varphi(k)$ tel que $u_{\varphi(k+1)} \in I_{k+1}$
  (possible car $I_{k+1}$ contient une infinité de termes).

  \textbf{Convergence.} Les suites $(a_k)$ et $(b_k)$ vérifient $b_k - a_k = \frac{b-a}{2^k} \to 0$.
  De plus $(a_k)$ est croissante et majorée, $(b_k)$ est décroissante et minorée. Par le
  théorème de la limite monotone, elles convergent vers une même limite $\ell$.

  Pour tout $k \in \N$, on a $a_k \leq u_{\varphi(k)} \leq b_k$. Par le théorème des
  gendarmes, $u_{\varphi(k)} \to \ell$.
\end{proof}

% ============================================================
\section{Suites de Cauchy}\label{sec:suites:cauchy}
\index{suite!de Cauchy}\index{Cauchy!suite de}

\begin{definition}[Suite de Cauchy]\label{def:suite-cauchy}
  Une suite $(u_n)$ est une \textbf{suite de Cauchy} si :
  \[
    \boxed{\forall\, \eps > 0, \quad \exists\, N \in \N, \quad
      \forall\, p, q \geq N, \quad \abs{u_p - u_q} < \eps.}
  \]
\end{definition}

\noindent L'intérêt de cette définition est qu'elle ne fait pas intervenir la limite :
on exprime le fait que les termes se <<~rapprochent les uns des autres~>> sans
avoir besoin de connaître la valeur vers laquelle ils tendent.

% --- TikZ : suite de Cauchy ---
\begin{center}
\begin{tikzpicture}[scale=0.9, >=Stealth]
  \draw[->] (-0.5,0) -- (10.5,0) node[right] {$n$};
  \draw[->] (0,-0.3) -- (0,4.5) node[above] {$u_n$};

  % Points Cauchy
  \foreach \n/\y in {0/1.0, 1/3.2, 2/1.8, 3/3.0, 4/2.3, 5/2.7, 6/2.45,
    7/2.55, 8/2.48, 9/2.52, 10/2.50} {
    \fill[blue!70!black] (\n, \y) circle (2.5pt);
  }
  \draw[blue!50!black, thin] (0,1.0) -- (1,3.2) -- (2,1.8) -- (3,3.0) --
    (4,2.3) -- (5,2.7) -- (6,2.45) -- (7,2.55) -- (8,2.48) -- (9,2.52) -- (10,2.50);

  % Bande pour grands n
  \fill[orange!15] (5,-0.1) rectangle (10.3,3.2);
  \draw[orange!70!black, dashed] (5,-0.1) -- (5,3.2);
  \node[orange!70!black, below] at (5,-0.1) {$N$};

  \draw[<->, orange!70!black, thick] (10.2,2.3) -- node[right] {\small $<\eps$} (10.2,2.7);

  \node[orange!70!black] at (7.5,3.8) {\small les termes se <<\,resserrent\,>>};
\end{tikzpicture}
\end{center}

\begin{proposition}\label{prop:cv-implique-cauchy}
  Toute suite convergente est une suite de Cauchy.
\end{proposition}

\begin{proof}
  Soit $(u_n)$ convergeant vers $\ell$ et soit $\eps > 0$. Il existe $N$ tel que
  pour tout $n \geq N$, $\abs{u_n - \ell} < \frac{\eps}{2}$.

  Pour tous $p, q \geq N$ :
  \[
    \abs{u_p - u_q} = \abs{(u_p - \ell) + (\ell - u_q)}
      \leq \abs{u_p - \ell} + \abs{u_q - \ell}
      < \frac{\eps}{2} + \frac{\eps}{2} = \eps. \qedhere
  \]
\end{proof}

\begin{lemme}\label{lem:cauchy-bornee}
  Toute suite de Cauchy est bornée.
\end{lemme}

\begin{proof}
  Soit $(u_n)$ de Cauchy. Avec $\eps = 1$, il existe $N$ tel que pour tout $p, q \geq N$,
  $\abs{u_p - u_q} < 1$. En particulier, pour tout $n \geq N$, $\abs{u_n - u_N} < 1$,
  donc $\abs{u_n} \leq \abs{u_N} + 1$.

  Posons $M = \max(\abs{u_0}, \ldots, \abs{u_{N-1}}, \abs{u_N} + 1)$. Alors
  $\abs{u_n} \leq M$ pour tout $n \in \N$.
\end{proof}

\begin{theoreme}[Complétude de $\R$]\label{thm:cauchy-R}
  \index{complétude de $\R$}\index{Cauchy!critère de convergence}
  Dans $\R$, une suite est convergente si et seulement si elle est de Cauchy.
\end{theoreme}

\begin{proof}
  $(\Rightarrow)$ C'est la proposition~\ref{prop:cv-implique-cauchy}.

  $(\Leftarrow)$ Soit $(u_n)$ une suite de Cauchy dans $\R$.

  \textit{Étape 1 :} $(u_n)$ est bornée (lemme~\ref{lem:cauchy-bornee}).

  \textit{Étape 2 :} Par le théorème de Bolzano-Weierstrass (théorème~\ref{thm:bolzano-weierstrass}),
  il existe une sous-suite $(u_{\varphi(k)})$ convergeant vers un réel $\ell$.

  \textit{Étape 3 :} Montrons que $(u_n)$ tout entière converge vers $\ell$.
  Soit $\eps > 0$.

  Puisque $(u_n)$ est de Cauchy, il existe $N_1$ tel que pour tous $p, q \geq N_1$ :
  $\abs{u_p - u_q} < \frac{\eps}{2}$.

  Puisque $u_{\varphi(k)} \to \ell$, il existe $K$ tel que pour tout $k \geq K$ :
  $\abs{u_{\varphi(k)} - \ell} < \frac{\eps}{2}$.

  Choisissons $k_0 \geq K$ tel que $\varphi(k_0) \geq N_1$ (possible car $\varphi(k) \to +\infty$).
  Posons $N = N_1$. Pour tout $n \geq N$ :
  \[
    \abs{u_n - \ell} \leq \abs{u_n - u_{\varphi(k_0)}} + \abs{u_{\varphi(k_0)} - \ell}
      < \frac{\eps}{2} + \frac{\eps}{2} = \eps.
  \]
  Donc $u_n \to \ell$.
\end{proof}

% ============================================================
\section{Suites adjacentes}\label{sec:suites:adjacentes}
\index{suites adjacentes}

\begin{definition}[Suites adjacentes]\label{def:suites-adjacentes}
  Deux suites $(u_n)$ et $(v_n)$ sont dites \textbf{adjacentes} si :
  \begin{enumerate}[label=(\roman*)]
    \item $(u_n)$ est croissante ;
    \item $(v_n)$ est décroissante ;
    \item $\lim_{n \to +\infty} (v_n - u_n) = 0$.
  \end{enumerate}
\end{definition}

\begin{theoreme}\label{thm:suites-adjacentes}
  Si $(u_n)$ et $(v_n)$ sont adjacentes, alors elles convergent vers une même limite
  $\ell$ et l'on a $u_n \leq \ell \leq v_n$ pour tout $n$.
\end{theoreme}

\begin{proof}
  Montrons d'abord que $u_n \leq v_n$ pour tout $n$. Posons $w_n = v_n - u_n$.
  La suite $(w_n)$ tend vers $0$. De plus, $w_{n+1} - w_n = (v_{n+1} - v_n) - (u_{n+1} - u_n) \leq 0$,
  donc $(w_n)$ est décroissante. Comme $w_n \to 0$, on a $w_n \geq 0$ pour tout $n$
  (si $w_{n_0} < 0$ pour un certain $n_0$, alors $w_n \leq w_{n_0} < 0$ pour tout $n \geq n_0$,
  contredisant $w_n \to 0$). Donc $u_n \leq v_n$.

  La suite $(u_n)$ est croissante et majorée (par $v_0$, car $u_n \leq v_n \leq v_0$).
  Par le théorème de la limite monotone, $u_n \to \ell_1$.

  La suite $(v_n)$ est décroissante et minorée (par $u_0$). Donc $v_n \to \ell_2$.

  Enfin, $0 = \lim(v_n - u_n) = \ell_2 - \ell_1$, donc $\ell_1 = \ell_2 =: \ell$.

  De $u_n \leq u_{n+1}$ pour tout $n$, on tire à la limite $u_n \leq \ell$.
  De même $\ell \leq v_n$.
\end{proof}

% ============================================================
\section{Suites récurrentes}\label{sec:suites:recurrentes}
\index{suite!récurrente}\index{point fixe}

Une suite récurrente est une suite définie par une relation de la forme
\[
  u_0 \in \R, \qquad u_{n+1} = f(u_n) \quad \text{pour tout } n \geq 0,
\]
où $f : I \to I$ est une fonction définie sur un intervalle $I$.

\begin{proposition}[Point fixe et limite]\label{prop:pt-fixe-limite}
  Si $(u_n)$ est définie par $u_{n+1} = f(u_n)$ avec $f$ continue sur un intervalle $I$,
  et si $(u_n)$ converge vers $\ell \in I$, alors $\ell$ est un \textbf{point fixe}
  de $f$, c'est-à-dire $f(\ell) = \ell$.
\end{proposition}

\begin{proof}
  Puisque $u_{n+1} = f(u_n)$ et que $u_n \to \ell$, on a d'une part $u_{n+1} \to \ell$
  et d'autre part $f(u_n) \to f(\ell)$ par continuité de $f$. Par unicité de la limite,
  $\ell = f(\ell)$.
\end{proof}

\begin{exemple}\label{ex:recurrence-sqrt}
  \textbf{Suite $u_{n+1} = \frac{1}{2}\bigl(u_n + \frac{2}{u_n}\bigr)$, $u_0 = 2$.}

  C'est la méthode de Héron (ou de Babylone) pour approcher $\sqrt{2}$.

  \textit{Étape 1 : Points fixes.} Si $\ell = \frac{1}{2}(\ell + \frac{2}{\ell})$,
  alors $2\ell = \ell + \frac{2}{\ell}$, soit $\ell = \frac{2}{\ell}$, d'où $\ell^2 = 2$
  et $\ell = \sqrt{2}$ (car $u_n > 0$ pour tout $n$).

  \textit{Étape 2 : Montrons que $u_n \geq \sqrt{2}$ pour tout $n \geq 1$.}
  Par l'inégalité arithmético-géométrique :
  $u_{n+1} = \frac{1}{2}\bigl(u_n + \frac{2}{u_n}\bigr) \geq \sqrt{u_n \cdot \frac{2}{u_n}} = \sqrt{2}$.

  \textit{Étape 3 : Montrons que $(u_n)$ est décroissante pour $n \geq 1$.}
  $u_{n+1} - u_n = \frac{1}{2}\bigl(\frac{2}{u_n} - u_n\bigr) = \frac{2 - u_n^2}{2u_n} \leq 0$
  car $u_n \geq \sqrt{2}$ donc $u_n^2 \geq 2$.

  \textit{Conclusion :} $(u_n)_{n \geq 1}$ est décroissante et minorée par $\sqrt{2}$.
  Elle converge, et sa limite est $\sqrt{2}$.
\end{exemple}

\begin{exemple}\label{ex:recurrence-contraction}
  \textbf{Suite $u_{n+1} = \frac{1}{3}u_n + 1$, $u_0 = 0$.}

  \textit{Point fixe :} $\ell = \frac{\ell}{3} + 1 \implies \frac{2\ell}{3} = 1 \implies \ell = \frac{3}{2}$.

  Posons $v_n = u_n - \frac{3}{2}$. Alors $v_{n+1} = u_{n+1} - \frac{3}{2} = \frac{1}{3}u_n + 1 - \frac{3}{2} = \frac{1}{3}(u_n - \frac{3}{2}) = \frac{1}{3}v_n$.

  Donc $v_n = \frac{1}{3^n}v_0 = \frac{1}{3^n}\bigl(-\frac{3}{2}\bigr) \to 0$.
  Ainsi $u_n \to \frac{3}{2}$.
\end{exemple}

% ============================================================
\section{Erreurs courantes}\label{sec:suites:erreurs}

\begin{enumerate}[label=\textbf{E\arabic*.}, leftmargin=2cm]
  \item \textbf{Confondre <<~bornée~>> et <<~convergente~>>.}
    La suite $((-1)^n)$ est bornée mais pas convergente. Une suite convergente est toujours
    bornée, mais la réciproque est fausse.

  \item \textbf{Inverser les quantificateurs.}
    La définition de la convergence est $\forall\, \eps > 0,\, \exists\, N,\, \ldots$
    L'écrire $\exists\, N,\, \forall\, \eps > 0,\, \ldots$ donne une propriété
    \emph{différente} et beaucoup plus forte (et fausse en général).

  \item \textbf{Oublier que $N$ dépend de $\eps$.}
    Dans une preuve $\eps$-$N$, on fixe d'abord $\eps$, \emph{puis} on exhibe $N$.
    Si $N$ ne dépend pas de $\eps$, c'est en général suspect.

  \item \textbf{Déduire la convergence de $u_{n+1} - u_n \to 0$.}
    C'est faux ! Contre-exemple : $u_n = \ln(n)$ vérifie $u_{n+1} - u_n = \ln(1 + 1/n) \to 0$,
    mais $(u_n)$ diverge.

  \item \textbf{Confondre monotone bornée $\Rightarrow$ convergente avec la réciproque.}
    Une suite convergente n'est pas nécessairement monotone : $u_n = (-1)^n/n$ converge
    vers $0$ sans être monotone.
\end{enumerate}

% ============================================================
\section{Résumé du chapitre}\label{sec:suites:resume}

\begin{center}
\fbox{\parbox{0.92\textwidth}{
\textbf{Points clés du chapitre 3}
\begin{itemize}
  \item \textbf{Convergence :} $\forall\, \eps > 0,\ \exists\, N,\ \forall\, n \geq N,\ \abs{u_n - \ell} < \eps$.
  \item La limite, si elle existe, est \textbf{unique}.
  \item Convergente $\Rightarrow$ bornée (la réciproque est \textbf{fausse}).
  \item \textbf{Théorème des gendarmes :} encadrement par deux suites de même limite.
  \item \textbf{Opérations :} somme, produit, quotient de limites.
  \item \textbf{Monotone bornée $\Rightarrow$ convergente} (utilise la complétude de $\R$).
  \item \textbf{Bolzano-Weierstrass :} toute suite bornée admet une sous-suite convergente.
  \item \textbf{Cauchy $\Leftrightarrow$ convergente} dans $\R$ (complétude).
  \item \textbf{Suites adjacentes :} convergent vers une même limite.
  \item \textbf{Suites récurrentes :} chercher les points fixes, montrer la monotonie et le caractère borné.
\end{itemize}
}}
\end{center}

% ============================================================
\section{Exercices}\label{sec:suites:exercices}

\begin{exercice}[$\star$]\label{exo:eps-large}
  Montrer que la définition de la convergence ne change pas si l'on remplace
  <<~$\abs{u_n - \ell} < \eps$~>> par <<~$\abs{u_n - \ell} \leq \eps$~>>.
\end{exercice}

\begin{exercice}[$\star$]\label{exo:cv-1-sur-n2}
  Montrer, à l'aide de la définition $\eps$-$N$, que $\lim_{n \to +\infty} \frac{1}{n^2} = 0$.
  \textit{Indication :} montrer que $\frac{1}{n^2} \leq \frac{1}{n}$ et utiliser l'exemple du cours,
  ou bien exhiber directement $N$.
\end{exercice}

\begin{exercice}[$\star$]\label{exo:cv-2n+1-sur-n}
  Montrer que $\lim_{n \to +\infty} \frac{2n+1}{n+3} = 2$ par la définition $\eps$-$N$.
  Détailler le brouillon et la rédaction.
\end{exercice}

\begin{exercice}[$\star$]\label{exo:suite-bornee-non-cv}
  Donner un exemple de suite bornée non convergente différent de $((-1)^n)$.
\end{exercice}

\begin{exercice}[$\star\star$]\label{exo:produit-bornee-cv0}
  Soit $(u_n)$ une suite bornée et $(v_n)$ une suite convergeant vers $0$.
  Montrer que $(u_n v_n)$ converge vers $0$.
\end{exercice}

\begin{exercice}[$\star\star$]\label{exo:cesaro}
  \textbf{(Lemme de Cesàro.)} Soit $(u_n)$ une suite convergeant vers $\ell$.
  On pose $w_n = \frac{1}{n+1}\sum_{k=0}^{n} u_k$. Montrer que $(w_n)$ converge vers $\ell$.

  \textit{Indication :} décomposer la somme en deux parties en utilisant le rang $N$
  fourni par la convergence de $(u_n)$.
\end{exercice}

\begin{exercice}[$\star\star$]\label{exo:suite-recurrente-cos}
  Soit $u_0 \in [0, 1]$ et $u_{n+1} = \frac{1}{2}(u_n + u_n^2)$ pour $n \geq 0$.
  \begin{enumerate}[label=(\alph*)]
    \item Montrer que $0 \leq u_n \leq 1$ pour tout $n$.
    \item Montrer que $(u_n)$ est croissante.
    \item En déduire que $(u_n)$ converge et déterminer sa limite.
  \end{enumerate}
\end{exercice}

\begin{exercice}[$\star\star$]\label{exo:sqrt-recurrence}
  Soit $u_0 > 0$ et $u_{n+1} = \frac{1}{2}\bigl(u_n + \frac{a}{u_n}\bigr)$, $a > 0$.
  Montrer que $(u_n)$ converge vers $\sqrt{a}$.
\end{exercice}

\begin{exercice}[$\star\star$]\label{exo:bolzano-application}
  Soit $(u_n)$ une suite bornée telle que toute sous-suite convergente a la même limite $\ell$.
  Montrer que $(u_n)$ converge vers $\ell$.

  \textit{Indication :} raisonner par l'absurde et utiliser Bolzano-Weierstrass.
\end{exercice}

\begin{exercice}[$\star\star$]\label{exo:suites-adjacentes-ex}
  On pose $u_n = \sum_{k=0}^{n} \frac{1}{k!}$ et $v_n = u_n + \frac{1}{n \cdot n!}$.
  Montrer que $(u_n)$ et $(v_n)$ sont adjacentes. (Leur limite commune est $e$.)
\end{exercice}

\begin{exercice}[$\star\star\star$]\label{exo:cauchy-pas-cv-Q}
  Donner un exemple de suite de Cauchy dans $\Q$ qui ne converge pas dans $\Q$.
  Que peut-on en conclure sur la complétude de $\Q$ ?
\end{exercice}

\begin{exercice}[$\star\star\star$]\label{exo:stolz-cesaro}
  \textbf{(Stolz-Cesàro.)} Soient $(a_n)$ et $(b_n)$ deux suites réelles telles que
  $(b_n)$ est strictement croissante, non bornée, et
  $\lim_{n \to +\infty} \frac{a_{n+1} - a_n}{b_{n+1} - b_n} = \ell \in \R$.
  Montrer que $\lim_{n \to +\infty} \frac{a_n}{b_n} = \ell$.
\end{exercice}

\begin{exercice}[$\star\star\star$]\label{exo:suite-contractive}
  Une suite $(u_n)$ est dite \textbf{contractive} s'il existe $k \in [0, 1[$ tel que
  $\abs{u_{n+2} - u_{n+1}} \leq k\abs{u_{n+1} - u_n}$ pour tout $n$.
  \begin{enumerate}[label=(\alph*)]
    \item Montrer que toute suite contractive est de Cauchy.
    \item En déduire qu'elle converge.
    \item Application : montrer que la suite définie par $u_0 = 1$,
      $u_{n+1} = \cos(u_n)$ converge. \textit{Indication :} utiliser le théorème des accroissements finis.
  \end{enumerate}
\end{exercice}


% ============================================================
% ============================================================
%  CHAPITRE 4 : SÉRIES NUMÉRIQUES
% ============================================================
% ============================================================

\chapter{Séries numériques}\label{chap:series}
\index{série numérique}

\section*{Introduction}

Que signifie <<~additionner une infinité de nombres~>> ? L'addition ordinaire est une
opération qui porte sur deux nombres, que l'on peut étendre par récurrence à un nombre
fini quelconque de termes. Mais rien ne garantit \emph{a priori} qu'une somme comportant
une infinité de termes ait un sens.

Considérons l'exemple suivant :
\[
  \frac{1}{2} + \frac{1}{4} + \frac{1}{8} + \frac{1}{16} + \cdots
\]
L'intuition suggère que cette somme vaut $1$ : à chaque étape, on comble la moitié de
l'écart restant. Les \textbf{sommes partielles} successives sont
$\frac{1}{2},\; \frac{3}{4},\; \frac{7}{8},\; \frac{15}{16},\; \ldots$
et <<~tendent~>> vers $1$. C'est précisément cette idée --- la convergence des sommes
partielles --- qui fonde la théorie des séries.

Mais attention : des sommes infinies peuvent aussi <<~exploser~>>. La somme
$1 + 1 + 1 + \cdots$ diverge manifestement. Plus subtil : la \textbf{série harmonique}
$1 + \frac{1}{2} + \frac{1}{3} + \frac{1}{4} + \cdots$ diverge également, bien que
ses termes tendent vers $0$. Distinguer les séries convergentes des séries divergentes
est l'objet principal de ce chapitre.

% ============================================================
\section{Définitions fondamentales}\label{sec:series:def}

\begin{definition}[Série et somme partielle]\label{def:serie}
  \index{série!définition}\index{somme partielle}
  Soit $(a_n)_{n \geq 0}$ une suite réelle. On appelle \textbf{série} de terme
  général $a_n$ la suite $(S_n)_{n \geq 0}$ des \textbf{sommes partielles} :
  \[
    S_n = \sum_{k=0}^{n} a_k = a_0 + a_1 + \cdots + a_n.
  \]
  On note la série $\sum a_n$ ou $\sum_{n \geq 0} a_n$.
\end{definition}

\begin{definition}[Convergence d'une série]\label{def:convergence-serie}
  \index{série!convergence}\index{convergence!d'une série}
  On dit que la série $\sum a_n$ \textbf{converge} si la suite $(S_n)$ des sommes
  partielles converge. Dans ce cas, la limite $S = \lim_{n \to +\infty} S_n$ est
  appelée \textbf{somme} de la série et l'on écrit
  \[
    \sum_{n=0}^{+\infty} a_n = S.
  \]
  Si $(S_n)$ diverge, on dit que la série $\sum a_n$ \textbf{diverge}.
\end{definition}

\begin{definition}[Reste d'une série]\label{def:reste-serie}
  \index{reste d'une série}
  Si $\sum a_n$ converge de somme $S$, le \textbf{reste d'ordre $n$} est
  \[
    R_n = S - S_n = \sum_{k=n+1}^{+\infty} a_k.
  \]
  On a $R_n \to 0$ quand $n \to +\infty$.
\end{definition}

% ============================================================
\section{Condition nécessaire de convergence}\label{sec:series:CN}

\begin{theoreme}[Condition nécessaire]\label{thm:CN-series}
  \index{série!condition nécessaire}
  Si la série $\sum a_n$ converge, alors $\lim_{n \to +\infty} a_n = 0$.
\end{theoreme}

\begin{proof}
  Si $\sum a_n$ converge, la suite $(S_n)$ converge vers un réel $S$. Or,
  pour $n \geq 1$,
  \[
    a_n = S_n - S_{n-1}.
  \]
  Donc $\lim_{n \to +\infty} a_n = \lim_{n \to +\infty}(S_n - S_{n-1}) = S - S = 0$.
\end{proof}

\begin{remarque}\label{rem:CN-pas-CS}
  \textbf{Attention !} La réciproque est \textbf{fausse}. C'est l'une des erreurs les
  plus fréquentes en analyse. La série harmonique $\sum \frac{1}{n}$ en est le
  contre-exemple classique.
\end{remarque}

\begin{proposition}[Divergence de la série harmonique]\label{prop:harmonique-diverge}
  \index{série!harmonique}\index{série harmonique}
  La série $\sum_{n=1}^{+\infty} \frac{1}{n}$ diverge.
\end{proposition}

\begin{proof}
  Groupons les termes par paquets de tailles $1, 2, 4, 8, \ldots$ :
  \begin{align*}
    S_{2^k}
      &= 1 + \underbrace{\frac{1}{2}}_{1\text{ terme}}
        + \underbrace{\frac{1}{3} + \frac{1}{4}}_{2\text{ termes}}
        + \underbrace{\frac{1}{5} + \cdots + \frac{1}{8}}_{4\text{ termes}}
        + \cdots
        + \underbrace{\frac{1}{2^{k-1}+1} + \cdots + \frac{1}{2^k}}_{2^{k-1}\text{ termes}}.
  \end{align*}
  Chaque paquet de $2^{j-1}$ termes (pour $j \geq 1$) vérifie
  \[
    \frac{1}{2^{j-1}+1} + \cdots + \frac{1}{2^j} \geq 2^{j-1} \cdot \frac{1}{2^j} = \frac{1}{2}.
  \]
  Donc $S_{2^k} \geq 1 + \frac{k}{2}$, et comme $\frac{k}{2} \to +\infty$,
  la suite $(S_{2^k})$ diverge vers $+\infty$. Puisque $(S_n)$ est croissante
  (les termes sont positifs), $(S_n)$ diverge vers $+\infty$.
\end{proof}

\begin{corollaire}\label{cor:CN-test-divergence}
  \index{test de divergence}
  Si $a_n \not\to 0$, la série $\sum a_n$ diverge (<<~test de divergence grossier~>>).
\end{corollaire}

% ============================================================
\section{Séries géométriques}\label{sec:series:geometriques}
\index{série!géométrique}

\begin{theoreme}[Série géométrique]\label{thm:serie-geometrique}
  Soit $r \in \R$.
  \begin{enumerate}[label=(\roman*)]
    \item Si $\abs{r} < 1$, la série $\sum_{n=0}^{+\infty} r^n$ converge et
      $\displaystyle\sum_{n=0}^{+\infty} r^n = \frac{1}{1 - r}$.
    \item Si $\abs{r} \geq 1$, la série $\sum r^n$ diverge.
  \end{enumerate}
\end{theoreme}

\begin{proof}
  \textit{Calcul de la somme partielle.} Pour $r \neq 1$ :
  \[
    S_n = \sum_{k=0}^{n} r^k = \frac{1 - r^{n+1}}{1 - r}.
  \]
  On le vérifie en multipliant les deux membres par $(1 - r)$ :
  $(1 - r)\sum_{k=0}^{n} r^k = \sum_{k=0}^{n} r^k - \sum_{k=0}^{n} r^{k+1}
  = 1 - r^{n+1}$ (somme télescopique).

  \textit{Cas $\abs{r} < 1$.} On a $\abs{r^{n+1}} = \abs{r}^{n+1} \to 0$, donc
  $S_n \to \frac{1 - 0}{1 - r} = \frac{1}{1 - r}$.

  \textit{Cas $\abs{r} > 1$.} On a $\abs{r^{n+1}} \to +\infty$, donc $(S_n)$ diverge.

  \textit{Cas $r = 1$.} $S_n = n + 1 \to +\infty$.

  \textit{Cas $r = -1$.} $S_n$ alterne entre $1$ et $0$, donc diverge.
\end{proof}

\begin{exemple}
  $\displaystyle\sum_{n=0}^{+\infty} \frac{1}{2^n} = \frac{1}{1 - 1/2} = 2$.
  Cela confirme notre intuition initiale : $\frac{1}{2} + \frac{1}{4} + \frac{1}{8} + \cdots = 1$,
  à quoi s'ajoute le terme $n=0$ qui vaut $1$.
\end{exemple}

% --- TikZ : sommes partielles d'une série géométrique ---
\begin{center}
\begin{tikzpicture}[scale=0.85, >=Stealth]
  \draw[->] (-0.5,0) -- (11,0) node[right] {$n$};
  \draw[->] (0,-0.3) -- (0,3.5) node[above] {$S_n$};

  % Limite
  \draw[dashed, red!70!black] (-0.3,2) -- (10.5,2) node[right] {$S = 2$};

  % Sommes partielles pour r=1/2
  \foreach \n/\y in {0/1.0, 1/1.5, 2/1.75, 3/1.875, 4/1.9375, 5/1.96875,
    6/1.984375, 7/1.992, 8/1.996, 9/1.998, 10/1.999} {
    \fill[blue!70!black] (\n, \y) circle (2.5pt);
  }
  \draw[blue!50!black, thin] (0,1.0) -- (1,1.5) -- (2,1.75) -- (3,1.875) --
    (4,1.9375) -- (5,1.96875) -- (6,1.984375) -- (7,1.992) -- (8,1.996) --
    (9,1.998) -- (10,1.999);

  \node[blue!70!black, left] at (0,1.0) {$S_0=1$};
  \node[blue!70!black, above left] at (2,1.75) {\small $S_2=\frac{7}{4}$};

  \node at (5.5, 0.5) {$\sum \bigl(\frac{1}{2}\bigr)^n$, $r = \frac{1}{2}$};
\end{tikzpicture}
\end{center}

% ============================================================
\section{Séries télescopiques}\label{sec:series:telescopiques}
\index{série!télescopique}

\begin{definition}[Série télescopique]\label{def:serie-telescopique}
  Une série $\sum a_n$ est dite \textbf{télescopique} si $a_n = b_n - b_{n+1}$
  pour une certaine suite $(b_n)$. Alors
  \[
    S_n = \sum_{k=0}^{n} (b_k - b_{k+1}) = b_0 - b_{n+1}.
  \]
  La série converge si et seulement si $(b_n)$ converge, auquel cas
  $\sum_{n=0}^{+\infty} a_n = b_0 - \lim b_n$.
\end{definition}

\begin{exemple}\label{ex:telescopique-1}
  $\displaystyle\sum_{n=1}^{+\infty} \frac{1}{n(n+1)}$.

  On décompose en éléments simples : $\frac{1}{n(n+1)} = \frac{1}{n} - \frac{1}{n+1}$.
  C'est une série télescopique avec $b_n = \frac{1}{n}$. Donc
  $S_N = 1 - \frac{1}{N+1} \to 1$, et $\sum_{n=1}^{+\infty} \frac{1}{n(n+1)} = 1$.
\end{exemple}

\begin{exemple}\label{ex:telescopique-2}
  $\displaystyle\sum_{n=0}^{+\infty} \frac{1}{(2n+1)(2n+3)}$.

  On écrit $\frac{1}{(2n+1)(2n+3)} = \frac{1}{2}\bigl(\frac{1}{2n+1} - \frac{1}{2n+3}\bigr)$.
  Donc $S_N = \frac{1}{2}\bigl(1 - \frac{1}{2N+3}\bigr) \to \frac{1}{2}$.
\end{exemple}

% ============================================================
\section{Séries à termes positifs : critères de convergence}\label{sec:series:termes-positifs}

Lorsque $a_n \geq 0$ pour tout $n$, la suite des sommes partielles est croissante.
Donc, par le théorème de la limite monotone (théorème~\ref{thm:limite-monotone}),
$\sum a_n$ converge si et seulement si $(S_n)$ est majorée.

\subsection{Critère de comparaison}

\begin{theoreme}[Comparaison]\label{thm:comparaison}
  \index{série!critère de comparaison}\index{comparaison (séries)}
  Soient $(a_n)$ et $(b_n)$ deux suites telles que $0 \leq a_n \leq b_n$ à partir
  d'un certain rang.
  \begin{enumerate}[label=(\roman*)]
    \item Si $\sum b_n$ converge, alors $\sum a_n$ converge.
    \item Si $\sum a_n$ diverge, alors $\sum b_n$ diverge.
  \end{enumerate}
\end{theoreme}

\begin{proof}
  (i) Notons $S_n^a = \sum_{k=0}^{n} a_k$ et $S_n^b = \sum_{k=0}^{n} b_k$.
  La suite $(S_n^a)$ est croissante (car $a_n \geq 0$). De plus, pour $n$ assez grand,
  $S_n^a \leq S_n^b \leq \sum_{k=0}^{+\infty} b_k < +\infty$.
  Donc $(S_n^a)$ est croissante et majorée, donc convergente.

  (ii) est la contraposée de (i).
\end{proof}

\subsection{Critère de d'Alembert (rapport)}

\begin{theoreme}[Critère de d'Alembert]\label{thm:dalembert}
  \index{série!critère de d'Alembert}\index{d'Alembert (critère)}
  Soit $\sum a_n$ une série à termes strictement positifs. On suppose que la limite
  \[
    L = \lim_{n \to +\infty} \frac{a_{n+1}}{a_n}
  \]
  existe (finie ou infinie).
  \begin{enumerate}[label=(\roman*)]
    \item Si $L < 1$, la série $\sum a_n$ converge.
    \item Si $L > 1$, la série $\sum a_n$ diverge.
    \item Si $L = 1$, on ne peut pas conclure.
  \end{enumerate}
\end{theoreme}

\begin{proof}
  \textit{Cas $L < 1$.} Choisissons $r$ tel que $L < r < 1$. Posons $\eps = r - L > 0$.
  Il existe $N$ tel que pour tout $n \geq N$, $\abs{\frac{a_{n+1}}{a_n} - L} < \eps$,
  c'est-à-dire $\frac{a_{n+1}}{a_n} < L + \eps = r$.

  Donc, pour tout $n \geq N$ : $a_{n+1} < r \cdot a_n$. Par récurrence, pour tout $n \geq N$ :
  \[
    a_n < r^{n-N} \cdot a_N = \frac{a_N}{r^N} \cdot r^n.
  \]
  Puisque $0 < r < 1$, la série géométrique $\sum r^n$ converge. Par le critère de
  comparaison (théorème~\ref{thm:comparaison}), $\sum a_n$ converge.

  \textit{Cas $L > 1$.} De même, il existe $N$ tel que pour tout $n \geq N$,
  $\frac{a_{n+1}}{a_n} > 1$, donc $(a_n)_{n \geq N}$ est croissante. En particulier,
  $a_n \geq a_N > 0$ pour tout $n \geq N$, donc $a_n \not\to 0$ et la série diverge
  par le test de divergence (corollaire~\ref{cor:CN-test-divergence}).

  \textit{Cas $L = 1$.} $\sum \frac{1}{n}$ diverge et $\sum \frac{1}{n^2}$ converge,
  mais les deux vérifient $\frac{a_{n+1}}{a_n} \to 1$. Donc le critère est non concluant.
\end{proof}

\subsection{Critère de Cauchy (racine)}

\begin{theoreme}[Critère de Cauchy]\label{thm:cauchy-racine}
  \index{série!critère de Cauchy (racine)}\index{Cauchy!critère de la racine}
  Soit $\sum a_n$ une série à termes positifs. On suppose que la limite
  \[
    L = \lim_{n \to +\infty} \sqrt[n]{a_n} = \lim_{n \to +\infty} (a_n)^{1/n}
  \]
  existe.
  \begin{enumerate}[label=(\roman*)]
    \item Si $L < 1$, la série $\sum a_n$ converge.
    \item Si $L > 1$, la série $\sum a_n$ diverge.
    \item Si $L = 1$, on ne peut pas conclure.
  \end{enumerate}
\end{theoreme}

\begin{proof}
  \textit{Cas $L < 1$.} Choisissons $r$ tel que $L < r < 1$. Il existe $N$ tel que pour
  tout $n \geq N$, $\sqrt[n]{a_n} < r$, c'est-à-dire $a_n < r^n$.
  Puisque $\sum r^n$ converge ($0 < r < 1$), la série $\sum a_n$ converge par comparaison.

  \textit{Cas $L > 1$.} Il existe $N$ tel que pour tout $n \geq N$, $\sqrt[n]{a_n} > 1$,
  donc $a_n > 1$. En particulier $a_n \not\to 0$ et la série diverge.
\end{proof}

\subsection{Comparaison intégrale et séries de Riemann}

\begin{theoreme}[Comparaison intégrale]\label{thm:comparaison-integrale}
  \index{série!comparaison intégrale}\index{comparaison intégrale}
  Soit $f : [1, +\infty[ \to [0, +\infty[$ une fonction décroissante. Alors la série
  $\sum_{n=1}^{+\infty} f(n)$ et l'intégrale $\int_1^{+\infty} f(t)\,dt$ sont de même
  nature (toutes deux convergentes ou toutes deux divergentes). Plus précisément :
  \[
    \int_1^{N+1} f(t)\,dt \leq \sum_{n=1}^{N} f(n) \leq f(1) + \int_1^{N} f(t)\,dt.
  \]
\end{theoreme}

\begin{corollaire}[Séries de Riemann]\label{cor:riemann}
  \index{série!de Riemann}\index{Riemann (série de)}
  La série $\sum_{n=1}^{+\infty} \frac{1}{n^\alpha}$ converge si et seulement si $\alpha > 1$.
\end{corollaire}

\begin{proof}
  On applique le théorème de comparaison intégrale avec $f(t) = t^{-\alpha}$.

  Si $\alpha \neq 1$ : $\int_1^N t^{-\alpha}\,dt = \frac{N^{1-\alpha} - 1}{1 - \alpha}$.
  \begin{itemize}
    \item Si $\alpha > 1$ : $1 - \alpha < 0$, donc $N^{1-\alpha} \to 0$ et l'intégrale converge.
    \item Si $\alpha < 1$ : $1 - \alpha > 0$, donc $N^{1-\alpha} \to +\infty$ et l'intégrale diverge.
  \end{itemize}
  Si $\alpha = 1$ : $\int_1^N \frac{dt}{t} = \ln N \to +\infty$, donc divergence
  (c'est la série harmonique).
\end{proof}

% ============================================================
\section{Convergence absolue}\label{sec:series:absolue}
\index{convergence absolue}\index{série!convergence absolue}

\begin{definition}[Convergence absolue]\label{def:convergence-absolue}
  La série $\sum a_n$ est dite \textbf{absolument convergente} si la série
  $\sum \abs{a_n}$ converge.
\end{definition}

\begin{theoreme}\label{thm:abs-implique-cv}
  \index{convergence absolue!implique convergence}
  Toute série absolument convergente est convergente. Autrement dit :
  \[
    \sum \abs{a_n} \text{ converge} \implies \sum a_n \text{ converge.}
  \]
\end{theoreme}

\begin{proof}
  Posons $b_n = a_n + \abs{a_n}$. On a $0 \leq b_n \leq 2\abs{a_n}$ pour tout $n$.

  Si $\sum \abs{a_n}$ converge, alors $\sum 2\abs{a_n}$ converge, et par comparaison
  $\sum b_n$ converge.

  Or $a_n = b_n - \abs{a_n}$, donc $\sum a_n = \sum b_n - \sum \abs{a_n}$ (différence
  de deux séries convergentes). Donc $\sum a_n$ converge.
\end{proof}

\begin{remarque}
  La réciproque est fausse. Une série peut converger sans converger absolument : on dit
  alors qu'elle est \textbf{semi-convergente} (ou conditionnellement convergente).
  L'exemple fondamental est la série alternée $\sum \frac{(-1)^n}{n+1}$.
\end{remarque}

\begin{definition}[Semi-convergence]\label{def:semi-convergence}
  \index{semi-convergence}\index{convergence conditionnelle}
  Une série $\sum a_n$ est dite \textbf{semi-convergente} (ou \textbf{conditionnellement
  convergente}) si elle converge mais ne converge pas absolument.
\end{definition}

% ============================================================
\section{Séries alternées et critère de Leibniz}\label{sec:series:alternees}
\index{série!alternée}\index{Leibniz (critère de)}

\begin{definition}[Série alternée]\label{def:serie-alternee}
  Une série $\sum a_n$ est dite \textbf{alternée} si ses termes sont alternativement
  positifs et négatifs, c'est-à-dire si $a_n = (-1)^n b_n$ avec $b_n \geq 0$ pour tout $n$.
\end{definition}

\begin{theoreme}[Critère de Leibniz]\label{thm:leibniz}
  \index{théorème!de Leibniz}
  Soit $(b_n)_{n \geq 0}$ une suite de réels positifs telle que :
  \begin{enumerate}[label=(\roman*)]
    \item $(b_n)$ est décroissante ;
    \item $\lim_{n \to +\infty} b_n = 0$.
  \end{enumerate}
  Alors la série alternée $\sum_{n=0}^{+\infty} (-1)^n b_n$ converge. De plus, si $S$
  désigne sa somme et $S_n$ la somme partielle d'ordre $n$, on a :
  \[
    \abs{S - S_n} \leq b_{n+1}
    \qquad \text{et} \qquad
    S_{2n+1} \leq S \leq S_{2n} \quad \text{pour tout } n.
  \]
\end{theoreme}

\begin{proof}
  Étudions les sous-suites $(S_{2n})$ et $(S_{2n+1})$.

  \textit{$(S_{2n})$ est décroissante.}
  $S_{2(n+1)} - S_{2n} = a_{2n+1} + a_{2n+2} = -b_{2n+1} + b_{2n+2} = -(b_{2n+1} - b_{2n+2}) \leq 0$
  car $(b_n)$ est décroissante.

  \textit{$(S_{2n+1})$ est croissante.}
  $S_{2(n+1)+1} - S_{2n+1} = a_{2n+2} + a_{2n+3} = b_{2n+2} - b_{2n+3} \geq 0$
  car $(b_n)$ est décroissante.

  \textit{Encadrement.}
  $S_{2n} - S_{2n+1} = a_{2n+1+1}\ldots$ Non, plus directement :
  $S_{2n+1} = S_{2n} - b_{2n+1} \leq S_{2n}$.
  De plus, $S_{2n+1} \geq S_1 = b_0 - b_1 \geq 0$ (si $b_0 \geq b_1$, ce qui est
  le cas car $(b_n)$ est décroissante). Et $S_{2n} \leq S_0 = b_0$.

  Donc $(S_{2n})$ est décroissante et minorée (par $S_1$), et $(S_{2n+1})$ est croissante
  et majorée (par $S_0$). Par le théorème de la limite monotone, elles convergent :
  \[
    S_{2n} \to \ell \qquad \text{et} \qquad S_{2n+1} \to \ell'.
  \]
  Or $S_{2n} - S_{2n+1} = b_{2n+1} \to 0$, donc $\ell = \ell' =: S$.

  Comme $(S_{2n})$ est décroissante de limite $S$ et $(S_{2n+1})$ est croissante de
  limite $S$, on a $S_{2n+1} \leq S \leq S_{2n}$ pour tout $n$.

  Enfin, pour la suite $(S_n)$ tout entière, les deux sous-suites
  $(S_{2n})$ et $(S_{2n+1})$ convergent vers la même limite $S$, donc $S_n \to S$.

  \textit{Majoration du reste.} Pour $n$ pair ($n = 2m$) :
  $\abs{S - S_{2m}} = S_{2m} - S \leq S_{2m} - S_{2m+1} = b_{2m+1} = b_{n+1}$.
  Pour $n$ impair ($n = 2m+1$) :
  $\abs{S - S_{2m+1}} = S - S_{2m+1} \leq S_{2(m+1)} - S_{2m+1} = b_{2m+2} = b_{n+1}$.
  Dans tous les cas, $\abs{S - S_n} \leq b_{n+1}$.
\end{proof}

\begin{exemple}\label{ex:serie-alternee-harmonique}
  La \textbf{série harmonique alternée} $\displaystyle\sum_{n=1}^{+\infty} \frac{(-1)^{n+1}}{n}
  = 1 - \frac{1}{2} + \frac{1}{3} - \frac{1}{4} + \cdots$
  converge par le critère de Leibniz (avec $b_n = \frac{1}{n}$, décroissante et tendant vers $0$).
  Sa somme est $\ln 2$.

  Pourtant, $\sum \frac{1}{n}$ diverge. C'est un exemple de série \emph{semi-convergente}.
\end{exemple}

% ============================================================
\section{Séries entières}\label{sec:series:entieres}
\index{série!entière}\index{série entière}

\begin{definition}[Série entière]\label{def:serie-entiere}
  Une \textbf{série entière} est une série de fonctions de la forme
  \[
    \sum_{n=0}^{+\infty} a_n x^n,
  \]
  où $(a_n)_{n \geq 0}$ est une suite de réels (les \textbf{coefficients}) et $x \in \R$
  est la variable.
\end{definition}

\begin{definition}[Rayon de convergence]\label{def:rayon-convergence}
  \index{rayon de convergence}
  Le \textbf{rayon de convergence} de la série entière $\sum a_n x^n$ est
  \[
    R = \sup\bigl\{\abs{x} \geq 0 : \text{la suite } (a_n x^n) \text{ est bornée}\bigr\}
      \in [0, +\infty].
  \]
  On a :
  \begin{itemize}
    \item la série converge absolument pour $\abs{x} < R$ ;
    \item la série diverge pour $\abs{x} > R$ ;
    \item pour $\abs{x} = R$, il faut étudier au cas par cas.
  \end{itemize}
\end{definition}

% --- TikZ : rayon de convergence ---
\begin{center}
\begin{tikzpicture}[scale=1.2, >=Stealth]
  \draw[->] (-4.5,0) -- (4.5,0) node[right] {$x$};

  % Intervalle de convergence
  \fill[green!15] (-3,-.4) rectangle (3,.4);
  \draw[green!60!black, very thick] (-3,0) -- (3,0);
  \draw[green!60!black, thick, fill=white] (-3,0) circle (3pt);
  \draw[green!60!black, thick, fill=white] (3,0) circle (3pt);

  % Zones de divergence
  \fill[red!10] (-4.5,-.4) rectangle (-3,.4);
  \fill[red!10] (3,-.4) rectangle (4.5,.4);

  \node[green!50!black, above] at (0,0.4) {convergence absolue};
  \node[red!60!black, above] at (-3.75,0.4) {div.};
  \node[red!60!black, above] at (3.75,0.4) {div.};

  \draw[<->, thick] (-3,-0.7) -- node[below] {$R$} (0,-0.7);
  \draw[<->, thick] (0,-0.7) -- node[below] {$R$} (3,-0.7);

  \node[below] at (-3,-0.3) {$-R$};
  \node[below] at (3,-0.3) {$R$};
  \node[below] at (0,-0.15) {$0$};

  \node[orange!70!black] at (-3,0.75) {\small ?};
  \node[orange!70!black] at (3,0.75) {\small ?};
\end{tikzpicture}
\end{center}

\begin{theoreme}[Formule de Hadamard]\label{thm:hadamard}
  \index{Hadamard (formule de)}\index{rayon de convergence!formule de Hadamard}
  Le rayon de convergence de $\sum a_n x^n$ est donné par
  \[
    \frac{1}{R} = \limsup_{n \to +\infty} \abs{a_n}^{1/n},
  \]
  avec les conventions $1/0 = +\infty$ et $1/(+\infty) = 0$.
\end{theoreme}

\begin{remarque}
  En pratique, lorsque la limite $\lim \abs{a_{n+1}/a_n}$ existe, on a aussi
  $R = \lim \abs{a_n / a_{n+1}}$ (conséquence du critère de d'Alembert).
\end{remarque}

\begin{proposition}[Convergence sur tout compact intérieur]\label{prop:cv-uniforme-compacts}
  \index{convergence uniforme}
  Si $\sum a_n x^n$ a un rayon de convergence $R > 0$, alors la série converge
  uniformément sur tout segment $[-r, r]$ avec $0 < r < R$. En particulier, la somme
  $f(x) = \sum_{n=0}^{+\infty} a_n x^n$ est continue sur $]-R, R[$.
\end{proposition}

\subsection{Exemples fondamentaux}

\begin{exemple}\label{ex:serie-exp}
  \index{exponentielle (série)}\index{exp@$\exp$}
  \textbf{Exponentielle.} $\displaystyle\exp(x) = \sum_{n=0}^{+\infty} \frac{x^n}{n!}$,
  rayon de convergence $R = +\infty$.

  En effet, $\abs{\frac{a_{n+1}}{a_n}} = \frac{1}{n+1} \to 0 < 1$ par le critère de d'Alembert.
\end{exemple}

\begin{exemple}\label{ex:serie-sin-cos}
  \index{sinus (série)}\index{cosinus (série)}
  \textbf{Sinus et cosinus.}
  \[
    \sin(x) = \sum_{n=0}^{+\infty} \frac{(-1)^n}{(2n+1)!}\,x^{2n+1}, \qquad
    \cos(x) = \sum_{n=0}^{+\infty} \frac{(-1)^n}{(2n)!}\,x^{2n},
  \]
  toutes deux de rayon de convergence $R = +\infty$.
\end{exemple}

\begin{exemple}\label{ex:serie-ln}
  \index{logarithme (série)}
  \textbf{Logarithme.} $\displaystyle\ln(1+x) = \sum_{n=1}^{+\infty} \frac{(-1)^{n+1}}{n}\,x^n$,
  rayon de convergence $R = 1$, valable pour $x \in \,]-1, 1]$.

  Le rayon se calcule par $\abs{a_{n+1}/a_n} = \frac{n}{n+1} \to 1$, donc $R = 1$.
  En $x = 1$, la série converge (série harmonique alternée). En $x = -1$, c'est $-\sum \frac{1}{n}$
  qui diverge.
\end{exemple}

\begin{exemple}\label{ex:serie-binomiale}
  \index{série binomiale}
  \textbf{Série binomiale.} Pour $\alpha \in \R$ :
  \[
    (1+x)^\alpha = \sum_{n=0}^{+\infty} \binom{\alpha}{n} x^n, \qquad \abs{x} < 1,
  \]
  où $\binom{\alpha}{n} = \frac{\alpha(\alpha-1)\cdots(\alpha-n+1)}{n!}$ est le coefficient
  binomial généralisé. Le rayon de convergence est $R = 1$.
\end{exemple}

% ============================================================
\section{Produit de Cauchy de séries}\label{sec:series:produit-cauchy}
\index{produit de Cauchy}\index{Cauchy!produit de séries}

\begin{definition}[Produit de Cauchy]\label{def:produit-cauchy}
  Soient $\sum a_n$ et $\sum b_n$ deux séries. Leur \textbf{produit de Cauchy} est la
  série $\sum c_n$ où
  \[
    c_n = \sum_{k=0}^{n} a_k \, b_{n-k} = a_0 b_n + a_1 b_{n-1} + \cdots + a_n b_0.
  \]
\end{definition}

\begin{theoreme}[Mertens]\label{thm:mertens}
  \index{théorème!de Mertens}
  Si $\sum a_n$ converge absolument (de somme $A$) et $\sum b_n$ converge (de somme $B$),
  alors le produit de Cauchy $\sum c_n$ converge et sa somme est $AB$.
\end{theoreme}

\begin{exemple}\label{ex:exp-produit}
  \textbf{$\exp(x) \cdot \exp(y) = \exp(x+y)$.}

  Les séries $\sum \frac{x^n}{n!}$ et $\sum \frac{y^n}{n!}$ sont absolument convergentes.
  Leur produit de Cauchy a pour terme général :
  \[
    c_n = \sum_{k=0}^{n} \frac{x^k}{k!} \cdot \frac{y^{n-k}}{(n-k)!}
        = \frac{1}{n!} \sum_{k=0}^{n} \binom{n}{k} x^k y^{n-k}
        = \frac{(x+y)^n}{n!}.
  \]
  Par le théorème de Mertens, $\exp(x) \cdot \exp(y) = \sum_{n=0}^{+\infty} \frac{(x+y)^n}{n!} = \exp(x+y)$.
\end{exemple}

% ============================================================
\section{Erreurs courantes}\label{sec:series:erreurs}

\begin{enumerate}[label=\textbf{E\arabic*.}, leftmargin=2cm]
  \item \textbf{<<~$a_n \to 0$ donc $\sum a_n$ converge.~>>}
    \textsc{Faux !} C'est la confusion entre condition nécessaire et condition suffisante.
    Contre-exemple : $\frac{1}{n} \to 0$ mais $\sum \frac{1}{n}$ diverge.

  \item \textbf{Appliquer le critère de d'Alembert ou de Cauchy lorsque $L = 1$.}
    Ces critères sont \textbf{non concluants} pour $L = 1$. Il faut alors utiliser un
    autre outil (comparaison, Leibniz, séries de Riemann, etc.).

  \item \textbf{Confondre convergence absolue et convergence.}
    La convergence absolue est plus forte. <<~$\sum a_n$ converge~>> ne signifie
    \emph{pas} que $\sum \abs{a_n}$ converge.

  \item \textbf{Oublier de vérifier les hypothèses de Leibniz.}
    Le critère de Leibniz exige que $(b_n)$ soit décroissante \emph{et} tende vers $0$.
    Si l'une de ces conditions manque, on ne peut pas conclure.

  \item \textbf{Confondre le terme général $a_n$ et la somme partielle $S_n$.}
    La série $\sum a_n$ converge signifie que $S_n = \sum_{k=0}^n a_k$ a une limite
    finie, \emph{pas} que $a_n$ a une limite finie (même si c'est une conséquence).
\end{enumerate}

% ============================================================
\section{Résumé du chapitre}\label{sec:series:resume}

\begin{center}
\fbox{\parbox{0.92\textwidth}{
\textbf{Points clés du chapitre 4}
\begin{itemize}
  \item $\sum a_n$ converge $\iff$ la suite des sommes partielles $S_n$ converge.
  \item \textbf{Condition nécessaire :} $\sum a_n$ cv $\Rightarrow$ $a_n \to 0$
    (réciproque \textbf{fausse}).
  \item \textbf{Série géométrique :} $\sum r^n = \frac{1}{1-r}$ si $\abs{r} < 1$.
  \item \textbf{Termes positifs :} comparaison, d'Alembert ($L < 1$ cv, $L > 1$ div),
    Cauchy racine, comparaison intégrale.
  \item \textbf{Séries de Riemann :} $\sum 1/n^\alpha$ cv $\iff$ $\alpha > 1$.
  \item \textbf{Convergence absolue $\Rightarrow$ convergence} (réciproque fausse).
  \item \textbf{Leibniz :} $\sum (-1)^n b_n$ cv si $b_n \searrow 0$, avec $\abs{R_n} \leq b_{n+1}$.
  \item \textbf{Séries entières :} rayon de convergence $R$, cv absolue pour $\abs{x} < R$.
  \item \textbf{Produit de Cauchy :} $(\sum a_n)(\sum b_n) = \sum c_n$ sous convergence absolue.
\end{itemize}
}}
\end{center}

% ============================================================
\section{Exercices}\label{sec:series:exercices}

\begin{exercice}[$\star$]\label{exo:serie-geo-somme}
  Calculer $\displaystyle\sum_{n=0}^{+\infty} \frac{3}{4^n}$
  et $\displaystyle\sum_{n=1}^{+\infty} \frac{(-1)^n}{3^n}$.
\end{exercice}

\begin{exercice}[$\star$]\label{exo:CN-divergence}
  Montrer que les séries suivantes divergent :
  $\sum \frac{n}{n+1}$, \quad $\sum \cos\bigl(\frac{1}{n}\bigr)$, \quad $\sum (-1)^n$.
\end{exercice}

\begin{exercice}[$\star$]\label{exo:telescopique-exo}
  Calculer $\displaystyle\sum_{n=1}^{+\infty} \frac{1}{n(n+2)}$.
  \textit{Indication :} décomposer en éléments simples.
\end{exercice}

\begin{exercice}[$\star$]\label{exo:telescopique-sqrt}
  Montrer que $\displaystyle\sum_{n=1}^{+\infty} \bigl(\sqrt{n+1} - \sqrt{n}\bigr)$ diverge.
\end{exercice}

\begin{exercice}[$\star\star$]\label{exo:dalembert-facto}
  Étudier la convergence de $\displaystyle\sum_{n=0}^{+\infty} \frac{n!}{n^n}$
  à l'aide du critère de d'Alembert.
\end{exercice}

\begin{exercice}[$\star\star$]\label{exo:cauchy-racine-exo}
  Étudier la convergence de $\displaystyle\sum_{n=1}^{+\infty} \Bigl(\frac{n}{2n+1}\Bigr)^n$
  à l'aide du critère de Cauchy (racine).
\end{exercice}

\begin{exercice}[$\star\star$]\label{exo:riemann-comparaison}
  En utilisant une comparaison avec une série de Riemann, montrer que
  $\displaystyle\sum_{n=1}^{+\infty} \frac{1}{n^2 + n + 1}$ converge.
\end{exercice}

\begin{exercice}[$\star\star$]\label{exo:leibniz-application}
  Montrer que la série $\displaystyle\sum_{n=1}^{+\infty} \frac{(-1)^{n+1}}{n^2}$ converge.
  Est-elle absolument convergente ?
\end{exercice}

\begin{exercice}[$\star\star$]\label{exo:abs-cv-non-reciproque}
  Montrer que $\displaystyle\sum_{n=1}^{+\infty} \frac{(-1)^{n+1}}{\sqrt{n}}$ converge
  mais n'est pas absolument convergente.
\end{exercice}

\begin{exercice}[$\star\star$]\label{exo:rayon-cv}
  Déterminer le rayon de convergence des séries entières suivantes :
  \begin{enumerate}[label=(\alph*)]
    \item $\displaystyle\sum n!\, x^n$ \qquad
    \item $\displaystyle\sum \frac{x^n}{n^2}$ \qquad
    \item $\displaystyle\sum \frac{n^n}{n!}\, x^n$ \qquad
    \item $\displaystyle\sum \frac{x^{2n}}{4^n}$
  \end{enumerate}
\end{exercice}

\begin{exercice}[$\star\star\star$]\label{exo:series-bertrand}
  \textbf{(Séries de Bertrand.)} Étudier la convergence de
  $\displaystyle\sum_{n=2}^{+\infty} \frac{1}{n^\alpha (\ln n)^\beta}$
  en fonction de $\alpha$ et $\beta$.

  \textit{Indication :} distinguer les cas $\alpha > 1$, $\alpha < 1$ et $\alpha = 1$.
\end{exercice}

\begin{exercice}[$\star\star\star$]\label{exo:produit-cauchy-exo}
  Soit $f(x) = \sum_{n=0}^{+\infty} \frac{x^n}{n!}$ pour $x \in \R$.
  En utilisant le produit de Cauchy, montrer directement que $f(x)f(-x) = 1$ pour tout $x$.
  En déduire que $f(x) \neq 0$ pour tout $x$.
\end{exercice}

\begin{exercice}[$\star\star\star$]\label{exo:abel}
  \textbf{(Transformation d'Abel.)} Soient $(a_n)$ et $(b_n)$ deux suites. On pose
  $B_n = \sum_{k=0}^{n} b_k$. Montrer la formule de sommation par parties :
  \[
    \sum_{k=0}^{n} a_k b_k = a_n B_n - \sum_{k=0}^{n-1} (a_{k+1} - a_k) B_k.
  \]
  En déduire le \textbf{critère d'Abel} : si $(a_n)$ est monotone convergente vers $0$
  et si les sommes partielles $B_n$ sont bornées, alors $\sum a_n b_n$ converge.
\end{exercice}

\begin{exercice}[$\star\star\star$]\label{exo:rearrangement}
  \textbf{(Théorème de Riemann sur les réarrangements.)}
  Soit $\sum a_n$ une série semi-convergente. On admet que pour tout $S \in \R \cup \{-\infty, +\infty\}$,
  il existe une bijection $\sigma : \N \to \N$ telle que $\sum a_{\sigma(n)} = S$.

  Vérifier ce résultat pour $a_n = \frac{(-1)^{n+1}}{n}$ : trouver un réarrangement
  dont la somme vaut $\frac{3}{2}\ln 2$.

  \textit{Indication :} prendre deux termes positifs, puis un terme négatif, et répéter.
\end{exercice}
% ============================================================
%  Analyse Réelle I — Partie 3 (Chapitres 5–6)
%  Licence L1, Mathématiques
% ============================================================

% ============================================================
%  CHAPITRE 5 : Limites et Continuité des Fonctions
% ============================================================
\chapter{Limites et Continuité des Fonctions}\label{chap:limites-continuite}
\index{continuité}\index{limite!de fonction}

\section*{Introduction}

Dans les chapitres précédents, nous avons étudié les suites réelles et leurs limites. Nous passons maintenant à l'étude des \emph{fonctions} $f \colon D \to \R$, où $D \subset \R$. La question centrale est la suivante : comment décrire le comportement de $f(x)$ lorsque $x$ s'approche d'un point~$a$ ? La réponse repose sur le formalisme $\varepsilon$-$\delta$ dû à Weierstrass, qui fournit un cadre rigoureux pour définir les limites de fonctions et la continuité. Ce formalisme est le socle de toute l'analyse moderne ; il est \emph{indispensable} de le maîtriser dès la première année.

% ============================================================
\section{Limite d'une fonction en un point}\label{sec:limite-fonction}

% --- Définition ε-δ ---
\begin{definition}[Limite d'une fonction en un point]\label{def:limite-fonction}
\index{limite!d'une fonction en un point}
Soient $D \subset \R$, $f \colon D \to \R$, $a \in \R$ un point \emph{adhérent} à $D \setminus \{a\}$ et $\ell \in \R$. On dit que $f$ admet pour \textbf{limite} $\ell$ lorsque $x$ tend vers $a$, et l'on écrit
\[
  \lim_{x \to a} f(x) = \ell,
\]
si et seulement si :
\[
  \forall\, \varepsilon > 0,\quad \exists\, \delta > 0,\quad
  \forall\, x \in D,\quad
  0 < |x - a| < \delta \;\Longrightarrow\; |f(x) - \ell| < \varepsilon.
\]

\medskip
\textbf{Décryptage.} Pour tout «~couloir~» de hauteur $2\varepsilon$ autour de $\ell$ sur l'axe des ordonnées, on peut trouver un «~couloir~» de largeur $2\delta$ autour de $a$ sur l'axe des abscisses tel que, dès que $x$ est dans ce couloir (sans être égal à $a$), l'image $f(x)$ tombe dans le couloir vertical. L'ordre des quantificateurs est crucial : $\varepsilon$ vient en premier (c'est l'adversaire qui choisit la précision), puis $\delta$ est notre réponse.
\end{definition}

\begin{remarque}\label{rem:limite-valeur-en-a}
La définition ne fait \emph{aucune} hypothèse sur $f(a)$ : le point $a$ peut ne pas appartenir à $D$, ou $f(a)$ peut différer de $\ell$. Seul le comportement de $f$ \emph{au voisinage} de $a$ importe.
\end{remarque}

% --- Caractérisation séquentielle ---
\begin{theoreme}[Caractérisation séquentielle de Heine]\label{thm:heine-limite}
\index{Heine!caractérisation séquentielle}
Soient $D \subset \R$, $f \colon D \to \R$, $a$ adhérent à $D \setminus \{a\}$ et $\ell \in \R$. Alors
\[
  \lim_{x \to a} f(x) = \ell
  \quad\Longleftrightarrow\quad
  \text{pour toute suite } (x_n) \subset D \setminus \{a\} \text{ telle que } x_n \to a,\;
  f(x_n) \to \ell.
\]
\end{theoreme}

\begin{proof}
\textbf{($\Rightarrow$)} Supposons $\lim_{x \to a} f(x) = \ell$. Soit $(x_n)$ une suite d'éléments de $D \setminus \{a\}$ telle que $x_n \to a$. Soit $\varepsilon > 0$. Par hypothèse, il existe $\delta > 0$ tel que
\[
  \forall\, x \in D,\quad 0 < |x - a| < \delta \;\Longrightarrow\; |f(x) - \ell| < \varepsilon.
\]
Puisque $x_n \to a$, il existe $N \in \N$ tel que pour tout $n \geqslant N$, $|x_n - a| < \delta$. Comme $x_n \neq a$ pour tout $n$, on a $0 < |x_n - a| < \delta$ pour $n \geqslant N$, d'où $|f(x_n) - \ell| < \varepsilon$. Ceci montre que $f(x_n) \to \ell$.

\medskip
\textbf{($\Leftarrow$)} Nous procédons par contraposée. Supposons que $\lim_{x \to a} f(x) \neq \ell$. Cela signifie :
\[
  \exists\, \varepsilon_0 > 0,\quad \forall\, \delta > 0,\quad
  \exists\, x \in D,\quad 0 < |x - a| < \delta \;\text{et}\; |f(x) - \ell| \geqslant \varepsilon_0.
\]
Appliquons cette propriété avec $\delta = 1/n$ pour chaque $n \geqslant 1$ : il existe $x_n \in D$ tel que $0 < |x_n - a| < 1/n$ et $|f(x_n) - \ell| \geqslant \varepsilon_0$. On obtient ainsi une suite $(x_n)$ d'éléments de $D \setminus \{a\}$ telle que $x_n \to a$ (car $|x_n - a| < 1/n \to 0$) mais $f(x_n) \not\to \ell$ (car $|f(x_n) - \ell| \geqslant \varepsilon_0$ pour tout $n$). La contraposée est établie.
\end{proof}

% --- TikZ : diagramme ε-δ ---
\begin{figure}[ht]
\centering
\begin{tikzpicture}[scale=1.1, >=Stealth]
  % Axes
  \draw[->] (-0.5,0) -- (6,0) node[right]{$x$};
  \draw[->] (0,-0.5) -- (0,5) node[above]{$y$};

  % Fonction
  \draw[blue, thick, domain=0.3:5.5, smooth, samples=100]
    plot(\x, {0.4*\x*\x - 1.6*\x + 3.5});

  % Point a et ℓ
  \def\aval{3}
  \def\lval{2.3}
  \def\epsval{0.6}
  \def\deltaval{0.75}

  % Bandes ε
  \fill[red!10] (-0.3,\lval-\epsval) rectangle (5.8,\lval+\epsval);
  \draw[red, dashed] (-0.3,\lval-\epsval) -- (5.8,\lval-\epsval);
  \draw[red, dashed] (-0.3,\lval+\epsval) -- (5.8,\lval+\epsval);

  % Bandes δ
  \fill[green!10] (\aval-\deltaval,0) rectangle (\aval+\deltaval,4.8);
  \draw[green!60!black, dashed] (\aval-\deltaval,0) -- (\aval-\deltaval,4.8);
  \draw[green!60!black, dashed] (\aval+\deltaval,0) -- (\aval+\deltaval,4.8);

  % Labels
  \draw[<->] (5.9,\lval-\epsval) -- node[right]{$2\varepsilon$} (5.9,\lval+\epsval);
  \draw[<->] (\aval-\deltaval,-0.35) -- node[below]{$2\delta$} (\aval+\deltaval,-0.35);

  \node[below] at (\aval,0) {$a$};
  \node[left] at (0,\lval) {$\ell$};

  % Point creux en a
  \draw[blue, fill=white] (\aval,\lval) circle (2pt);
\end{tikzpicture}
\caption{Illustration de la définition $\varepsilon$-$\delta$ de la limite.
La bande horizontale (rouge) correspond à $]\ell - \varepsilon, \ell + \varepsilon[$,
la bande verticale (verte) à $]a - \delta, a + \delta[$.}
\label{fig:epsilon-delta-limite}
\end{figure}

% --- Exemple travaillé 1 ---
\begin{exemple}\label{ex:limite-3x-1}
\textbf{Montrer que} $\displaystyle\lim_{x \to 2} (3x - 1) = 5$.

\medskip
\textbf{Brouillon (travail préparatoire).} On veut $|f(x) - \ell| < \varepsilon$, c'est-à-dire $|(3x-1) - 5| < \varepsilon$, soit $|3x - 6| < \varepsilon$, soit $3|x-2| < \varepsilon$, soit $|x - 2| < \varepsilon/3$. On prendra donc $\delta = \varepsilon/3$.

\medskip
\textbf{Rédaction.} Soit $\varepsilon > 0$. Posons $\delta = \varepsilon/3 > 0$. Soit $x \in \R$ tel que $0 < |x - 2| < \delta$. Alors :
\[
  |f(x) - 5| = |(3x - 1) - 5| = |3x - 6| = 3|x - 2| < 3\delta = 3 \cdot \frac{\varepsilon}{3} = \varepsilon.
\]
Donc $\lim_{x \to 2}(3x - 1) = 5$. \qed
\end{exemple}

% --- Exemple travaillé 2 ---
\begin{exemple}\label{ex:limite-sinx-sur-x}
\textbf{Montrer que} $\displaystyle\lim_{x \to 0} \frac{\sin x}{x} = 1$.

\medskip
\textbf{Méthode : théorème d'encadrement.} Pour $x \in \,]0, \pi/2[$, on sait (par argument géométrique sur le cercle trigonométrique) que
\[
  \cos x \leqslant \frac{\sin x}{x} \leqslant 1.
\]
Comme la fonction $x \mapsto \sin(x)/x$ est paire, cette encadrement vaut aussi pour $x \in \,]{-\pi/2}, 0[\,$. Or $\lim_{x \to 0} \cos x = 1$ et $\lim_{x \to 0} 1 = 1$. Par le théorème d'encadrement (des gendarmes) pour les fonctions :
\[
  \lim_{x \to 0} \frac{\sin x}{x} = 1.  \qed
\]
\end{exemple}

% ============================================================
\section{Limites à gauche, à droite, limites infinies}\label{sec:limites-laterales}
\index{limite!à gauche}\index{limite!à droite}\index{limite!infinie}

\begin{definition}[Limite à gauche et à droite]\label{def:limites-laterales}
Soient $f \colon D \to \R$ et $a \in \R$ adhérent à $D \cap \,]{-\infty}, a[\,$.

On dit que $f$ admet $\ell$ pour \textbf{limite à gauche} en $a$ si :
\[
  \forall\, \varepsilon > 0,\quad \exists\, \delta > 0,\quad
  \forall\, x \in D,\quad a - \delta < x < a \;\Longrightarrow\; |f(x) - \ell| < \varepsilon.
\]
On note $\lim_{x \to a^-} f(x) = \ell$ ou $f(a^-) = \ell$.

La \textbf{limite à droite} est définie de manière analogue en remplaçant $a - \delta < x < a$ par $a < x < a + \delta$.
\end{definition}

\begin{proposition}\label{prop:limite-laterale}
On a $\displaystyle\lim_{x \to a} f(x) = \ell$ si et seulement si $\displaystyle\lim_{x \to a^-} f(x) = \lim_{x \to a^+} f(x) = \ell$.
\end{proposition}

\begin{proof}
($\Rightarrow$) Immédiat : si la condition $\varepsilon$-$\delta$ vaut pour tout $x$ avec $0 < |x-a| < \delta$, elle vaut a fortiori pour $x < a$ ou $x > a$.

($\Leftarrow$) Soit $\varepsilon > 0$. Par hypothèse il existe $\delta_1, \delta_2 > 0$ tels que :
\begin{itemize}
  \item $a - \delta_1 < x < a \Rightarrow |f(x) - \ell| < \varepsilon$,
  \item $a < x < a + \delta_2 \Rightarrow |f(x) - \ell| < \varepsilon$.
\end{itemize}
Posons $\delta = \min(\delta_1, \delta_2) > 0$. Pour tout $x \in D$ avec $0 < |x-a| < \delta$, on a soit $a - \delta < x < a$, soit $a < x < a + \delta$. Dans les deux cas, $|f(x) - \ell| < \varepsilon$.
\end{proof}

\begin{definition}[Limite infinie]\label{def:limite-infinie}
On dit que $\lim_{x \to a} f(x) = +\infty$ si :
\[
  \forall\, M > 0,\quad \exists\, \delta > 0,\quad
  \forall\, x \in D,\quad 0 < |x - a| < \delta \;\Longrightarrow\; f(x) > M.
\]
On définit de même $\lim_{x \to a} f(x) = -\infty$ en remplaçant $f(x) > M$ par $f(x) < -M$.
\end{definition}

\begin{definition}[Limite en l'infini]\label{def:limite-infini}
\index{limite!en l'infini}
Soit $f \colon D \to \R$ avec $D$ non majoré. On dit que $\lim_{x \to +\infty} f(x) = \ell$ si :
\[
  \forall\, \varepsilon > 0,\quad \exists\, A > 0,\quad
  \forall\, x \in D,\quad x > A \;\Longrightarrow\; |f(x) - \ell| < \varepsilon.
\]
\end{definition}

% ============================================================
\section{Opérations sur les limites}\label{sec:operations-limites}
\index{limite!opérations}

\begin{theoreme}[Opérations algébriques sur les limites]\label{thm:operations-limites}
Soient $f, g \colon D \to \R$ et $a$ adhérent à $D \setminus \{a\}$. Supposons $\lim_{x \to a} f(x) = \ell$ et $\lim_{x \to a} g(x) = m$. Alors :
\begin{enumerate}[label=\textup{(\roman*)}]
  \item $\lim_{x \to a} (f(x) + g(x)) = \ell + m$ ;
  \item $\lim_{x \to a} (\lambda f(x)) = \lambda \ell$ pour tout $\lambda \in \R$ ;
  \item $\lim_{x \to a} f(x)\,g(x) = \ell\, m$ ;
  \item si $m \neq 0$, $\lim_{x \to a} \dfrac{f(x)}{g(x)} = \dfrac{\ell}{m}$.
\end{enumerate}
\end{theoreme}

\begin{proof}
Nous démontrons (i) et (iii) ; les autres sont similaires.

\textbf{(i) Somme.} Soit $\varepsilon > 0$. Il existe $\delta_1 > 0$ tel que $0 < |x - a| < \delta_1 \Rightarrow |f(x) - \ell| < \varepsilon/2$ et $\delta_2 > 0$ tel que $0 < |x - a| < \delta_2 \Rightarrow |g(x) - m| < \varepsilon/2$. Posons $\delta = \min(\delta_1, \delta_2)$. Pour $0 < |x - a| < \delta$ :
\[
  |(f(x) + g(x)) - (\ell + m)| \leqslant |f(x) - \ell| + |g(x) - m|
  < \frac{\varepsilon}{2} + \frac{\varepsilon}{2} = \varepsilon.
\]

\textbf{(iii) Produit.} Écrivons :
\[
  f(x)\,g(x) - \ell\, m = f(x)\bigl(g(x) - m\bigr) + m\bigl(f(x) - \ell\bigr).
\]
Soit $\varepsilon > 0$. Puisque $f$ admet une limite en $a$, $f$ est bornée au voisinage de $a$ : il existe $\delta_0 > 0$ et $C > 0$ tels que $0 < |x-a| < \delta_0 \Rightarrow |f(x)| \leqslant C$. Posons $M = \max(C, |m|) + 1$ (pour éviter $M = 0$). Il existe $\delta_1 > 0$ tel que $0 < |x-a| < \delta_1 \Rightarrow |g(x) - m| < \varepsilon/(2M)$ et $\delta_2 > 0$ tel que $0 < |x-a| < \delta_2 \Rightarrow |f(x) - \ell| < \varepsilon/(2M)$. Posons $\delta = \min(\delta_0, \delta_1, \delta_2)$. Pour $0 < |x-a| < \delta$ :
\[
  |f(x)\,g(x) - \ell\, m|
  \leqslant |f(x)|\,|g(x) - m| + |m|\,|f(x) - \ell|
  < C \cdot \frac{\varepsilon}{2M} + |m| \cdot \frac{\varepsilon}{2M}
  \leqslant \frac{\varepsilon}{2} + \frac{\varepsilon}{2} = \varepsilon.
\]
\end{proof}

% ============================================================
\section{Continuité en un point}\label{sec:continuite-point}
\index{continuité!en un point}

\begin{definition}[Continuité en un point]\label{def:continuite-point}
Soient $D \subset \R$, $f \colon D \to \R$ et $a \in D$. On dit que $f$ est \textbf{continue en $a$} si :
\[
  \forall\, \varepsilon > 0,\quad \exists\, \delta > 0,\quad
  \forall\, x \in D,\quad |x - a| < \delta \;\Longrightarrow\; |f(x) - f(a)| < \varepsilon.
\]
De manière équivalente : $\lim_{x \to a} f(x) = f(a)$.

\medskip
\textbf{Remarque essentielle.} La différence avec la définition de limite est double : (i)~on exige $a \in D$, et (ii)~la condition porte sur $|x - a| < \delta$ (et non $0 < |x - a| < \delta$), c'est-à-dire que $x = a$ est autorisé. Cela revient à imposer que la limite de $f$ en $a$ existe et vaut $f(a)$.
\end{definition}

\begin{proposition}[Caractérisation séquentielle de la continuité]\label{prop:continuite-sequentielle}
La fonction $f \colon D \to \R$ est continue en $a \in D$ si et seulement si, pour toute suite $(x_n)$ d'éléments de $D$ convergeant vers $a$, la suite $(f(x_n))$ converge vers $f(a)$.
\end{proposition}

\begin{proof}
Conséquence immédiate du théorème~\ref{thm:heine-limite}, en notant que la continuité en $a$ équivaut à $\lim_{x \to a} f(x) = f(a)$.
\end{proof}

\begin{exemple}\label{ex:polynome-continu}
\textbf{Tout polynôme est continu sur $\R$.} En effet, la fonction constante $x \mapsto c$ et la fonction identité $x \mapsto x$ sont continues (exercice immédiat). Par le théorème~\ref{thm:operations-limites}, les sommes et produits de fonctions continues sont continus. Tout polynôme étant combinaison linéaire de puissances de $x$, il est continu.
\end{exemple}

\begin{exemple}\label{ex:valeur-absolue-continue}
\textbf{La fonction $x \mapsto |x|$ est continue sur $\R$ (preuve $\varepsilon$-$\delta$).}

\medskip
\textbf{Brouillon.} On veut $\bigl||x| - |a|\bigr| < \varepsilon$. Par l'inégalité triangulaire inversée, $\bigl||x| - |a|\bigr| \leqslant |x - a|$. Il suffit donc de prendre $\delta = \varepsilon$.

\medskip
\textbf{Rédaction.} Soit $a \in \R$ et $\varepsilon > 0$. Posons $\delta = \varepsilon$. Pour tout $x \in \R$ tel que $|x - a| < \delta$ :
\[
  \bigl||x| - |a|\bigr| \leqslant |x - a| < \delta = \varepsilon.
\]
Donc $x \mapsto |x|$ est continue en $a$. Comme $a$ est arbitraire, elle est continue sur $\R$. \qed
\end{exemple}

\begin{exemple}[Fonction de Dirichlet]\label{ex:dirichlet}
\index{Dirichlet!fonction de}
Soit $\mathbf{1}_{\Q} \colon \R \to \R$ définie par
\[
  \mathbf{1}_{\Q}(x) =
  \begin{cases}
    1 & \text{si } x \in \Q, \\
    0 & \text{si } x \notin \Q.
  \end{cases}
\]

\textbf{Affirmation : $\mathbf{1}_{\Q}$ n'est continue en aucun point de $\R$.}

\medskip
\textbf{Preuve.} Soit $a \in \R$. Nous utilisons la caractérisation séquentielle.

\emph{Cas 1 : $a \in \Q$.} Alors $\mathbf{1}_{\Q}(a) = 1$. Par densité de $\R \setminus \Q$ dans $\R$, il existe une suite $(x_n)$ d'irrationnels convergeant vers $a$. On a $\mathbf{1}_{\Q}(x_n) = 0$ pour tout $n$, donc $\mathbf{1}_{\Q}(x_n) \to 0 \neq 1 = \mathbf{1}_{\Q}(a)$.

\emph{Cas 2 : $a \notin \Q$.} Alors $\mathbf{1}_{\Q}(a) = 0$. Par densité de $\Q$ dans $\R$, il existe une suite $(r_n)$ de rationnels convergeant vers $a$. On a $\mathbf{1}_{\Q}(r_n) = 1 \to 1 \neq 0 = \mathbf{1}_{\Q}(a)$.

Dans les deux cas, la caractérisation séquentielle n'est pas satisfaite. Donc $\mathbf{1}_{\Q}$ n'est continue en aucun point.
\end{exemple}

% ============================================================
\section{Continuité sur un intervalle — Théorème des valeurs intermédiaires}\label{sec:TVI}
\index{théorème des valeurs intermédiaires}\index{TVI|see{théorème des valeurs intermédiaires}}

\begin{definition}[Continuité sur un ensemble]\label{def:continuite-ensemble}
On dit que $f \colon D \to \R$ est \textbf{continue sur $D$} si $f$ est continue en tout point de $D$.
\end{definition}

\begin{theoreme}[Théorème des valeurs intermédiaires]\label{thm:TVI}
Soient $a < b$ deux réels et $f \colon [a, b] \to \R$ une fonction continue. Si $f(a) \leqslant 0 \leqslant f(b)$ (ou $f(b) \leqslant 0 \leqslant f(a)$), alors il existe $c \in [a, b]$ tel que $f(c) = 0$.

Plus généralement, pour toute valeur $\gamma$ comprise entre $f(a)$ et $f(b)$, il existe $c \in [a, b]$ tel que $f(c) = \gamma$.
\end{theoreme}

\begin{proof}
Traitons le cas $f(a) \leqslant 0 \leqslant f(b)$ (l'autre cas s'obtient en considérant $-f$). Si $f(a) = 0$, on prend $c = a$. Supposons donc $f(a) < 0 < f(b)$.

Considérons l'ensemble
\[
  E = \{x \in [a, b] : f(x) \leqslant 0\}.
\]
Cet ensemble est non vide (car $a \in E$) et majoré par $b$. Posons $c = \sup E$. On a $c \in [a, b]$.

\medskip
\textbf{Étape 1 : $f(c) \leqslant 0$.} Par définition du supremum, il existe une suite $(x_n)$ d'éléments de $E$ telle que $x_n \to c$. Pour tout $n$, $f(x_n) \leqslant 0$. Par continuité de $f$ en $c$ et caractérisation séquentielle, $f(x_n) \to f(c)$. Comme $f(x_n) \leqslant 0$ pour tout $n$, le passage à la limite donne $f(c) \leqslant 0$.

\medskip
\textbf{Étape 2 : $f(c) \geqslant 0$.} Supposons par l'absurde que $f(c) < 0$. Par continuité de $f$ en $c$, il existe $\delta > 0$ tel que pour tout $x \in [a, b]$ avec $|x - c| < \delta$, on a $f(x) < 0$ (en prenant $\varepsilon = |f(c)|/2$ dans la définition). En particulier, si $c < b$, alors pour $x = \min(c + \delta/2, b)$, on aurait $x > c$, $x \in [a, b]$ et $f(x) < 0$, c'est-à-dire $x \in E$, ce qui contredit $c = \sup E$. Si $c = b$, alors $f(b) < 0$, ce qui contredit $f(b) > 0$.

\medskip
\textbf{Conclusion.} On a $f(c) \leqslant 0$ et $f(c) \geqslant 0$, donc $f(c) = 0$.

Pour la forme générale, on applique ce résultat à la fonction $g(x) = f(x) - \gamma$.
\end{proof}

% --- TikZ : illustration du TVI ---
\begin{figure}[ht]
\centering
\begin{tikzpicture}[scale=1, >=Stealth]
  \draw[->] (-0.5,0) -- (7,0) node[right]{$x$};
  \draw[->] (0,-2) -- (0,3.5) node[above]{$y$};

  % Fonction continue
  \draw[blue, thick, domain=0.5:6, smooth, samples=100]
    plot(\x, {0.15*(\x-1)*(\x-3.5)*(\x-5.5)+0.5});

  % Points a, b
  \filldraw (0.5,{0.15*(0.5-1)*(0.5-3.5)*(0.5-5.5)+0.5}) circle(2pt);
  \filldraw (6,{0.15*(6-1)*(6-3.5)*(6-5.5)+0.5}) circle(2pt);

  \node[below] at (0.5,0) {$a$};
  \node[below] at (6,0) {$b$};

  % Valeur γ
  \draw[red, dashed] (-0.3,0.8) -- (6.5,0.8) node[right]{$\gamma$};

  % Intersection
  \draw[green!60!black, dashed] (1.45,0) -- (1.45,0.8);
  \node[below, green!60!black] at (1.45,0) {$c$};
  \filldraw[red] (1.45,0.8) circle(2pt);

  \node[left, blue] at (0.5,{0.15*(0.5-1)*(0.5-3.5)*(0.5-5.5)+0.5}) {$f(a)$};
  \node[right, blue] at (6,{0.15*(6-1)*(6-3.5)*(6-5.5)+0.5}) {$f(b)$};
\end{tikzpicture}
\caption{Théorème des valeurs intermédiaires : toute valeur $\gamma$ entre $f(a)$ et $f(b)$ est atteinte.}
\label{fig:TVI}
\end{figure}

\begin{corollaire}[Existence des racines $n$-ièmes]\label{cor:racine-n}
\index{racine $n$-ième}
Pour tout $y \geqslant 0$ et tout $n \geqslant 1$ entier, il existe un unique $x \geqslant 0$ tel que $x^n = y$.
\end{corollaire}

\begin{proof}
\emph{Existence.} Si $y = 0$, on prend $x = 0$. Si $y > 0$, considérons $f(x) = x^n - y$ sur $[0, M]$ où $M = \max(1, y)$. La fonction $f$ est continue, $f(0) = -y < 0$ et $f(M) = M^n - y \geqslant M - y \geqslant 0$ (car $M \geqslant y$ et $M \geqslant 1$). Par le TVI, il existe $c \in [0, M]$ tel que $f(c) = 0$, i.e.\ $c^n = y$.

\emph{Unicité.} La fonction $x \mapsto x^n$ est strictement croissante sur $[0, +\infty[$, donc l'équation $x^n = y$ a au plus une solution.
\end{proof}

\begin{corollaire}\label{cor:polynome-impair}
Tout polynôme de degré impair à coefficients réels admet au moins une racine réelle.
\end{corollaire}

\begin{proof}
Soit $P(x) = a_n x^n + \cdots + a_1 x + a_0$ avec $n$ impair et $a_n \neq 0$. On a $P(x) = a_n x^n\bigl(1 + o(1)\bigr)$ lorsque $|x| \to +\infty$. Puisque $n$ est impair, $\lim_{x \to +\infty} P(x) = +\infty \cdot \operatorname{signe}(a_n)$ et $\lim_{x \to -\infty} P(x) = -\infty \cdot \operatorname{signe}(a_n)$. Ainsi $P$ prend des valeurs de signes opposés pour $|x|$ assez grand. Par le TVI, $P$ s'annule.
\end{proof}

% ============================================================
\section{Continuité sur un compact}\label{sec:continuite-compact}
\index{continuité!sur un compact}\index{compact}

Dans $\R$, les compacts sont exactement les fermés bornés (théorème de Borel-Lebesgue, ou de Heine-Borel). En particulier, tout intervalle fermé borné $[a, b]$ est compact.

\begin{theoreme}[Théorème des bornes atteintes]\label{thm:bornes-atteintes}
\index{théorème des bornes atteintes}
Soit $f \colon [a, b] \to \R$ continue. Alors $f$ est bornée et atteint ses bornes : il existe $c, d \in [a, b]$ tels que
\[
  f(c) = \inf_{[a,b]} f \quad\text{et}\quad f(d) = \sup_{[a,b]} f.
\]
\end{theoreme}

\begin{proof}
\textbf{Étape 1 : $f$ est majorée.} Supposons par l'absurde que $f$ n'est pas majorée. Alors pour tout $n \geqslant 1$, il existe $x_n \in [a, b]$ tel que $f(x_n) > n$. La suite $(x_n)$ est bornée (car $x_n \in [a, b]$). Par le théorème de Bolzano-Weierstrass, il existe une sous-suite $(x_{\varphi(n)})$ convergeant vers un certain $c \in [a, b]$ (le segment $[a, b]$ est fermé, donc la limite y appartient). Par continuité de $f$ en $c$, $f(x_{\varphi(n)}) \to f(c)$. Mais $f(x_{\varphi(n)}) > \varphi(n) \geqslant n \to +\infty$, contradiction avec la convergence.

\textbf{Étape 2 : le supremum est atteint.} Posons $M = \sup_{[a,b]} f$ (qui est fini d'après l'étape 1). Par définition du supremum, pour tout $n \geqslant 1$, il existe $x_n \in [a, b]$ tel que $M - 1/n < f(x_n) \leqslant M$. Par Bolzano-Weierstrass, il existe une sous-suite $(x_{\varphi(n)})$ convergeant vers $d \in [a, b]$. Par continuité, $f(x_{\varphi(n)}) \to f(d)$. D'autre part, $f(x_{\varphi(n)}) \to M$ par encadrement. Par unicité de la limite, $f(d) = M$.

Le raisonnement pour l'infimum est analogue (ou appliquer le résultat à $-f$).
\end{proof}

\begin{theoreme}[Théorème de Heine]\label{thm:heine-unif}
\index{Heine!théorème de}\index{continuité!uniforme}
Soit $f \colon [a, b] \to \R$ continue. Alors $f$ est \textbf{uniformément continue} sur $[a, b]$.
\end{theoreme}

\begin{proof}
\textbf{Brouillon.} La continuité dit : pour chaque point $a$, pour tout $\varepsilon$, il existe $\delta$ (dépendant de $a$ et de $\varepsilon$). La continuité uniforme demande un $\delta$ qui marche pour \emph{tous} les points à la fois. Sur un compact, on va extraire une sous-suite convergente pour obtenir une contradiction.

\medskip
Procédons par l'absurde. Supposons que $f$ ne soit pas uniformément continue. Alors :
\[
  \exists\, \varepsilon_0 > 0,\quad \forall\, \delta > 0,\quad
  \exists\, x, y \in [a, b],\quad |x - y| < \delta \;\text{et}\; |f(x) - f(y)| \geqslant \varepsilon_0.
\]
Appliquons ceci avec $\delta = 1/n$ pour tout $n \geqslant 1$ : il existe $x_n, y_n \in [a, b]$ tels que $|x_n - y_n| < 1/n$ et $|f(x_n) - f(y_n)| \geqslant \varepsilon_0$.

Les suites $(x_n)$ et $(y_n)$ sont bornées (dans $[a, b]$). Par Bolzano-Weierstrass, il existe une sous-suite $(x_{\varphi(n)})$ convergeant vers $c \in [a, b]$. Puisque $|y_{\varphi(n)} - x_{\varphi(n)}| < 1/\varphi(n) \to 0$, on a aussi $y_{\varphi(n)} \to c$.

Par continuité de $f$ en $c$ :
\[
  f(x_{\varphi(n)}) \to f(c) \quad\text{et}\quad f(y_{\varphi(n)}) \to f(c).
\]
Donc $|f(x_{\varphi(n)}) - f(y_{\varphi(n)})| \to 0$, ce qui contredit $|f(x_{\varphi(n)}) - f(y_{\varphi(n)})| \geqslant \varepsilon_0 > 0$.
\end{proof}

% --- TikZ : continuité uniforme vs ponctuelle ---
\begin{figure}[ht]
\centering
\begin{tikzpicture}[scale=0.9, >=Stealth]
  % Gauche : uniforme
  \begin{scope}
    \node[above, font=\bfseries] at (3,4.5) {Continuité uniforme};
    \draw[->] (-0.3,0) -- (6.5,0) node[right]{$x$};
    \draw[->] (0,-0.3) -- (0,4.5);
    \draw[blue, thick, domain=0.2:6, smooth] plot(\x, {0.5*\x + 0.5});

    % Même δ partout
    \foreach \a in {1, 3, 5} {
      \draw[green!60!black, dashed] (\a-0.4, 0) -- (\a-0.4, 4.2);
      \draw[green!60!black, dashed] (\a+0.4, 0) -- (\a+0.4, 4.2);
      \draw[<->, green!60!black] (\a-0.4, -0.2) -- (\a+0.4, -0.2);
    }
    \node[below, green!60!black] at (3, -0.4) {même $\delta$ partout};
  \end{scope}

  % Droite : non uniforme
  \begin{scope}[xshift=9cm]
    \node[above, font=\bfseries] at (3,4.5) {Continue, non uniformément};
    \draw[->] (-0.3,0) -- (6.5,0) node[right]{$x$};
    \draw[->] (0,-0.3) -- (0,4.5);
    \draw[blue, thick, domain=0.5:6, smooth, samples=100]
      plot(\x, {3/\x});

    % δ décroissant
    \draw[green!60!black, dashed] (0.7,0) -- (0.7,4.2);
    \draw[green!60!black, dashed] (1.1,0) -- (1.1,4.2);
    \draw[<->, green!60!black] (0.7,-0.2) -- (1.1,-0.2);

    \draw[orange, dashed] (2.5,0) -- (2.5,4.2);
    \draw[orange, dashed] (3.5,0) -- (3.5,4.2);
    \draw[<->, orange] (2.5,-0.2) -- (3.5,-0.2);

    \node[below] at (3, -0.4) {$\delta$ dépend du point};
  \end{scope}
\end{tikzpicture}
\caption{Continuité uniforme (à gauche) : le même $\delta$ convient en tout point.
À droite : $f(x) = 1/x$ sur $]0, +\infty[$, le $\delta$ nécessaire diminue près de $0$.}
\label{fig:unif-continuite}
\end{figure}

% ============================================================
\section{Continuité uniforme}\label{sec:continuite-uniforme}

\begin{definition}[Continuité uniforme]\label{def:continuite-uniforme}
\index{continuité!uniforme}
Soit $f \colon D \to \R$. On dit que $f$ est \textbf{uniformément continue} sur $D$ si :
\[
  \forall\, \varepsilon > 0,\quad \exists\, \delta > 0,\quad
  \forall\, x, y \in D,\quad |x - y| < \delta \;\Longrightarrow\; |f(x) - f(y)| < \varepsilon.
\]
\end{definition}

\begin{remarque}\label{rem:unif-vs-ponctuelle}
La différence avec la continuité (ponctuelle) est que $\delta$ ne dépend que de $\varepsilon$, et non du point considéré. Toute fonction uniformément continue est continue, mais la réciproque est fausse en général (elle est vraie sur les compacts, par le théorème de Heine).
\end{remarque}

\begin{exemple}\label{ex:x2-non-unif}
\textbf{$f(x) = x^2$ n'est pas uniformément continue sur $\R$.}

\medskip
\textbf{Brouillon.} On a $|f(x) - f(y)| = |x^2 - y^2| = |x+y|\,|x-y|$. Même si $|x-y|$ est petit, le facteur $|x+y|$ peut être arbitrairement grand.

\medskip
\textbf{Preuve.} Prenons $\varepsilon_0 = 1$. Soit $\delta > 0$ quelconque. Choisissons $x = 1/\delta$ et $y = 1/\delta + \delta/2$. Alors $|x - y| = \delta/2 < \delta$, mais
\[
  |f(x) - f(y)| = |x^2 - y^2| = |x - y|\,|x + y|
  = \frac{\delta}{2}\left(\frac{2}{\delta} + \frac{\delta}{2}\right)
  = 1 + \frac{\delta^2}{4} \geqslant 1 = \varepsilon_0.
\]
Donc la condition de continuité uniforme n'est pas satisfaite.
\end{exemple}

% ============================================================
\section{Fonctions lipschitziennes}\label{sec:lipschitz}
\index{Lipschitz}

\begin{definition}[Fonction lipschitzienne]\label{def:lipschitz}
Soit $f \colon D \to \R$. On dit que $f$ est \textbf{$K$-lipschitzienne} (ou lipschitzienne de constante $K$) s'il existe $K \geqslant 0$ tel que :
\[
  \forall\, x, y \in D,\quad |f(x) - f(y)| \leqslant K\,|x - y|.
\]
\end{definition}

\begin{proposition}\label{prop:lipschitz-unif}
Toute fonction lipschitzienne est uniformément continue.
\end{proposition}

\begin{proof}
Soit $f$ une fonction $K$-lipschitzienne sur $D$. Si $K = 0$, $f$ est constante et le résultat est immédiat. Supposons $K > 0$. Soit $\varepsilon > 0$. Posons $\delta = \varepsilon/K > 0$. Pour tous $x, y \in D$ avec $|x - y| < \delta$ :
\[
  |f(x) - f(y)| \leqslant K\,|x - y| < K \cdot \frac{\varepsilon}{K} = \varepsilon. \qedhere
\]
\end{proof}

\begin{exemple}\label{ex:sin-lipschitz}
La fonction $\sin$ est $1$-lipschitzienne sur $\R$. En effet, pour tous $x, y \in \R$ :
\[
  |\sin x - \sin y| = \left|2\cos\frac{x+y}{2}\sin\frac{x-y}{2}\right|
  \leqslant 2 \cdot 1 \cdot \frac{|x-y|}{2} = |x - y|.
\]
(On a utilisé $|\sin u| \leqslant |u|$ pour tout $u$.) Donc $\sin$ est uniformément continue sur $\R$.
\end{exemple}

% ============================================================
\section{Erreurs fréquentes}\label{sec:erreurs-ch5}

\begin{enumerate}[label=\textbf{\arabic*.}]
  \item \textbf{Confondre continuité et continuité uniforme.} La continuité est une propriété \emph{locale} (en chaque point), la continuité uniforme est \emph{globale}. La fonction $x \mapsto x^2$ est continue sur $\R$ mais pas uniformément continue.

  \item \textbf{Mauvaise structure $\varepsilon$-$\delta$.} L'ordre des quantificateurs est $\forall\,\varepsilon > 0$, $\exists\,\delta > 0$, $\forall\,x$. On ne peut pas choisir $\delta$ \emph{après} $x$.

  \item \textbf{Oublier $0 < |x - a|$ dans la définition de limite.} La condition $|x - a| < \delta$ (sans $0 <$) définit la continuité, pas la limite.

  \item \textbf{Confondre la caractérisation séquentielle dans le bon sens.} Pour montrer que $f$ n'a \emph{pas} de limite, il suffit de trouver \emph{une} suite $(x_n) \to a$ telle que $(f(x_n))$ ne converge pas vers $\ell$ (ou deux suites donnant des limites différentes).
\end{enumerate}

% ============================================================
\section{Exercices}\label{sec:exercices-ch5}

\begin{exercice}\label{exo:limite-x2+1}
Montrer par un argument $\varepsilon$-$\delta$ que $\displaystyle\lim_{x \to 1}(x^2 + 1) = 2$.
\textit{Indication : écrire $|x^2 - 1| = |x-1|\,|x+1|$ et borner $|x+1|$ en se restreignant à $|x-1| < 1$.}
\end{exercice}

\begin{exercice}\label{exo:limite-1surx}
Montrer que $\displaystyle\lim_{x \to 2} \frac{1}{x} = \frac{1}{2}$.
\textit{Indication : $\left|\frac{1}{x} - \frac{1}{2}\right| = \frac{|x-2|}{2|x|}$. Borner $|x|$ inférieurement en se restreignant à $|x-2| < 1$.}
\end{exercice}

\begin{exercice}\label{exo:continuite-sqrt}
Montrer que $f(x) = \sqrt{x}$ est continue sur $[0, +\infty[$ par un argument $\varepsilon$-$\delta$.
\textit{Indication : $|\sqrt{x} - \sqrt{a}| = \frac{|x - a|}{|\sqrt{x} + \sqrt{a}|}$. Distinguer les cas $a = 0$ et $a > 0$.}
\end{exercice}

\begin{exercice}\label{exo:sqrt-unif-continue}
Montrer que $f(x) = \sqrt{x}$ est uniformément continue sur $[0, +\infty[$.
\textit{Indication : utiliser $|\sqrt{x} - \sqrt{y}| \leqslant \sqrt{|x-y|}$.}
\end{exercice}

\begin{exercice}\label{exo:TVI-cos}
Montrer qu'il existe $x \in [0, \pi/2]$ tel que $\cos x = x$.
\end{exercice}

\begin{exercice}\label{exo:TVI-polynome}
Montrer que le polynôme $P(x) = x^5 - 3x + 1$ admet au moins une racine réelle dans $[0, 1]$.
\end{exercice}

\begin{exercice}\label{exo:non-unif-sin1x}
Montrer que $f(x) = \sin(1/x)$ n'est pas uniformément continue sur $]0, 1]$.
\end{exercice}

\begin{exercice}\label{exo:lipschitz-sin}
Soit $f \colon [0, +\infty[ \to \R$ définie par $f(x) = \sin(\sqrt{x})$. Montrer que $f$ est uniformément continue sur $[0, +\infty[$. \textit{Indication : montrer d'abord que $f$ est lipschitzienne sur $[1, +\infty[$ et utiliser Heine sur $[0, 2]$.}
\end{exercice}

\begin{exercice}\label{exo:image-intervalle}
Soit $f \colon [a, b] \to \R$ continue. Montrer que $f([a, b])$ est un intervalle fermé borné. \textit{Indication : combiner le TVI et le théorème des bornes atteintes.}
\end{exercice}

\begin{exercice}\label{exo:thomae}
\index{Thomae!fonction de}
Soit $f \colon [0, 1] \to \R$ définie par $f(x) = 1/q$ si $x = p/q$ (fraction irréductible avec $q \geqslant 1$) et $f(x) = 0$ si $x \notin \Q$. Montrer que $f$ est continue en tout point irrationnel et discontinue en tout point rationnel.
\end{exercice}

\begin{exercice}\label{exo:point-fixe}
Soit $f \colon [0, 1] \to [0, 1]$ continue. Montrer que $f$ admet un point fixe, c'est-à-dire qu'il existe $c \in [0, 1]$ tel que $f(c) = c$. \textit{Indication : appliquer le TVI à $g(x) = f(x) - x$.}
\end{exercice}

\begin{exercice}\label{exo:cauchy-fonctionnel}
Soit $f \colon \R \to \R$ continue telle que $f(x+y) = f(x) + f(y)$ pour tous $x, y \in \R$. Montrer que $f(x) = f(1)\cdot x$ pour tout $x$.
\textit{Indication : montrer d'abord le résultat pour $x \in \N$, puis $x \in \Z$, puis $x \in \Q$, puis conclure par densité et continuité.}
\end{exercice}

% ============================================================
\section*{Résumé du chapitre~\thechapter}
\addcontentsline{toc}{section}{Résumé du chapitre~\thechapter}

\begin{itemize}
  \item La \textbf{limite} $\lim_{x \to a} f(x) = \ell$ est définie par le critère $\varepsilon$-$\delta$ : $\forall\,\varepsilon > 0$, $\exists\,\delta > 0$, $0 < |x-a| < \delta \Rightarrow |f(x) - \ell| < \varepsilon$.
  \item La \textbf{caractérisation séquentielle} (Heine) permet de ramener les limites de fonctions aux limites de suites.
  \item La \textbf{continuité} en $a$ signifie $\lim_{x \to a} f(x) = f(a)$.
  \item Le \textbf{TVI} affirme qu'une fonction continue sur $[a, b]$ prend toutes les valeurs entre $f(a)$ et $f(b)$.
  \item Sur un compact $[a, b]$, une fonction continue est \textbf{bornée et atteint ses bornes} (théorème des bornes atteintes).
  \item Sur un compact, continuité $\Rightarrow$ \textbf{continuité uniforme} (théorème de Heine).
  \item Lipschitzienne $\Rightarrow$ uniformément continue $\Rightarrow$ continue.
\end{itemize}

% ============================================================
% ============================================================
%  CHAPITRE 6 : Dérivation
% ============================================================
% ============================================================
\chapter{Dérivation --- Définition, Règles, Théorème des Accroissements Finis, Formule de Taylor}\label{chap:derivation}
\index{dérivation}\index{dérivée}

\section*{Introduction}

La notion de dérivée est l'un des piliers du calcul différentiel. Géométriquement, la dérivée $f'(a)$ mesure la pente de la tangente au graphe de $f$ au point $(a, f(a))$. Physiquement, elle représente un taux de variation instantané (vitesse, accélération, etc.). Nous allons définir rigoureusement cette notion, établir les règles de calcul, puis démontrer les théorèmes fondamentaux : Rolle, accroissements finis, et Taylor.

% ============================================================
\section{Dérivée en un point}\label{sec:derivee-point}

\begin{definition}[Dérivée en un point]\label{def:derivee}
\index{dérivée!en un point}
Soient $I$ un intervalle ouvert de $\R$, $f \colon I \to \R$ et $a \in I$. On dit que $f$ est \textbf{dérivable en $a$} si la limite
\[
  f'(a) = \lim_{h \to 0} \frac{f(a + h) - f(a)}{h}
  = \lim_{x \to a} \frac{f(x) - f(a)}{x - a}
\]
existe et est finie. Le nombre $f'(a) \in \R$ est appelé la \textbf{dérivée} de $f$ en $a$.

De manière équivalente, $f$ est dérivable en $a$ de dérivée $f'(a)$ si et seulement si :
\[
  f(a + h) = f(a) + f'(a)\,h + o(h) \quad (h \to 0),
\]
c'est-à-dire qu'il existe une fonction $\varepsilon(h)$ avec $\lim_{h \to 0} \varepsilon(h) = 0$ telle que $f(a+h) = f(a) + f'(a)\,h + h\,\varepsilon(h)$.
\end{definition}

\begin{theoreme}[Dérivable implique continue]\label{thm:derivable-continue}
\index{dérivabilité!implique continuité}
Si $f$ est dérivable en $a$, alors $f$ est continue en $a$.
\end{theoreme}

\begin{proof}
Supposons $f$ dérivable en $a$. Alors, d'après la caractérisation ci-dessus :
\[
  f(a + h) - f(a) = f'(a)\,h + h\,\varepsilon(h),
\]
où $\varepsilon(h) \to 0$ quand $h \to 0$. Donc :
\[
  \lim_{h \to 0} \bigl(f(a+h) - f(a)\bigr) = \lim_{h \to 0} h\bigl(f'(a) + \varepsilon(h)\bigr) = 0 \cdot \bigl(f'(a) + 0\bigr) = 0.
\]
Ainsi $\lim_{h \to 0} f(a+h) = f(a)$, ce qui signifie que $f$ est continue en $a$.
\end{proof}

\begin{remarque}\label{rem:reciproque-fausse}
La réciproque est fausse : une fonction peut être continue en un point sans y être dérivable. Voir l'exemple suivant.
\end{remarque}

% --- Contre-exemple : |x| ---
\begin{exemple}\label{ex:valeur-absolue-non-derivable}
\index{valeur absolue!non dérivable en $0$}
La fonction $f(x) = |x|$ est continue en $0$ (cf.\ exemple~\ref{ex:valeur-absolue-continue}) mais n'est pas dérivable en $0$.

\medskip
\textbf{Preuve.} Le taux d'accroissement en $0$ vaut :
\[
  \frac{f(h) - f(0)}{h} = \frac{|h|}{h} =
  \begin{cases}
    1 & \text{si } h > 0, \\
    -1 & \text{si } h < 0.
  \end{cases}
\]
Donc $\lim_{h \to 0^+} |h|/h = 1$ et $\lim_{h \to 0^-} |h|/h = -1$. Les limites à gauche et à droite étant distinctes, la limite n'existe pas. Donc $f$ n'est pas dérivable en $0$. \qed
\end{exemple}

% --- TikZ : |x| et la tangente ---
\begin{figure}[ht]
\centering
\begin{tikzpicture}[scale=1.2, >=Stealth]
  \draw[->] (-3,0) -- (3,0) node[right]{$x$};
  \draw[->] (0,-0.5) -- (0,3) node[above]{$y$};

  \draw[blue, thick] (-2.5,2.5) -- (0,0) -- (2.5,2.5);

  \draw[red, dashed, thick] (-2,1) -- (2,1) node[right, font=\small]{pente $+1$};
  \draw[orange, dashed, thick] (-2,-1+4) -- (0.5,0.5-0.5+2)
    node[right, font=\small]{pente $-1$};

  \filldraw[blue] (0,0) circle(1.5pt);
  \node[below right] at (0,0) {$0$};
  \node[above left, blue] at (-1.5,1.5) {$y = |x|$};

  \node[below, font=\small\itshape, text width=5cm, align=center] at (0,-0.8)
    {Point anguleux en $0$ : pas de tangente unique};
\end{tikzpicture}
\caption{La fonction $x \mapsto |x|$ : continue mais non dérivable en $0$.}
\label{fig:valeur-absolue}
\end{figure}

% --- Contre-exemple : x²sin(1/x) ---
\begin{exemple}\label{ex:x2-sin-1-sur-x}
\index{dérivée!discontinue}
Soit $f \colon \R \to \R$ définie par $f(x) = x^2 \sin(1/x)$ si $x \neq 0$ et $f(0) = 0$. Alors :
\begin{enumerate}[label=\textup{(\alph*)}]
  \item $f$ est dérivable en $0$ et $f'(0) = 0$.
  \item $f'$ n'est pas continue en $0$.
\end{enumerate}

\medskip
\textbf{Preuve de (a).} On calcule le taux d'accroissement :
\[
  \frac{f(h) - f(0)}{h} = \frac{h^2 \sin(1/h)}{h} = h \sin(1/h).
\]
Or $|h \sin(1/h)| \leqslant |h| \to 0$ quand $h \to 0$, donc $\lim_{h \to 0} h\sin(1/h) = 0$. Ainsi $f'(0) = 0$.

\textbf{Preuve de (b).} Pour $x \neq 0$, par les règles de dérivation :
\[
  f'(x) = 2x \sin(1/x) + x^2 \cdot \cos(1/x) \cdot (-1/x^2) = 2x\sin(1/x) - \cos(1/x).
\]
Le terme $2x\sin(1/x) \to 0$ quand $x \to 0$, mais $\cos(1/x)$ n'a pas de limite en $0$ (il oscille entre $-1$ et $1$). Donc $f'(x)$ n'a pas de limite en $0$, ce qui montre que $f'$ n'est pas continue en $0$.
\end{exemple}

% ============================================================
\section{Règles de dérivation}\label{sec:regles-derivation}
\index{dérivation!règles}

\begin{theoreme}[Opérations sur les dérivées]\label{thm:regles-derivation}
Soient $I$ un intervalle ouvert et $f, g \colon I \to \R$ dérivables en $a \in I$. Alors :
\begin{enumerate}[label=\textup{(\roman*)}]
  \item \textbf{Somme :} $f + g$ est dérivable en $a$ et $(f + g)'(a) = f'(a) + g'(a)$.
  \item \textbf{Produit (Leibniz)\index{Leibniz!formule de} :} $fg$ est dérivable en $a$ et $(fg)'(a) = f'(a)\,g(a) + f(a)\,g'(a)$.
  \item \textbf{Quotient :} si $g(a) \neq 0$, $f/g$ est dérivable en $a$ et
  \[
    \left(\frac{f}{g}\right)'(a) = \frac{f'(a)\,g(a) - f(a)\,g'(a)}{g(a)^2}.
  \]
\end{enumerate}
\end{theoreme}

\begin{proof}
\textbf{(i) Somme.} C'est immédiat par linéarité de la limite :
\[
  \frac{(f+g)(a+h) - (f+g)(a)}{h} = \frac{f(a+h)-f(a)}{h} + \frac{g(a+h)-g(a)}{h}
  \xrightarrow[h \to 0]{} f'(a) + g'(a).
\]

\textbf{(ii) Produit.} On écrit :
\begin{align*}
  \frac{(fg)(a+h) - (fg)(a)}{h}
  &= \frac{f(a+h)\,g(a+h) - f(a)\,g(a)}{h} \\
  &= \frac{f(a+h)\,g(a+h) - f(a+h)\,g(a) + f(a+h)\,g(a) - f(a)\,g(a)}{h} \\
  &= f(a+h)\,\frac{g(a+h) - g(a)}{h} + g(a)\,\frac{f(a+h) - f(a)}{h}.
\end{align*}
Quand $h \to 0$ : $f(a+h) \to f(a)$ (car $f$ dérivable, donc continue en $a$), et les taux d'accroissement convergent vers $g'(a)$ et $f'(a)$ respectivement. Donc :
\[
  (fg)'(a) = f(a)\,g'(a) + g(a)\,f'(a) = f'(a)\,g(a) + f(a)\,g'(a).
\]

\textbf{(iii) Quotient.} Il suffit de montrer que $1/g$ est dérivable et d'appliquer (ii). Puisque $g(a) \neq 0$ et $g$ est continue en $a$, $g$ ne s'annule pas au voisinage de $a$, et :
\begin{align*}
  \frac{1/g(a+h) - 1/g(a)}{h}
  &= \frac{g(a) - g(a+h)}{h\,g(a+h)\,g(a)}
  = \frac{-1}{g(a+h)\,g(a)} \cdot \frac{g(a+h) - g(a)}{h}.
\end{align*}
Quand $h \to 0$, $g(a+h) \to g(a)$, donc :
\[
  \left(\frac{1}{g}\right)'(a) = \frac{-g'(a)}{g(a)^2}.
\]
Puis $(f/g)'(a) = (f \cdot (1/g))'(a) = f'(a)/g(a) + f(a) \cdot (-g'(a)/g(a)^2) = (f'(a)\,g(a) - f(a)\,g'(a))/g(a)^2$.
\end{proof}

\begin{theoreme}[Dérivée d'une composée --- Règle de la chaîne]\label{thm:chaine}
\index{règle de la chaîne}\index{dérivation!composée}
Soient $I, J$ des intervalles ouverts, $g \colon I \to J$ dérivable en $a$ et $f \colon J \to \R$ dérivable en $g(a)$. Alors $f \circ g$ est dérivable en $a$ et
\[
  (f \circ g)'(a) = f'(g(a)) \cdot g'(a).
\]
\end{theoreme}

\begin{proof}
Puisque $g$ est dérivable en $a$, on a $g(a+h) = g(a) + g'(a)\,h + h\,\varepsilon_1(h)$ avec $\varepsilon_1(h) \to 0$.

Posons $k(h) = g(a+h) - g(a) = g'(a)\,h + h\,\varepsilon_1(h)$. Alors $k(h) \to 0$ quand $h \to 0$.

Puisque $f$ est dérivable en $b = g(a)$, on peut écrire : pour tout $k$,
\[
  f(b + k) - f(b) = f'(b)\,k + k\,\eta(k),
\]
où $\eta(k) \to 0$ quand $k \to 0$, en posant $\eta(0) = 0$ par convention. La fonction $\eta$ ainsi prolongée est continue en $0$.

Alors :
\begin{align*}
  f(g(a+h)) - f(g(a))
  &= f(b + k(h)) - f(b) \\
  &= f'(b)\,k(h) + k(h)\,\eta(k(h)) \\
  &= \bigl(f'(b) + \eta(k(h))\bigr)\,k(h) \\
  &= \bigl(f'(b) + \eta(k(h))\bigr)\bigl(g'(a)\,h + h\,\varepsilon_1(h)\bigr).
\end{align*}
En divisant par $h$ (pour $h \neq 0$) :
\[
  \frac{f(g(a+h)) - f(g(a))}{h}
  = \bigl(f'(b) + \eta(k(h))\bigr)\bigl(g'(a) + \varepsilon_1(h)\bigr).
\]
Quand $h \to 0$ : $k(h) \to 0$, donc $\eta(k(h)) \to 0$, et $\varepsilon_1(h) \to 0$. Ainsi :
\[
  (f \circ g)'(a) = f'(g(a)) \cdot g'(a). \qedhere
\]
\end{proof}

\begin{theoreme}[Dérivée de la fonction réciproque]\label{thm:derivee-reciproque}
\index{dérivée!de la réciproque}
Soit $f \colon I \to J$ bijective et continue, avec $I$ un intervalle. Si $f$ est dérivable en $a \in I$ et $f'(a) \neq 0$, alors $f^{-1}$ est dérivable en $b = f(a)$ et
\[
  (f^{-1})'(b) = \frac{1}{f'(a)} = \frac{1}{f'(f^{-1}(b))}.
\]
\end{theoreme}

\begin{proof}
Posons $g = f^{-1}$. Pour $k \neq 0$ avec $b + k \in J$, posons $h = g(b+k) - g(b)$. Comme $g$ est continue (théorème admis : la réciproque d'une bijection continue monotone sur un intervalle est continue), $h \to 0$ quand $k \to 0$. De plus, $h \neq 0$ car $g$ est injective. Par définition, $f(g(b) + h) = f(g(b+k)) = b + k$, donc $k = f(a + h) - f(a)$, d'où :
\[
  \frac{g(b+k) - g(b)}{k} = \frac{h}{f(a+h) - f(a)}
  = \frac{1}{\dfrac{f(a+h) - f(a)}{h}}.
\]
Quand $k \to 0$, $h \to 0$ et le dénominateur tend vers $f'(a) \neq 0$. Donc :
\[
  g'(b) = \frac{1}{f'(a)}. \qedhere
\]
\end{proof}

% ============================================================
\section{Extrema locaux}\label{sec:extrema}
\index{extremum local}

\begin{theoreme}[Condition nécessaire de Fermat]\label{thm:fermat}
\index{Fermat!condition nécessaire}
Soit $f \colon I \to \R$ où $I$ est un intervalle ouvert. Si $f$ admet un extremum local en $a \in I$ et si $f$ est dérivable en $a$, alors $f'(a) = 0$.
\end{theoreme}

\begin{proof}
Supposons que $f$ admet un maximum local en $a$ (le cas du minimum est analogue). Il existe $\delta > 0$ tel que $]a - \delta, a + \delta[ \subset I$ et $f(x) \leqslant f(a)$ pour tout $x \in \,]a - \delta, a + \delta[\,$.

Pour $h \in \,]0, \delta[\,$, on a $f(a+h) \leqslant f(a)$, donc $\dfrac{f(a+h) - f(a)}{h} \leqslant 0$. En passant à la limite $h \to 0^+$ :
\[
  f'(a) = \lim_{h \to 0^+} \frac{f(a+h) - f(a)}{h} \leqslant 0.
\]

Pour $h \in \,]{-\delta}, 0[\,$, on a $f(a+h) \leqslant f(a)$ et $h < 0$, donc $\dfrac{f(a+h) - f(a)}{h} \geqslant 0$. En passant à la limite $h \to 0^-$ :
\[
  f'(a) = \lim_{h \to 0^-} \frac{f(a+h) - f(a)}{h} \geqslant 0.
\]

Des deux inégalités, on déduit $f'(a) = 0$.
\end{proof}

\begin{theoreme}[Théorème de Rolle]\label{thm:rolle}
\index{Rolle!théorème de}
Soient $a < b$ deux réels et $f \colon [a, b] \to \R$ telle que :
\begin{enumerate}[label=\textup{(\roman*)}]
  \item $f$ est continue sur $[a, b]$ ;
  \item $f$ est dérivable sur $]a, b[$ ;
  \item $f(a) = f(b)$.
\end{enumerate}
Alors il existe $c \in \,]a, b[\,$ tel que $f'(c) = 0$.
\end{theoreme}

\begin{proof}
Par le théorème des bornes atteintes (théorème~\ref{thm:bornes-atteintes}), $f$ étant continue sur le compact $[a, b]$, elle atteint son maximum $M$ et son minimum $m$ sur $[a, b]$.

\emph{Cas 1 :} $m = M$. Alors $f$ est constante sur $[a, b]$, et $f'(c) = 0$ pour tout $c \in \,]a, b[\,$.

\emph{Cas 2 :} $m < M$. Puisque $f(a) = f(b)$, au moins l'une des valeurs $m$ ou $M$ est atteinte en un point intérieur $c \in \,]a, b[\,$. (En effet, si $m$ et $M$ étaient tous deux atteints uniquement aux extrémités, on aurait $f(a) = m$ et $f(b) = M$, ou l'inverse, ce qui contredirait $f(a) = f(b)$ puisque $m \neq M$.) Le point $c$ est donc un extremum local en un point intérieur où $f$ est dérivable. Par le théorème de Fermat (théorème~\ref{thm:fermat}), $f'(c) = 0$.
\end{proof}

% ============================================================
\section{Théorème des accroissements finis}\label{sec:TAF}
\index{théorème des accroissements finis}\index{TAF|see{théorème des accroissements finis}}

\begin{theoreme}[Théorème des accroissements finis (TAF)]\label{thm:TAF}
Soient $a < b$ deux réels et $f \colon [a, b] \to \R$ telle que :
\begin{enumerate}[label=\textup{(\roman*)}]
  \item $f$ est continue sur $[a, b]$ ;
  \item $f$ est dérivable sur $]a, b[\,$.
\end{enumerate}
Alors il existe $c \in \,]a, b[\,$ tel que
\[
  f'(c) = \frac{f(b) - f(a)}{b - a}.
\]
\end{theoreme}

\begin{proof}
Posons
\[
  \varphi(x) = f(x) - f(a) - \frac{f(b) - f(a)}{b - a}(x - a).
\]
La fonction $\varphi$ est continue sur $[a, b]$ (combinaison linéaire de fonctions continues), dérivable sur $]a, b[$, et :
\[
  \varphi(a) = f(a) - f(a) - 0 = 0, \qquad
  \varphi(b) = f(b) - f(a) - (f(b) - f(a)) = 0.
\]
Donc $\varphi(a) = \varphi(b)$. Par le théorème de Rolle (théorème~\ref{thm:rolle}), il existe $c \in \,]a, b[\,$ tel que $\varphi'(c) = 0$. Or :
\[
  \varphi'(x) = f'(x) - \frac{f(b) - f(a)}{b - a},
\]
donc $\varphi'(c) = 0$ donne $f'(c) = \dfrac{f(b) - f(a)}{b - a}$.
\end{proof}

% --- TikZ : illustration du TAF ---
\begin{figure}[ht]
\centering
\begin{tikzpicture}[scale=1, >=Stealth]
  \draw[->] (-0.5,0) -- (7,0) node[right]{$x$};
  \draw[->] (0,-0.5) -- (0,5) node[above]{$y$};

  % Fonction
  \draw[blue, thick, domain=0.8:6, smooth, samples=100]
    plot(\x, {0.08*(\x-1)*(\x-3)*(\x-6)+3});

  % Sécante
  \def\fa{0.08*(0.8-1)*(0.8-3)*(0.8-6)+3}
  \def\fb{0.08*(6-1)*(6-3)*(6-6)+3}
  \draw[red, thick] (0.8,{0.08*(0.8-1)*(0.8-3)*(0.8-6)+3})
    -- (6,{0.08*(6-1)*(6-3)*(6-6)+3});

  % Tangente parallèle (approximation c ≈ 4.1)
  \draw[green!60!black, thick, dashed]
    (2.5,{0.08*(4.1-1)*(4.1-3)*(4.1-6)+3 + (0.08*((4.1-3)*(4.1-6)+(4.1-1)*(4.1-6)+(4.1-1)*(4.1-3)))*(2.5-4.1)})
    -- (5.7,{0.08*(4.1-1)*(4.1-3)*(4.1-6)+3 + (0.08*((4.1-3)*(4.1-6)+(4.1-1)*(4.1-6)+(4.1-1)*(4.1-3)))*(5.7-4.1)});

  % Points
  \filldraw (0.8,{0.08*(0.8-1)*(0.8-3)*(0.8-6)+3}) circle(2pt);
  \filldraw (6,3) circle(2pt);

  \node[below] at (0.8,0) {$a$};
  \node[below] at (6,0) {$b$};
  \node[below, green!60!black] at (4.1,0) {$c$};

  \draw[green!60!black, dashed] (4.1,0) -- (4.1,{0.08*(4.1-1)*(4.1-3)*(4.1-6)+3});
  \filldraw[green!60!black] (4.1,{0.08*(4.1-1)*(4.1-3)*(4.1-6)+3}) circle(2pt);

  \node[right, red, font=\small] at (5.5,{0.08*(0.8-1)*(0.8-3)*(0.8-6)+3+0.5}) {sécante};
  \node[right, green!60!black, font=\small] at (5.2,1.5) {tangente en $c$};
\end{tikzpicture}
\caption{Théorème des accroissements finis : il existe $c \in \,]a,b[\,$ où la tangente est parallèle à la sécante.}
\label{fig:TAF}
\end{figure}

\begin{corollaire}\label{cor:derivee-nulle-constante}
Soit $f \colon I \to \R$ dérivable sur un intervalle $I$. Si $f'(x) = 0$ pour tout $x \in I$, alors $f$ est constante sur $I$.
\end{corollaire}

\begin{proof}
Soient $x, y \in I$ avec $x < y$. Par le TAF, il existe $c \in \,]x, y[\,$ tel que $f(y) - f(x) = f'(c)(y - x) = 0$. Donc $f(y) = f(x)$.
\end{proof}

\begin{corollaire}\label{cor:derivee-positive-croissante}
Soit $f \colon I \to \R$ dérivable sur un intervalle $I$.
\begin{enumerate}[label=\textup{(\roman*)}]
  \item Si $f'(x) \geqslant 0$ pour tout $x \in I$, alors $f$ est croissante sur $I$.
  \item Si $f'(x) > 0$ pour tout $x \in I$, alors $f$ est strictement croissante sur $I$.
\end{enumerate}
\end{corollaire}

\begin{proof}
Soient $x, y \in I$ avec $x < y$. Par le TAF, $f(y) - f(x) = f'(c)(y-x)$ pour un certain $c \in \,]x,y[\,$.

(i) Si $f' \geqslant 0$, alors $f'(c) \geqslant 0$ et $y - x > 0$, d'où $f(y) - f(x) \geqslant 0$, i.e.\ $f(y) \geqslant f(x)$.

(ii) Si $f' > 0$, alors $f'(c) > 0$, d'où $f(y) - f(x) > 0$, i.e.\ $f(y) > f(x)$.
\end{proof}

% ============================================================
\section{Théorème de L'Hôpital}\label{sec:hopital}
\index{L'Hôpital!théorème de}

\begin{theoreme}[Règle de L'Hôpital, cas $0/0$]\label{thm:hopital}
Soient $f, g \colon \,]a, b[\, \to \R$ dérivables avec $g'(x) \neq 0$ pour tout $x \in \,]a, b[\,$. On suppose que :
\begin{enumerate}[label=\textup{(\roman*)}]
  \item $\lim_{x \to a^+} f(x) = \lim_{x \to a^+} g(x) = 0$ ;
  \item $\lim_{x \to a^+} \dfrac{f'(x)}{g'(x)} = L$ \quad (où $L \in \R \cup \{-\infty, +\infty\}$).
\end{enumerate}
Alors $\lim_{x \to a^+} \dfrac{f(x)}{g(x)} = L$.
\end{theoreme}

\begin{proof}
On prolonge $f$ et $g$ par continuité en $a$ en posant $f(a) = g(a) = 0$. Ainsi $f$ et $g$ sont continues sur $[a, x]$ pour tout $x \in \,]a, b[\,$ et dérivables sur $]a, x[\,$. Pour $x \in \,]a, b[\,$, on a $g(x) \neq g(a) = 0$ (sinon, par Rolle, il existerait $c \in \,]a, x[\,$ avec $g'(c) = 0$, contradiction).

Par le \emph{théorème des accroissements finis généralisé} (théorème de Cauchy)\index{Cauchy!théorème des accroissements finis}, pour tout $x \in \,]a, b[\,$, il existe $c_x \in \,]a, x[\,$ tel que :
\[
  \frac{f(x)}{g(x)} = \frac{f(x) - f(a)}{g(x) - g(a)} = \frac{f'(c_x)}{g'(c_x)}.
\]
Quand $x \to a^+$, on a $c_x \to a^+$ (car $a < c_x < x$), donc $\dfrac{f'(c_x)}{g'(c_x)} \to L$. D'où $\dfrac{f(x)}{g(x)} \to L$.
\end{proof}

\begin{remarque}\label{rem:hopital-attention}
\textbf{Attention :} la réciproque est fausse. De plus, on ne peut appliquer L'Hôpital que si le quotient des dérivées admet une limite. Si $f'(x)/g'(x)$ n'a pas de limite, la règle ne permet pas de conclure (mais $f(x)/g(x)$ peut tout de même converger).
\end{remarque}

% ============================================================
\section{Formule de Taylor}\label{sec:taylor}
\index{Taylor!formule de}

\subsection{Polynôme de Taylor}

\begin{definition}[Polynôme de Taylor]\label{def:polynome-taylor}
\index{Taylor!polynôme de}
Soit $f \colon I \to \R$ de classe $\mathcal{C}^n$ au voisinage de $a \in I$. Le \textbf{polynôme de Taylor d'ordre $n$ de $f$ en $a$} est :
\[
  T_n(x) = \sum_{k=0}^{n} \frac{f^{(k)}(a)}{k!}(x - a)^k
  = f(a) + f'(a)(x-a) + \frac{f''(a)}{2}(x-a)^2 + \cdots + \frac{f^{(n)}(a)}{n!}(x-a)^n.
\]
\end{definition}

\subsection{Formule de Taylor-Lagrange}

\begin{theoreme}[Formule de Taylor-Lagrange]\label{thm:taylor-lagrange}
\index{Taylor!Lagrange}
Soit $f \colon [a, b] \to \R$ de classe $\mathcal{C}^n$ sur $[a, b]$, dérivable $n+1$ fois sur $]a, b[\,$. Alors il existe $c \in \,]a, b[\,$ tel que :
\[
  f(b) = \sum_{k=0}^{n} \frac{f^{(k)}(a)}{k!}(b - a)^k + \frac{f^{(n+1)}(c)}{(n+1)!}(b - a)^{n+1}.
\]
\end{theoreme}

\begin{proof}
Posons $h = b - a$ et définissons le \emph{reste} :
\[
  R = f(b) - \sum_{k=0}^{n} \frac{f^{(k)}(a)}{k!} h^k.
\]
Nous cherchons $c \in \,]a, b[\,$ tel que $R = \dfrac{f^{(n+1)}(c)}{(n+1)!}\,h^{n+1}$. Considérons la fonction auxiliaire :
\[
  g(t) = f(b) - \sum_{k=0}^{n} \frac{f^{(k)}(t)}{k!}(b - t)^k - \frac{R}{h^{n+1}}(b-t)^{n+1}.
\]
Calculons $g(a)$ et $g(b)$ :
\begin{itemize}
  \item $g(a) = f(b) - \displaystyle\sum_{k=0}^{n} \frac{f^{(k)}(a)}{k!}h^k - R = 0$ (par définition de $R$).
  \item $g(b) = f(b) - f(b) - 0 = 0$.
\end{itemize}
La fonction $g$ est continue sur $[a, b]$ et dérivable sur $]a, b[\,$, avec $g(a) = g(b) = 0$. Par le théorème de Rolle, il existe $c \in \,]a, b[\,$ tel que $g'(c) = 0$.

Calculons $g'(t)$. En dérivant terme à terme (et par télescopage) :
\begin{align*}
  g'(t) &= -\sum_{k=0}^{n}\left[\frac{f^{(k+1)}(t)}{k!}(b-t)^k - \frac{f^{(k)}(t)}{(k-1)!}(b-t)^{k-1}\right] + (n+1)\frac{R}{h^{n+1}}(b-t)^n,
\end{align*}
où le terme avec $(k-1)!$ n'apparaît que pour $k \geqslant 1$. En développant, les termes se simplifient par télescopage et il reste :
\[
  g'(t) = -\frac{f^{(n+1)}(t)}{n!}(b-t)^n + (n+1)\frac{R}{h^{n+1}}(b-t)^n.
\]
L'égalité $g'(c) = 0$ donne (avec $(b-c)^n \neq 0$ car $c < b$) :
\[
  \frac{f^{(n+1)}(c)}{n!} = (n+1)\frac{R}{h^{n+1}},
  \quad\text{d'où}\quad
  R = \frac{f^{(n+1)}(c)}{(n+1)!}\,h^{n+1}. \qedhere
\]
\end{proof}

\subsection{Formule de Taylor-Young}

\begin{theoreme}[Formule de Taylor-Young]\label{thm:taylor-young}
\index{Taylor!Young}
Soit $f$ de classe $\mathcal{C}^n$ au voisinage de $a$. Alors :
\[
  f(a + h) = \sum_{k=0}^{n} \frac{f^{(k)}(a)}{k!}\,h^k + o(h^n) \quad (h \to 0).
\]
\end{theoreme}

\begin{proof}
Posons $R_n(h) = f(a+h) - \displaystyle\sum_{k=0}^{n}\frac{f^{(k)}(a)}{k!}\,h^k$. Nous devons montrer que $R_n(h)/h^n \to 0$ quand $h \to 0$.

Notons que $R_n(0) = 0$. De plus, $R_n'(h) = f'(a+h) - \sum_{k=1}^{n}\frac{f^{(k)}(a)}{(k-1)!}\,h^{k-1}$, qui est le reste d'ordre $n-1$ de la formule de Taylor de $f'$ en $a$. Plus généralement, pour $0 \leqslant j \leqslant n$, $R_n^{(j)}(0) = 0$.

Appliquons la règle de L'Hôpital $n$ fois successivement (ce qui est licite car numérateur et dénominateur tendent vers $0$) :
\[
  \lim_{h \to 0} \frac{R_n(h)}{h^n}
  = \lim_{h \to 0} \frac{R_n'(h)}{n\,h^{n-1}}
  = \lim_{h \to 0} \frac{R_n''(h)}{n(n-1)\,h^{n-2}}
  = \cdots
  = \lim_{h \to 0} \frac{R_n^{(n)}(h)}{n!}.
\]
Or $R_n^{(n)}(h) = f^{(n)}(a+h) - f^{(n)}(a)$, qui tend vers $0$ quand $h \to 0$ par continuité de $f^{(n)}$. Donc $R_n(h)/h^n \to 0/n! = 0$.
\end{proof}

\subsection{Applications : développements de Taylor usuels}\label{subsec:DL-usuels}
\index{développement de Taylor}

En appliquant la formule de Taylor-Young en $a = 0$, on obtient les développements classiques (valables au voisinage de $0$) :

\begin{align}
  e^x &= \sum_{k=0}^{n} \frac{x^k}{k!} + o(x^n)
  = 1 + x + \frac{x^2}{2} + \frac{x^3}{6} + \cdots + \frac{x^n}{n!} + o(x^n), \label{eq:DL-exp} \\[6pt]
  \sin x &= \sum_{k=0}^{n} \frac{(-1)^k\,x^{2k+1}}{(2k+1)!} + o(x^{2n+2})
  = x - \frac{x^3}{6} + \frac{x^5}{120} - \cdots + o(x^{2n+2}), \label{eq:DL-sin} \\[6pt]
  \cos x &= \sum_{k=0}^{n} \frac{(-1)^k\,x^{2k}}{(2k)!} + o(x^{2n+1})
  = 1 - \frac{x^2}{2} + \frac{x^4}{24} - \cdots + o(x^{2n+1}), \label{eq:DL-cos} \\[6pt]
  \ln(1+x) &= \sum_{k=1}^{n} \frac{(-1)^{k+1}\,x^k}{k} + o(x^n)
  = x - \frac{x^2}{2} + \frac{x^3}{3} - \cdots + \frac{(-1)^{n+1}\,x^n}{n} + o(x^n), \label{eq:DL-ln} \\[6pt]
  (1+x)^\alpha &= \sum_{k=0}^{n} \binom{\alpha}{k} x^k + o(x^n), \quad
  \text{où } \binom{\alpha}{k} = \frac{\alpha(\alpha-1)\cdots(\alpha-k+1)}{k!}. \label{eq:DL-puissance}
\end{align}

\textbf{Justification pour $e^x$.} On a $f(x) = e^x$, $f^{(k)}(x) = e^x$ pour tout $k$, $f^{(k)}(0) = 1$. La formule de Taylor-Young donne directement~\eqref{eq:DL-exp}. La formule de Taylor-Lagrange donne de plus, pour le reste :
\[
  R_n(x) = \frac{e^c}{(n+1)!}\,x^{n+1} \quad\text{pour un certain $c$ entre $0$ et $x$.}
\]
Si $|x| \leqslant A$, alors $|R_n(x)| \leqslant \dfrac{e^A \cdot A^{n+1}}{(n+1)!} \to 0$ quand $n \to +\infty$, ce qui justifie l'identité $e^x = \sum_{k=0}^{+\infty} x^k/k!$.

% --- TikZ : approximations successives de sin ---
\begin{figure}[ht]
\centering
\begin{tikzpicture}[scale=0.75, >=Stealth]
  \draw[->] (-7,0) -- (7,0) node[right]{$x$};
  \draw[->] (0,-3) -- (0,3) node[above]{$y$};

  % sin x
  \draw[blue, very thick, domain=-6.5:6.5, smooth, samples=200]
    plot(\x, {sin(\x r)}) node[right]{$\sin x$};

  % T1 = x
  \draw[red, thin, domain=-2.5:2.5, smooth]
    plot(\x, {\x}) node[right, font=\scriptsize]{$T_1$};

  \begin{scope}
  \clip (-7,-3) rectangle (7,3);

  % T3 = x - x^3/6
  \draw[orange, thin, domain=-3.8:3.8, smooth, samples=100]
    plot(\x, {\x - \x*\x*\x/6}) node[right, font=\scriptsize]{$T_3$};

  % T5 = x - x^3/6 + x^5/120
  \draw[green!60!black, thin, domain=-4.2:4.2, smooth, samples=100]
    plot(\x, {\x - \x*\x*\x/6 + \x*\x*\x*\x*\x/120}) node[above right, font=\scriptsize]{$T_5$};

  % T7
  \draw[purple, thin, domain=-3.8:3.8, smooth, samples=150]
    plot(\x, {\x - \x*\x*\x/6 + \x*\x*\x*\x*\x/120 - \x*\x*\x*\x*\x*\x*\x/5040})
    node[below right, font=\scriptsize]{$T_7$};

  \end{scope}
\end{tikzpicture}
\caption{Approximations de Taylor successives $T_1, T_3, T_5, T_7$ de $\sin x$ en $0$.}
\label{fig:taylor-sin}
\end{figure}

% ============================================================
\section{Développements limités}\label{sec:DL}
\index{développement limité}

\begin{definition}[Développement limité]\label{def:DL}
Soit $f$ définie au voisinage de $a$. On dit que $f$ admet un \textbf{développement limité} (DL) d'ordre $n$ en $a$ s'il existe des réels $c_0, c_1, \ldots, c_n$ tels que :
\[
  f(a + h) = c_0 + c_1\,h + c_2\,h^2 + \cdots + c_n\,h^n + o(h^n) \quad (h \to 0).
\]
Si $f$ est $\mathcal{C}^n$ en $a$, alors $c_k = f^{(k)}(a)/k!$ et le DL coïncide avec la formule de Taylor-Young.
\end{definition}

\begin{proposition}[Unicité du DL]\label{prop:unicite-DL}
Le DL d'ordre $n$, s'il existe, est unique.
\end{proposition}

\begin{proof}
Supposons $f(a+h) = \sum_{k=0}^n c_k h^k + o(h^n) = \sum_{k=0}^n d_k h^k + o(h^n)$. Alors $\sum_{k=0}^n (c_k - d_k)h^k = o(h^n)$. Si tous les $c_k - d_k$ ne sont pas nuls, soit $j$ le plus petit indice avec $c_j \neq d_j$. En divisant par $h^j$, on obtient $(c_j - d_j) + o(1) = o(h^{n-j})$, d'où $c_j - d_j = 0$ en passant à la limite, contradiction.
\end{proof}

\subsection{Opérations sur les DL}

\begin{proposition}\label{prop:operations-DL}
Soient $f$ et $g$ admettant des DL d'ordre $n$ en $a$. Alors :
\begin{enumerate}[label=\textup{(\roman*)}]
  \item $f + g$ admet un DL d'ordre $n$ en $a$, obtenu en sommant les DL.
  \item $f \cdot g$ admet un DL d'ordre $n$ en $a$, obtenu en multipliant les DL et en tronquant à l'ordre $n$.
  \item Si $g(a) \neq 0$, $f/g$ admet un DL d'ordre $n$ en $a$, obtenu par division selon les puissances croissantes.
  \item Si $f$ admet un DL en $a$ et $g$ un DL en $f(a)$ avec $f(a)$ bien défini, alors $g \circ f$ admet un DL en $a$ (par substitution et troncature).
\end{enumerate}
\end{proposition}

\subsection{Applications des DL au calcul de limites}

\begin{exemple}\label{ex:DL-limite-1}
\textbf{Calculer} $\displaystyle\lim_{x \to 0} \frac{e^x - 1 - x}{x^2}$.

En $0$ : $e^x = 1 + x + x^2/2 + o(x^2)$, donc $e^x - 1 - x = x^2/2 + o(x^2)$, d'où :
\[
  \frac{e^x - 1 - x}{x^2} = \frac{x^2/2 + o(x^2)}{x^2} = \frac{1}{2} + o(1)
  \xrightarrow[x \to 0]{} \frac{1}{2}.
\]
\end{exemple}

\begin{exemple}\label{ex:DL-limite-2}
\textbf{Calculer} $\displaystyle\lim_{x \to 0} \frac{\sin x - x + x^3/6}{x^5}$.

En $0$ : $\sin x = x - x^3/6 + x^5/120 + o(x^5)$, donc $\sin x - x + x^3/6 = x^5/120 + o(x^5)$, d'où :
\[
  \frac{\sin x - x + x^3/6}{x^5} = \frac{1}{120} + o(1)
  \xrightarrow[x \to 0]{} \frac{1}{120}.
\]
\end{exemple}

\begin{exemple}\label{ex:DL-limite-3}
\textbf{Calculer} $\displaystyle\lim_{x \to 0} \frac{\ln(1+x) - x}{x^2}$.

En $0$ : $\ln(1+x) = x - x^2/2 + o(x^2)$, donc $\ln(1+x) - x = -x^2/2 + o(x^2)$, d'où :
\[
  \frac{\ln(1+x) - x}{x^2} = -\frac{1}{2} + o(1)
  \xrightarrow[x \to 0]{} -\frac{1}{2}.
\]
\end{exemple}

% ============================================================
\section{Erreurs fréquentes}\label{sec:erreurs-ch6}

\begin{enumerate}[label=\textbf{\arabic*.}]
  \item \textbf{Appliquer L'Hôpital sans vérifier les hypothèses.} Avant d'appliquer la règle, on doit vérifier que l'on est bien dans une forme indéterminée $0/0$ ou $\infty/\infty$, et que le quotient des dérivées admet une limite.

  \item \textbf{Oublier les hypothèses du TAF.} Le théorème des accroissements finis requiert la continuité sur le \emph{fermé} $[a, b]$ et la dérivabilité sur l'\emph{ouvert} $]a, b[\,$. Certains étudiants oublient la continuité aux bords.

  \item \textbf{Confondre condition nécessaire et suffisante pour un extremum.} $f'(a) = 0$ est une condition \emph{nécessaire} d'extremum local (si $f$ est dérivable), pas suffisante. Exemple : $f(x) = x^3$ satisfait $f'(0) = 0$ mais $0$ n'est pas un extremum.

  \item \textbf{Tronquer un DL trop tard ou trop tôt.} Dans un produit de DL, il faut tronquer à l'ordre voulu \emph{après} la multiplication, pas avant. Garder des termes d'ordre supérieur à $n$ est inutile et source d'erreurs.

  \item \textbf{Écrire $o(x^n) + o(x^n) = o(x^{2n})$.} C'est faux ! On a $o(x^n) + o(x^n) = o(x^n)$.
\end{enumerate}

% ============================================================
\section{Exercices}\label{sec:exercices-ch6}

\begin{exercice}\label{exo:derivable-x3}
Calculer la dérivée de $f(x) = x^3$ en tout point par la définition (limite du taux d'accroissement).
\end{exercice}

\begin{exercice}\label{exo:derive-sqrt-def}
Montrer par la définition que $f(x) = \sqrt{x}$ est dérivable sur $]0, +\infty[$ et que $f'(x) = 1/(2\sqrt{x})$.
\end{exercice}

\begin{exercice}\label{exo:derivee-xn}
Montrer par récurrence que $(x^n)' = n\,x^{n-1}$ pour tout $n \in \N^*$.
\end{exercice}

\begin{exercice}\label{exo:rolle-application}
Soit $P(x) = x^3 - 3x + 1$. Montrer que $P$ a exactement trois racines réelles distinctes.
\textit{Indication : étudier les variations de $P$ et utiliser le TVI.}
\end{exercice}

\begin{exercice}\label{exo:TAF-inegalite}
En utilisant le TAF, montrer que pour tout $x > 0$ :
\[
  \frac{x}{1+x} \leqslant \ln(1 + x) \leqslant x.
\]
\end{exercice}

\begin{exercice}\label{exo:TAF-sin}
Montrer que $|\sin x - \sin y| \leqslant |x - y|$ pour tous $x, y \in \R$, en utilisant le TAF.
\end{exercice}

\begin{exercice}\label{exo:hopital-exemples}
Calculer les limites suivantes en justifiant l'utilisation de la règle de L'Hôpital :
\begin{enumerate}[label=\textup{(\alph*)}]
  \item $\displaystyle\lim_{x \to 0} \frac{e^x - 1 - x}{x^2}$ ;
  \item $\displaystyle\lim_{x \to 0} \frac{x - \sin x}{x^3}$ ;
  \item $\displaystyle\lim_{x \to 0^+} x\ln x$.
\end{enumerate}
\end{exercice}

\begin{exercice}\label{exo:DL-calcul}
Calculer les DL à l'ordre $3$ en $0$ des fonctions suivantes :
\begin{enumerate}[label=\textup{(\alph*)}]
  \item $f(x) = e^{\sin x}$ ;
  \item $f(x) = \sqrt{1 + x}$ ;
  \item $f(x) = \dfrac{1}{1 - x + x^2}$.
\end{enumerate}
\end{exercice}

\begin{exercice}\label{exo:DL-limite}
À l'aide de DL, calculer :
\[
  \lim_{x \to 0} \frac{\tan x - \sin x}{x^3}.
\]
\end{exercice}

\begin{exercice}\label{exo:taylor-reste}
Soit $f(x) = \cos x$. Écrire la formule de Taylor-Lagrange d'ordre $2n$ en $0$ et en déduire que, pour tout $x \in \R$,
\[
  \left|\cos x - \sum_{k=0}^{n} \frac{(-1)^k x^{2k}}{(2k)!}\right| \leqslant \frac{|x|^{2n+2}}{(2n+2)!}.
\]
\end{exercice}

\begin{exercice}\label{exo:point-inflexion}
Soit $f \colon I \to \R$ de classe $\mathcal{C}^3$ et $a \in I$ tel que $f''(a) = 0$ et $f'''(a) \neq 0$. Montrer que $a$ est un point d'inflexion de $f$.
\textit{Indication : utiliser le DL d'ordre $3$.}
\end{exercice}

\begin{exercice}\label{exo:inegalite-convexite}
Soit $f \colon I \to \R$ de classe $\mathcal{C}^2$ avec $f'' \geqslant 0$. Montrer que pour tous $x, a \in I$ :
\[
  f(x) \geqslant f(a) + f'(a)(x - a).
\]
\textit{Indication : formule de Taylor-Lagrange à l'ordre $1$.}
\end{exercice}

% ============================================================
\section*{Résumé du chapitre~\thechapter}
\addcontentsline{toc}{section}{Résumé du chapitre~\thechapter}

\begin{itemize}
  \item La \textbf{dérivée} $f'(a) = \lim_{h \to 0} (f(a+h) - f(a))/h$ mesure le taux de variation instantané.
  \item Dérivable en $a$ $\Rightarrow$ continue en $a$ (réciproque fausse).
  \item \textbf{Règles :} $(f+g)' = f'+g'$, $(fg)' = f'g+fg'$, $(f/g)' = (f'g-fg')/g^2$, $(f\circ g)' = (f'\circ g)\cdot g'$.
  \item \textbf{Fermat :} si $f$ dérivable a un extremum local intérieur, alors $f' = 0$.
  \item \textbf{Rolle :} $f(a) = f(b)$, $f$ continue sur $[a,b]$, dérivable sur $]a,b[$ $\Rightarrow$ $\exists\,c$, $f'(c) = 0$.
  \item \textbf{TAF :} $f(b)-f(a) = f'(c)(b-a)$ pour un certain $c \in \,]a,b[\,$.
  \item \textbf{Conséquences :} $f'=0 \Rightarrow f$ constante ; $f'>0 \Rightarrow f$ strictement croissante.
  \item \textbf{L'Hôpital :} dans les formes indéterminées, $f/g$ a la même limite que $f'/g'$ (si celle-ci existe).
  \item \textbf{Taylor-Lagrange :} reste exact $f^{(n+1)}(c)/(n+1)! \cdot h^{n+1}$.
  \item \textbf{Taylor-Young :} reste en $o(h^n)$, fondement des développements limités.
  \item \textbf{DL usuels :} $e^x$, $\sin x$, $\cos x$, $\ln(1+x)$, $(1+x)^\alpha$.
\end{itemize}

% ============================================================
%  Bibliographie
% ============================================================
\begin{thebibliography}{99}

\bibitem{rudin}
W.~\textsc{Rudin},
\textit{Principles of Mathematical Analysis},
3\ieme{} éd., McGraw-Hill, 1976.
\index{Rudin}

\bibitem{zuily}
C.~\textsc{Zuily} et H.~\textsc{Queffélec},
\textit{Analyse pour l'agrégation},
Dunod, 2013.
\index{Zuily}\index{Queffélec}

\bibitem{ramis}
E.~\textsc{Ramis}, C.~\textsc{Deschamps} et J.~\textsc{Odoux},
\textit{Cours de mathématiques spéciales},
tomes~1--4, Masson, 1991.
\index{Ramis}\index{Deschamps}\index{Odoux}

\bibitem{tao}
T.~\textsc{Tao},
\textit{Analysis~I},
3\ieme{} éd., Springer, Texts and Readings in Mathematics, 2016.
\index{Tao}

\bibitem{pompeiani}
F.~\textsc{Pompeïani} et P.~\textsc{Tarrago},
\textit{Analyse},
Ellipses, 1998.
\index{Pompeïani}\index{Tarrago}

\end{thebibliography}

% ============================================================
\printindex
\end{document}
